% !TeX program = xelatex
% !TeX encoding = UTF-8
\documentclass[UTF8,10pt,a4paper,fontset=fandol]{ctexart}
\usepackage{amsmath,amssymb,graphicx,tabularx,array,multirow,enumitem,titlesec,xcolor,ragged2e,etoolbox}
\usepackage[a4paper,top=1.6cm,bottom=1.6cm,left=1.6cm,right=1.6cm]{geometry}
\setlength{\parindent}{2em}
\setlength{\parskip}{0.15em}
\setlength{\tabcolsep}{4pt}
\renewcommand{\arraystretch}{1.15}
\setlist[enumerate]{leftmargin=2.6em,itemsep=0.1em,topsep=0.2em}
\titleformat{\section}{\large\bfseries}{\thesection}{0.6em}{}
\titleformat{\subsection}{\normalsize\bfseries}{\thesubsection}{0.6em}{}
\titleformat{\subsubsection}{\normalsize\bfseries}{\thesubsubsection}{0.6em}{}
\titlespacing*{\section}{0pt}{0.7em}{0.25em}
\titlespacing*{\subsection}{0pt}{0.55em}{0.2em}
\titlespacing*{\subsubsection}{0pt}{0.4em}{0.15em}
\raggedbottom
\BeforeBeginEnvironment{tabularx}{\par\noindent}
\newcommand{\fieldvalue}[1]{#1}
\newcommand{\handwritten}[1]{#1}
\let\mathcheckmark\checkmark
\renewcommand{\checkmark}{\ensuremath{\mathcheckmark}}
\begin{document}

% TODO: The company logo and company name are visible in the upper-left corner.

\title{艾米迈托赛注射液暂停放行通知单}
\date{}
\maketitle
\vspace{-2.2em}

\begin{center}
\small
\begin{tabularx}{\textwidth}{|>{\centering\arraybackslash}p{1.35cm}|X|>{\centering\arraybackslash}p{1.55cm}|X|}
\hline
产品名称 &
% #VALUE_ID: LEX-P0001-V0001
% #FIELD_VALUE: 产品名称
\fieldvalue{艾米迈托赛注射液}
&
代码 &
% #VALUE_ID: LEX-P0001-V0002
% #FIELD_VALUE: 代码
\fieldvalue{C304}
\\ \hline
规\quad 格 &
% #VALUE_ID: LEX-P0001-V0003
% #FIELD_VALUE: 规格
\fieldvalue{12ml：$6.0\times 10^{7}$个细胞}
&
产品批号 &
% #VALUE_ID: LEX-P0001-V0004
% #FIELD_VALUE: 产品批号
\fieldvalue{A33Z2202506001}
\\ \hline
数\quad 量 &
% #VALUE_ID: LEX-P0001-V0005
% #FIELD_VALUE: 数量
\fieldvalue{541袋}
&
检验报告编号 &
% #VALUE_ID: LEX-P0001-V0006
% #FIELD_VALUE: 检验报告编号
\fieldvalue{C3042025080703}
\\ \hline
\multirow{1}{*}{效\quad 期} &
% #VALUE_ID: LEX-P0001-V0007
% #FIELD_VALUE: 有效期选择
% #HANDWRITTEN: √
\fieldvalue{\handwritten{\ensuremath{\checkmark}}}\quad 有效期至：&
\multicolumn{2}{l|}{
% #VALUE_ID: LEX-P0001-V0008
% #FIELD_VALUE: 有效期日期
\fieldvalue{2026年06月17日}
}
\\ \hline
\multicolumn{4}{|l|}{
质量评价：\quad
% #VALUE_ID: LEX-P0001-V0009
% #FIELD_VALUE: 质量评价
\fieldvalue{\textcolor{red}{检测合格，但不满足放行标准，暂停放行！}}
}
\\ \hline
质量受权人 &
% #VALUE_ID: LEX-P0001-V0010
% #FIELD_VALUE: 质量受权人
% #HANDWRITTEN: 冯春艳
\fieldvalue{\handwritten{冯春艳}}
&
日\quad 期 &
% #VALUE_ID: LEX-P0001-V0011
% #FIELD_VALUE: 日期
% #HANDWRITTEN: 2025年6月23日
\fieldvalue{\handwritten{2025年6月23日}}
\\ \hline
\end{tabularx}
\end{center}

\vspace{0.3em}
% #VALUE_ID: LEX-P0001-V0012
% #FIELD_VALUE: 质量评价旁注
% TODO #HANDWRITTEN: 手写内容难辨；蓝色批注覆盖在质量评价栏右侧
% #HANDWRITTEN: 手写内容难辨
\noindent\fieldvalue{\handwritten{手写内容难辨}}

\vspace{0.2em}
% #VALUE_ID: LEX-P0001-V0013
% #FIELD_VALUE: 底部旁注
% TODO #HANDWRITTEN: 手写内容难辨；位于日期栏下方
% #HANDWRITTEN: 手写内容难辨
\noindent\fieldvalue{\handwritten{手写内容难辨}}

% LEXOID_PAGE_COMPLETED: 1/59

\newpage

\begin{center}
\textbf{艾米迩托赛注射液审核放行订单}
\end{center}

\begin{tabularx}{\textwidth}{|p{2.2cm}|X|p{2.2cm}|p{2.2cm}|}
\hline
产品名称 &
% #VALUE_ID: LEX-P0002-V0001
% #FIELD_VALUE: 产品名称
\fieldvalue{艾米迩托赛注射液}
& 代码 &
% #VALUE_ID: LEX-P0002-V0002
% #FIELD_VALUE: 代码
\fieldvalue{C304}
\\ \hline
产品批号 &
% #VALUE_ID: LEX-P0002-V0003
% #FIELD_VALUE: 产品批号
\fieldvalue{A33Z2202506001}
& 规格 &
% #VALUE_ID: LEX-P0002-V0004
% #FIELD_VALUE: 规格
\fieldvalue{12ml：$6.0\times10^{7}$个细胞}
\\ \hline
&& 数量 &
% #VALUE_ID: LEX-P0002-V0005
% #FIELD_VALUE: 数量
\fieldvalue{541袋}
\\ \hline
\multicolumn{1}{|c|}{工序/项目}
& \multicolumn{2}{c|}{审核内容及标准} & \multicolumn{1}{c|}{结果}
\\ \hline
\multirow{1}{*}{物料检查}
& \multicolumn{2}{l|}{生产所用原辅料、耗材、包装材料等均为经放行的合格物料。}
& 符合规定：
% #VALUE_ID: LEX-P0002-V0006
% #FIELD_VALUE: 物料检查结果
\fieldvalue{\(\checkmark\) 是}\quad $\square$ 否
\\ \hline
\multirow{5}{*}{生产}
& \multicolumn{2}{l|}{生产过程工艺参数是否与工艺规程一致，生产操作是否符合相应的操作规程要求。}
& 符合规定：
% #VALUE_ID: LEX-P0002-V0007
% #FIELD_VALUE: 生产过程工艺参数及操作结果
\fieldvalue{\(\checkmark\) 是}\quad $\square$ 否
\\ \cline{2-4}
& \multicolumn{2}{l|}{物料平衡是否符合要求。}
& 符合规定：
% #VALUE_ID: LEX-P0002-V0008
% #FIELD_VALUE: 物料平衡结果
\fieldvalue{\(\checkmark\) 是}\quad $\square$ 否
\\ \cline{2-4}
& \multicolumn{2}{l|}{主要生产工艺和生产设备已经验证/确认。}
& 符合规定：
% #VALUE_ID: LEX-P0002-V0009
% #FIELD_VALUE: 生产工艺和设备验证结果
\fieldvalue{\(\checkmark\) 是}\quad $\square$ 否
\\ \cline{2-4}
& \multicolumn{2}{l|}{设备状态、清场等均符合要求。}
& 符合规定：
% #VALUE_ID: LEX-P0002-V0010
% #FIELD_VALUE: 设备状态及清场结果
\fieldvalue{\(\checkmark\) 是}\quad $\square$ 否
\\ \cline{2-4}
& \multicolumn{2}{l|}{}
& 符合规定：
% #VALUE_ID: LEX-P0002-V0011
% #FIELD_VALUE: 生产审核结果
\fieldvalue{\(\checkmark\) 是}\quad $\square$ 否
\\ \hline
\multirow{5}{*}{偏差}
& \multicolumn{2}{l|}{生产过程是否发生偏差；\newline
如发生偏差是否已关闭；\newline
如偏差未关闭，放行前确认是否对产品质量无影响。}
& $\square$ 无偏差发生\\
& & &
% #VALUE_ID: LEX-P0002-V0012
% #FIELD_VALUE: 偏差发生情况
\fieldvalue{\(\checkmark\) 有偏差发生，\newline 偏差编号为：}
\\ \cline{4-4}
& \multicolumn{2}{l|}{}
&
% #VALUE_ID: LEX-P0002-V0013
% #FIELD_VALUE: 偏差编号
% #HANDWRITTEN: DE2025085
\fieldvalue{\handwritten{DE2025085}}
\\ \cline{4-4}
& \multicolumn{2}{l|}{}
&
% #VALUE_ID: LEX-P0002-V0014
% #FIELD_VALUE: 偏差编号
% #HANDWRITTEN: DE2025087
\fieldvalue{\handwritten{DE2025087}}
\\ \cline{4-4}
& \multicolumn{2}{l|}{}
& $\square$ 偏差已关闭
\\ \cline{4-4}
& \multicolumn{2}{l|}{}
&
% #VALUE_ID: LEX-P0002-V0015
% #FIELD_VALUE: 偏差处理状态
\fieldvalue{\(\checkmark\) 偏差未关闭，但对产品质量无影响}
\\ \cline{4-4}
& \multicolumn{2}{l|}{}
& $\square$ 偏差未关闭，对产品质量有影响
\\ \hline
\multirow{6}{*}{变更}
& \multicolumn{2}{l|}{生产过程是否发生变更；\newline
如发生变更是否已关闭；\newline
如变更未关闭，放行前确认是否对产品质量无影响。}
&
% #VALUE_ID: LEX-P0002-V0016
% #FIELD_VALUE: 变更发生情况
\fieldvalue{\(\checkmark\) 无变更发生}
\\ \cline{4-4}
& \multicolumn{2}{l|}{}
& $\square$ 有变更发生，变更编号为：
\\ \cline{4-4}
& \multicolumn{2}{l|}{}
& $\square$ 变更已关闭
\\ \cline{4-4}
& \multicolumn{2}{l|}{}
& $\square$ 变更未关闭，但对产品质量无影响
\\ \cline{4-4}
& \multicolumn{2}{l|}{}
& $\square$ 变更未关闭，对产品质量有影响
\\ \hline
记录完整性检查
& \multicolumn{2}{l|}{批生产记录完整，并完成批生产记录审核签批。}
& 符合规定：
% #VALUE_ID: LEX-P0002-V0017
% #FIELD_VALUE: 记录完整性检查结果
\fieldvalue{\(\checkmark\) 是}\quad $\square$ 否
\\ \hline
\end{tabularx}

\vfill
\begin{center}
第 1 页 共 3 页
\end{center}
% LEXOID_PAGE_COMPLETED: 2/59

\newpage

\begin{center}
文件编码：REC-YZ-SMP-QA-01-007-a \qquad 版本号：1.6
\end{center}

\begin{tabularx}{\textwidth}{|p{0.14\textwidth}|X|}
\hline
生产部负责人签字： &
% #VALUE_ID: LEX-P0003-V0001
% #FIELD_VALUE: 生产部负责人签字
% #HANDWRITTEN: 江中伟 2025.10.30
\fieldvalue{\handwritten{江中伟 2025.10.30}} \\
\cline{1-2}
生产负责人签字： &
% #VALUE_ID: LEX-P0003-V0002
% #FIELD_VALUE: 生产负责人签字
% #HANDWRITTEN: 江中伟 2025.02.30
\fieldvalue{\handwritten{江中伟 2025.02.30}} \\
\hline
\end{tabularx}

\begin{tabularx}{\textwidth}{|p{0.11\textwidth}|X|p{0.28\textwidth}|}
\hline
\multirow{4}{*}{检验}
& 取样、检验是否符合相应的操作规程要求。 &
符合规定：
% #VALUE_ID: LEX-P0003-V0003
% #FIELD_VALUE: 取样、检验是否符合相应的操作规程要求——符合规定
\fieldvalue{是}\quad $\square$ 否 \\ \cline{2-3}
& 报告是否按照经批准的质量标准出具。 &
符合规定：
% #VALUE_ID: LEX-P0003-V0004
% #FIELD_VALUE: 报告是否按照经批准的质量标准出具——符合规定
\fieldvalue{是}\quad $\square$ 否 \\ \cline{2-3}
& 检验过程中发生的OOS是否关闭。 &
符合规定：
% #VALUE_ID: LEX-P0003-V0005
% #FIELD_VALUE: 检验过程中发生的OOS是否关闭——符合规定
\fieldvalue{是}\quad $\square$ 否 \\ \cline{2-3}
& 检验方法和检验仪器、设备已经经过验证或确认。 &
符合规定：
% #VALUE_ID: LEX-P0003-V0006
% #FIELD_VALUE: 检验方法和检验仪器、设备已经经过验证或确认——符合规定
\fieldvalue{是}\quad $\square$ 否 \\
\hline
\multirow{3}{*}{偏差}
& \multicolumn{2}{l|}{%
检验过程是否发生偏差；\newline
如发生偏差是否已关闭；\newline
如偏差未关闭，放行前确认是否对产品质量无影响。}
\\[-2pt]
& \multicolumn{2}{l|}{%
% #VALUE_ID: LEX-P0003-V0007
% #FIELD_VALUE: 偏差——是否发生偏差
\fieldvalue{无偏差发生}}\\
& \multicolumn{2}{l|}{$\square$ 有偏差发生，偏差编号为：\rule{3cm}{0.4pt}
\quad $\square$ 偏差已关闭}\\
& \multicolumn{2}{l|}{\rule{3cm}{0.4pt}\quad $\square$ 偏差未关闭}\\
\hline
\multirow{3}{*}{变更}
& \multicolumn{2}{l|}{%
检验过程是否发生变更；\newline
如发生变更是否已关闭；\newline
如变更未关闭，放行前确认是否对产品质量无影响。}
\\[-2pt]
& \multicolumn{2}{l|}{%
% #VALUE_ID: LEX-P0003-V0008
% #FIELD_VALUE: 变更——是否发生变更
\fieldvalue{无变更发生}}\\
& \multicolumn{2}{l|}{$\square$ 有变更发生，变更编号为：\rule{3cm}{0.4pt}
\quad $\square$ 变更已关闭}\\
& \multicolumn{2}{l|}{\rule{3cm}{0.4pt}\quad $\square$ 变更未关闭}\\
\hline
记录完整性检查
& 批检验记录完整，并完成批检验记录审核签批。 &
符合规定：
% #VALUE_ID: LEX-P0003-V0009
% #FIELD_VALUE: 批检验记录完整，并完成批检验记录审核签批——符合规定
\fieldvalue{是}\quad $\square$ 否 \\
\hline
\end{tabularx}

质量控制部经理签字：

% #VALUE_ID: LEX-P0003-V0010
% #FIELD_VALUE: 质量控制部经理签字
% #HANDWRITTEN: 周成 2025.10.30
\fieldvalue{\handwritten{周成 2025.10.30}}

\vspace{1em}

质量控制部负责人签字：

% #VALUE_ID: LEX-P0003-V0011
% #FIELD_VALUE: 质量控制部负责人签字
% #HANDWRITTEN: 孙明军 2025.10.20
\fieldvalue{\handwritten{孙明军 2025.10.20}}

\begin{center}
第 2 页 共 3 页
\end{center}
% LEXOID_PAGE_COMPLETED: 3/59

\newpage

\begin{center}
\small
\begin{tabularx}{\textwidth}{|p{2.1cm}|X|p{3.4cm}|}
\hline
\multicolumn{3}{|c|}{人脉带} \\ \hline
\multirow{2}{*}{人脉带}
& 人脉带是否已正式放行
&
符合规定：
% #VALUE_ID: LEX-P0004-V0001
% #FIELD_VALUE: 人脉带是否已正式放行
\fieldvalue{是}\quad $\square$ 否
\\ \cline{2-3}
& 窗口期病毒检测是否完成，检测结果合格
&
符合规定：
% #VALUE_ID: LEX-P0004-V0002
% #FIELD_VALUE: 窗口期病毒检测是否完成
\fieldvalue{是}\quad $\square$ 否
\\ \hline

质量保证部审核人签字：
&
% #VALUE_ID: LEX-P0004-V0003
% #FIELD_VALUE: 质量保证部审核人签字
% #TODO #HANDWRITTEN: 审核人签字，2025-10-30；签名文字难以辨认
\fieldvalue{\handwritten{审核人签字，2025-10-30}}
& \\ \hline

质量保证部负责人签字：
&
% #VALUE_ID: LEX-P0004-V0004
% #FIELD_VALUE: 质量保证部负责人签字
% #TODO #HANDWRITTEN: 负责人签字，2025-10-30；签名文字难以辨认
\fieldvalue{\handwritten{负责人签字，2025-10-30}}
& \\ \hline

\multirow{2}{*}{偏差、变更}
& 以上偏差、变更是否涵盖本批次所有偏差、变更
&
$\square$ 是，其他偏差、变更另为
\\ \cline{2-3}
&
确认是否所有偏差、变更已关闭或经风险评估对产品质量无影响，不能影响放行
&
$\square$ 是\\
& & $\square$ 否，对产品质量产生影响的偏差、变更号：
% #VALUE_ID: LEX-P0004-V005
% #FIELD_VALUE: 偏差、变更是否已关闭或经风险评估无影响
\fieldvalue{否}
\\ \hline

偏差变更管理员签字：
&
% #VALUE_ID: LEX-P0004-V006
% #FIELD_VALUE: 偏差变更管理员签字
% #TODO #HANDWRITTEN: 管理员签字，2025-10-30；签名文字难以辨认
\fieldvalue{\handwritten{管理员签字，2025-10-30}}
& \\ \hline

\multirow{1}{*}{OOS}
& 确认检验过程中发生的OOS已关闭。
&
$\square$ 是\\
& & $\square$ 否，未关闭的OOS编号：
% #VALUE_ID: LEX-P0004-V007
% #FIELD_VALUE: OOS是否已关闭
\fieldvalue{否}
\\ \hline

OOS管理员签字：
&
% #VALUE_ID: LEX-P0004-V008
% #FIELD_VALUE: OOS管理员签字
% #TODO #HANDWRITTEN: OOS管理员签字，2025-10-30；签名文字难以辨认
\fieldvalue{\handwritten{OOS管理员签字，2025-10-30}}
& \\ \hline

\multirow{3}{*}{记录审核}
& 批生产记录：批生产记录齐全、书写正确、数据完整无空项，所用物料、设备、工艺符合要求。
&
符合规定：
% #VALUE_ID: LEX-P0004-V009
% #FIELD_VALUE: 批生产记录审核
\fieldvalue{是}\quad $\square$ 否
\\ \cline{2-3}
& 批检验记录：批检验记录齐全、书写正确、数据完整无空项；所用物料、设备、工艺符合要求。
&
符合规定：
% #VALUE_ID: LEX-P0004-V010
% #FIELD_VALUE: 批检验记录审核
\fieldvalue{是}\quad $\square$ 否
\\ \cline{2-3}
& 检验报告单：检验项目完整、与检验记录数据一致，报告结果经过质量控制部负责人批准。
&
符合规定：
% #VALUE_ID: LEX-P0004-V011
% #FIELD_VALUE: 检验报告单审核
\fieldvalue{是}\quad $\square$ 否
\\ \hline

质量保证部审核人签字：
&
% #VALUE_ID: LEX-P0004-V012
% #FIELD_VALUE: 质量保证部审核人签字
% #TODO #HANDWRITTEN: 审核人签字，2025-10-30；签名文字难以辨认
\fieldvalue{\handwritten{审核人签字，2025-10-30}}
& \\ \hline

质量保证部负责人签字：
&
% #VALUE_ID: LEX-P0004-V013
% #FIELD_VALUE: 质量保证部负责人签字
% #TODO #HANDWRITTEN: 负责人签字，2025-10-30；签名文字难以辨认
\fieldvalue{\handwritten{负责人签字，2025-10-30}}
& \\ \hline

\multirow{3}{*}{质量评价及放行结论}
& $\square$ 合格，满足放行标准，准予放行
& \multirow{3}{*}{质量受检人批准/日期：} \\ \cline{2-2}
& $\square$ 不合格，不满足放行标准，拒绝放行
& \\ \cline{2-2}
& 
% #VALUE_ID: LEX-P0004-V014
% #FIELD_VALUE: 质量评价及放行结论
\fieldvalue{合格，不满足放行标准，暂停放行}
& \\ \hline
\end{tabularx}
\end{center}

% #VALUE_ID: LEX-P0004-V015
% #FIELD_VALUE: 质量评价及放行结论旁批注
% #TODO #HANDWRITTEN: 关于暂停放行及后续处理的手写批注；内容部分难以辨认
\fieldvalue{\handwritten{关于暂停放行及后续处理的手写批注}}

% #VALUE_ID: LEX-P0004-V016
% #FIELD_VALUE: 质量评价及放行结论旁批注
% #TODO #HANDWRITTEN: 手写补充说明；内容难以辨认
\fieldvalue{\handwritten{手写补充说明}}

% #VALUE_ID: LEX-P0004-V017
% #FIELD_VALUE: 质量评价及放行结论旁批注
% #TODO #HANDWRITTEN: 手写补充说明，2024-10-30；内容难以辨认
\fieldvalue{\handwritten{手写补充说明，2024-10-30}}

% #VALUE_ID: LEX-P0004-V018
% #FIELD_VALUE: 质量受检人批准
% #TODO #HANDWRITTEN: 批准人签字；签名文字难以辨认
\fieldvalue{\handwritten{批准人签字}}

% #VALUE_ID: LEX-P0004-V019
% #FIELD_VALUE: 质量受检人批准日期
% #HANDWRITTEN: 2025-10-30
\fieldvalue{\handwritten{2025-10-30}}
% LEXOID_PAGE_COMPLETED: 4/59

\newpage

\begin{center}
\small
\begin{tabularx}{\textwidth}{@{}lX@{}}
文件编码：&
% #VALUE_ID: LEX-P0005-V0001
% #FIELD_VALUE: 文件编码
\fieldvalue{REC-YZ-SMP-QA-01-007-e}
\\
版本号：&
% #VALUE_ID: LEX-P0005-V0002
% #FIELD_VALUE: 版本号
\fieldvalue{1.5}
\end{tabularx}

\vspace{0.5em}

\textbf{艾米迈托赛注射液暂停放行通知单}

\vspace{0.5em}

\renewcommand{\arraystretch}{1.6}
\begin{tabularx}{\textwidth}{|>{\centering\arraybackslash}p{2.2cm}|X|
>{\centering\arraybackslash}p{2.2cm}|X|}
\hline
产品名称 &
% #VALUE_ID: LEX-P0005-V0003
% #FIELD_VALUE: 产品名称
\fieldvalue{艾米迈托赛注射液}
&
代码 &
% #VALUE_ID: LEX-P0005-V0004
% #FIELD_VALUE: 代码
\fieldvalue{C304}
\\ \hline

规　格 &
% #VALUE_ID: LEX-P0005-V0005
% #FIELD_VALUE: 规格
\fieldvalue{12ml：6.0$\times10^{7}$个细胞}
&
产品批号 &
% #VALUE_ID: LEX-P0005-V0006
% #FIELD_VALUE: 产品批号
\fieldvalue{A33Z2202506001}
\\ \hline

数　量 &
% #VALUE_ID: LEX-P0005-V0007
% #FIELD_VALUE: 数量
\fieldvalue{541袋}
&
检验报告编号 &
% #VALUE_ID: LEX-P0005-V0008
% #FIELD_VALUE: 检验报告编号
\fieldvalue{C30420250807032}
\\ \hline

效　期 &
% #VALUE_ID: LEX-P0005-V0009
% #FIELD_VALUE: 有效期至
\fieldvalue{\fbox{\checkmark}\ 2026年06月17日}
&
\multicolumn{2}{l|}{}
\\ \hline

\multicolumn{4}{|l|}{%
质量评价：
% #VALUE_ID: LEX-P0005-V0010
% #FIELD_VALUE: 质量评价
\fieldvalue{\textcolor{red}{检测合格，但不满足放行标准，暂停放行！}}
\hfill
% #VALUE_ID: LEX-P0005-V0011
% #TODO #HANDWRITTEN: 手写批注，字迹难辨；蓝色手写文字无法可靠辨识
\handwritten{手写批注（字迹难辨）}
}
\\ \hline

质量受权人 &
% #VALUE_ID: LEX-P0005-V0012
% #FIELD_VALUE: 质量受权人
% #HANDWRITTEN: 签名
\fieldvalue{\handwritten{签名}}
&
日　期 &
% #VALUE_ID: LEX-P0005-V0013
% #FIELD_VALUE: 日期
% #HANDWRITTEN: 2025年9月3日
\fieldvalue{\handwritten{2025年9月3日}}
\\ \hline
\end{tabularx}

% TODO: The lower portion of the page contains only the top edge of a partially visible dashed element; its content is not legible.
\end{center}
% LEXOID_PAGE_COMPLETED: 5/59

\newpage

\begin{center}
卓越细胞卓越细胞（北京）有限公司\\[-2pt]
\textit{PLATINUM LIFE SCIENCE BIOTECH}\\[6pt]
艾莱迈托巫注射液审核放行单
\end{center}

\begin{flushright}
文件编码：REC-YZ-SMP-QA-01-007-a\qquad 版本号：1.5
\end{flushright}

\renewcommand{\arraystretch}{1.35}
\small
\begin{tabularx}{\textwidth}{|p{0.11\textwidth}|X|p{0.11\textwidth}|X|}
\hline
产品名称 &
% #VALUE_ID: LEX-P0006-V0001
% #FIELD_VALUE: 产品名称
\fieldvalue{艾莱迈托巫注射液}
& 代码 &
% #VALUE_ID: LEX-P0006-V0002
% #FIELD_VALUE: 代码
\fieldvalue{C304}
\\ \hline
产品批号 &
% #VALUE_ID: LEX-P0006-V0003
% #FIELD_VALUE: 产品批号
\fieldvalue{A33Z2022506001}
& 规格 &
% #VALUE_ID: LEX-P0006-V0004
% #FIELD_VALUE: 规格
\fieldvalue{12ml：$6.0\times 10^{7}$个细胞}
\\ \hline
& & 数量 &
% #VALUE_ID: LEX-P0006-V0005
% #FIELD_VALUE: 数量
\fieldvalue{541袋}
\\ \hline
\multicolumn{1}{|c|}{工序/项目}
& \multicolumn{1}{c|}{审核内容及标准}
& \multicolumn{2}{c|}{结果}
\\ \hline
物料检查
& 生产所用原辅料、耗材、包装材料等均为经放行的合格物料。
& \multicolumn{2}{l|}{符合规定：%
% #VALUE_ID: LEX-P0006-V0006
% #FIELD_VALUE: 物料检查结果
% #HANDWRITTEN: \checkmark
\fieldvalue{\handwritten{\checkmark}} 是\quad $\square$ 否}
\\ \hline
生产
& 生产过程工艺参数是否与工艺规程一致，生产操作是否符合相应的操作规程要求。
& \multicolumn{2}{l|}{符合规定：%
% #VALUE_ID: LEX-P0006-V0007
% #FIELD_VALUE: 生产过程审核结果
% #HANDWRITTEN: \checkmark
\fieldvalue{\handwritten{\checkmark}} 是\quad $\square$ 否}
\\ \cline{2-4}
&
物料平衡是否符合要求。
& \multicolumn{2}{l|}{符合规定：%
% #VALUE_ID: LEX-P0006-V0008
% #FIELD_VALUE: 物料平衡审核结果
% #HANDWRITTEN: \checkmark
\fieldvalue{\handwritten{\checkmark}} 是\quad $\square$ 否}
\\ \cline{2-4}
&
主要生产工艺和生产设备已经过验证/确认。
& \multicolumn{2}{l|}{符合规定：%
% #VALUE_ID: LEX-P0006-V0009
% #FIELD_VALUE: 生产工艺及设备验证结果
% #HANDWRITTEN: \checkmark
\fieldvalue{\handwritten{\checkmark}} 是\quad $\square$ 否}
\\ \cline{2-4}
&
设备状态、清场等均符合要求。
& \multicolumn{2}{l|}{符合规定：%
% #VALUE_ID: LEX-P0006-V0010
% #FIELD_VALUE: 设备状态及清场结果
% #HANDWRITTEN: \checkmark
\fieldvalue{\handwritten{\checkmark}} 是\quad $\square$ 否}
\\ \hline
偏差
& \begin{minipage}[t]{\linewidth}
生产过程是否发生偏差；\\
如发生偏差是否已关闭；\\
如偏差未关闭，放行前确认是否对产品质量有影响。
\end{minipage}
&
\begin{minipage}[t]{\linewidth}
$\square$ 无偏差发生\\[2pt]
% #VALUE_ID: LEX-P0006-V0011
% #FIELD_VALUE: 偏差发生情况
% #HANDWRITTEN: \checkmark
\fieldvalue{\handwritten{\checkmark}} 有偏差发生，\\
偏差编号为：\\[-2pt]
% #VALUE_ID: LEX-P0006-V0012
% #FIELD_VALUE: 偏差编号
% #TODO #HANDWRITTEN: DE20?37；手写内容部分不清晰
\fieldvalue{\handwritten{DE20?37}}\\[-2pt]
% #VALUE_ID: LEX-P0006-V0013
% #FIELD_VALUE: 偏差编号
% #TODO #HANDWRITTEN: DE20?87；手写内容部分不清晰
\fieldvalue{\handwritten{DE20?87}}
\end{minipage}
&
\begin{minipage}[t]{\linewidth}
$\square$ 偏差已关闭\\
% #VALUE_ID: LEX-P0006-V0014
% #FIELD_VALUE: 偏差关闭状态
% #HANDWRITTEN: \checkmark
\fieldvalue{\handwritten{\checkmark}} 偏差未关闭\\
$\square$ 偏差未关闭\\
对产品质量有影响\\
$\square$ 偏差未关闭\\
对产品质量有影响
\end{minipage}
\\ \hline
变更
& \begin{minipage}[t]{\linewidth}
生产过程是否发生变更；\\
如发生变更是否已关闭；\\
如变更未关闭，放行前确认是否对产品质量无影响。
\end{minipage}
&
\begin{minipage}[t]{\linewidth}
% #VALUE_ID: LEX-P0006-V0015
% #FIELD_VALUE: 变更发生情况
% #HANDWRITTEN: \checkmark
\fieldvalue{\handwritten{\checkmark}} 无变更发生\\
$\square$ 有变更发生，\\
变更编号为：\\
\rule{0.8\linewidth}{0.4pt}\\
\rule{0.8\linewidth}{0.4pt}
\end{minipage}
&
\begin{minipage}[t]{\linewidth}
$\square$ 变更已关闭\\
$\square$ 变更未关闭\\
但对产品质量无影响\\
$\square$ 变更未关闭\\
但对产品质量无影响
\end{minipage}
\\ \hline
\begin{minipage}[t]{\linewidth}\centering
记录完整性\\
检查
\end{minipage}
&
批生产记录完整，并完成批生产记录审核签字批准。
& \multicolumn{2}{l|}{符合规定：%
% #VALUE_ID: LEX-P0006-V0016
% #FIELD_VALUE: 记录完整性检查结果
% #HANDWRITTEN: \checkmark
\fieldvalue{\handwritten{\checkmark}} 是\quad $\square$ 否}
\\ \hline
\end{tabularx}

\vfill
\begin{center}
第 1 页 共 3 页
\end{center}
% LEXOID_PAGE_COMPLETED: 6/59

\newpage

\noindent
\begin{tabularx}{\textwidth}{@{}Xr@{}}
\small\textbf{塑料包装用聚乙烯吹塑薄膜}（或类似固定页眉标识）&
文件编码：REC-YZ-SMP-QA-01-007-a \quad 版本号：1.5
\end{tabularx}

\vspace{2pt}

\noindent
\begin{tabularx}{\textwidth}{|p{0.22\textwidth}|X|}
\hline
生产部负责人签字：
&
% #VALUE_ID: LEX-P0007-V0001
% #FIELD_VALUE: 生产部负责人签字
% #TODO #HANDWRITTEN: [签名] 2025.08.20；签名文字难以辨认
\fieldvalue{\handwritten{[签名] 2025.08.20}}\\
\hline
生产负责人签字：
&
% #VALUE_ID: LEX-P0007-V0002
% #FIELD_VALUE: 生产负责人签字
% #TODO #HANDWRITTEN: [签名] 2025.08.20；签名文字难以辨认
\fieldvalue{\handwritten{[签名] 2025.08.20}}\\
\hline
\end{tabularx}

\vspace{2pt}

\renewcommand{\arraystretch}{1.35}
\small
\begin{tabularx}{\textwidth}{|p{0.11\textwidth}|X|p{0.27\textwidth}|}
\hline
\multirow{4}{*}{检验}
&
取样、检验是否符合相应的操作规程要求。
&
符合规定：
% #VALUE_ID: LEX-P0007-V0003
% #FIELD_VALUE: 取样、检验是否符合相应的操作规程要求——符合规定
% #HANDWRITTEN: 是（勾选）
\fieldvalue{\handwritten{是（勾选）}} \quad □否
\\ \cline{2-3}
&
报告是否按照经批准的质量标准出具。
&
符合规定：
% #VALUE_ID: LEX-P0007-V0004
% #FIELD_VALUE: 报告是否按照经批准的质量标准出具——符合规定
% #HANDWRITTEN: 是（勾选）
\fieldvalue{\handwritten{是（勾选）}} \quad □否
\\ \cline{2-3}
&
检验过程发生的OOS是否关闭。
&
符合规定：
% #VALUE_ID: LEX-P0007-V0005
% #FIELD_VALUE: 检验过程发生的OOS是否关闭——符合规定
% #HANDWRITTEN: 是（勾选）
\fieldvalue{\handwritten{是（勾选）}} \quad □否
\\ \cline{2-3}
&
检验方法和检验仪器、设备已经经过验证或确认。
&
符合规定：
% #VALUE_ID: LEX-P0007-V0006
% #FIELD_VALUE: 检验方法和检验仪器、设备已经经过验证或确认——符合规定
% #HANDWRITTEN: 是（勾选）
\fieldvalue{\handwritten{是（勾选）}} \quad □否
\\ \hline

\multirow{1}{*}{偏差}
&
\begin{tabular}[t]{@{}l@{}}
检验过程是否发生偏差；\\
如发生偏差是否已关闭；\\
如偏差未关闭，放行前确认是否对产品质量无影响。
\end{tabular}
&
\begin{tabular}[t]{@{}l@{}}
% #VALUE_ID: LEX-P0007-V0007
% #FIELD_VALUE: 偏差——是否发生偏差
% #HANDWRITTEN: 无偏差发生（勾选）
\fieldvalue{\handwritten{无偏差发生（勾选）}}\\
□有偏差发生，\\
偏差编号为：\rule{2.2cm}{0.4pt}\\
\rule{2.2cm}{0.4pt}\\
\rule{2.2cm}{0.4pt}\\
\rule{2.2cm}{0.4pt}\\
\rule{2.2cm}{0.4pt}
\end{tabular}
\\ \hline

\multirow{1}{*}{变更}
&
\begin{tabular}[t]{@{}l@{}}
检验过程是否发生变更；\\
如发生变更是否已关闭；\\
如变更未关闭，放行前确认是否对产品质量无影响。
\end{tabular}
&
\begin{tabular}[t]{@{}l@{}}
% #VALUE_ID: LEX-P0007-V0008
% #FIELD_VALUE: 变更——是否发生变更
% #HANDWRITTEN: 无变更发生（勾选）
\fieldvalue{\handwritten{无变更发生（勾选）}}\\
□有变更发生，\\
变更编号为：\rule{2.2cm}{0.4pt}\\
\rule{2.2cm}{0.4pt}\\
\rule{2.2cm}{0.4pt}\\
\rule{2.2cm}{0.4pt}\\
\rule{2.2cm}{0.4pt}
\end{tabular}
\\ \hline

记录完整性检查
&
批检验记录完整，并完成批检验记录审核单签批。
&
符合规定：
% #VALUE_ID: LEX-P0007-V0009
% #FIELD_VALUE: 批检验记录完整，并完成批检验记录审核单签批——符合规定
% #HANDWRITTEN: 是（勾选）
\fieldvalue{\handwritten{是（勾选）}} \quad □否
\\ \hline
\end{tabularx}

\vspace{2pt}

\noindent
\begin{tabularx}{\textwidth}{|p{0.22\textwidth}|X|}
\hline
质量控制部经理签字：
&
% #VALUE_ID: LEX-P0007-V0010
% #FIELD_VALUE: 质量控制部经理签字
% #TODO #HANDWRITTEN: [签名] 2025.08.20；签名文字难以辨认
\fieldvalue{\handwritten{[签名] 2025.08.20}}\\
\hline
质量控制部负责人签字：
&
% #VALUE_ID: LEX-P0007-V0011
% #FIELD_VALUE: 质量控制部负责人签字
% #TODO #HANDWRITTEN: [签名] 2025.08.21；签名文字难以辨认
\fieldvalue{\handwritten{[签名] 2025.08.21}}\\
\hline
\end{tabularx}

\vfill
\begin{center}
第 2 页 共 3 页
\end{center}
% LEXOID_PAGE_COMPLETED: 7/59

\newpage

\begin{flushright}
文件编码：REC-YZ-SMR-QA-01-007-a \qquad 版本号：1.5
\end{flushright}

\begin{tabularx}{\textwidth}{|p{2.2cm}|X|p{4.0cm}|}
\hline
\multicolumn{1}{|c|}{人员带} &
人脉带是否已正式放行 &
符合规定：%
% #VALUE_ID: LEX-P0008-V0001
% #FIELD_VALUE: 人脉带是否已正式放行
\fieldvalue{\ensuremath{\checkmark}\,是}\quad $\square$ 否
\\ \hline
&
窗口期病害检测是否完成，检测结果合格 &
符合规定：%
% #VALUE_ID: LEX-P0008-V0002
% #FIELD_VALUE: 窗口期病害检测是否完成，检测结果合格
\fieldvalue{\ensuremath{\checkmark}\,是}\quad $\square$ 否

% #VALUE_ID: LEX-P0008-V0003
% #FIELD_VALUE: 检测结果备注
% #TODO #HANDWRITTEN: 无异常，2025.08.21；部分字迹难以辨认
\fieldvalue{\handwritten{无异常，2025.08.21}}
\\ \hline
\multicolumn{1}{|c|}{质量保证证书} &
质量保证证书审核人签字：%
% #VALUE_ID: LEX-P0008-V0004
% #FIELD_VALUE: 质量保证证书审核人签字
% #HANDWRITTEN: 赵祥宇 2025.08.21
\fieldvalue{\handwritten{赵祥宇\quad 2025.08.21}} &
\\ \cline{2-3}
&
质量保证证书负责人签字：%
% #VALUE_ID: LEX-P0008-V0005
% #FIELD_VALUE: 质量保证证书负责人签字
% #HANDWRITTEN: 陈敏敏 2025.08.21
\fieldvalue{\handwritten{陈敏敏\quad 2025.08.21}} &
\\ \hline
\multicolumn{1}{|c|}{偏差、变更} &
以上偏差、变更是否涵盖本批次所有偏差、变更 &
$\square$ 是\quad $\square$ 否，其他偏差、变更另为

% #VALUE_ID: LEX-P0008-V006
% #FIELD_VALUE: 是否涵盖本批次所有偏差、变更
\fieldvalue{\ensuremath{\checkmark}\,是}
\\ \cline{2-3}
&
确认是否所有偏差、变更已关闭或经风险评估对产品质量无影响，不影响放行 &
$\square$ 是\quad
% #VALUE_ID: LEX-P0008-V007
% #FIELD_VALUE: 偏差、变更是否已关闭或经风险评估对产品质量无影响
\fieldvalue{\ensuremath{\checkmark}\,否}，对产品质量产生影响的偏差、变更另为

% #TODO #HANDWRITTEN: 相关说明及日期，字迹难以辨认
% #VALUE_ID: LEX-P0008-V008
% #FIELD_VALUE: 偏差、变更说明
% #HANDWRITTEN: 相关说明，2025.08.21
\fieldvalue{\handwritten{相关说明\quad 2025.08.21}}
\\ \cline{2-3}
&
偏差变更管理员签字：%
% #VALUE_ID: LEX-P0008-V009
% #FIELD_VALUE: 偏差变更管理员签字
% #TODO #HANDWRITTEN: 杨伟 2025.08.21；后续字迹难以辨认
\fieldvalue{\handwritten{杨伟\quad 2025.08.21\quad [字迹不清]}} &
\\ \hline
\multicolumn{1}{|c|}{OOS} &
确认检验过程发生的 OOS 已关闭。 &
$\square$ 是\quad
% #VALUE_ID: LEX-P0008-V010
% #FIELD_VALUE: OOS 是否已关闭
\fieldvalue{\ensuremath{\checkmark}\,否}，未关闭的 OOS 需按规定处理。

% #TODO #HANDWRITTEN: OOS处理说明及日期，字迹难以辨认
% #VALUE_ID: LEX-P0008-V011
% #FIELD_VALUE: 未关闭OOS处理说明
% #HANDWRITTEN: 已关闭 2025.08.21
\fieldvalue{\handwritten{已关闭\quad 2025.08.21}}
\\ \cline{2-3}
&
OOS 管理员签字：%
% #VALUE_ID: LEX-P0008-V012
% #FIELD_VALUE: OOS管理员签字
% #TODO #HANDWRITTEN: 冯某 2025.08.21；姓名字迹难以辨认
\fieldvalue{\handwritten{冯某\quad 2025.08.21}} &
\\ \hline
\multicolumn{1}{|c|}{记录审核} &
批生产记录：批检验记录齐全、书写正确、数据完整无空项；所用物料、设备、工艺符合要求。 &
符合规定：%
% #VALUE_ID: LEX-P0008-V013
% #FIELD_VALUE: 批生产记录审核结果
\fieldvalue{\ensuremath{\checkmark}\,是}\quad $\square$ 否
\\ \cline{2-3}
&
批检验记录：批检验记录齐全、书写正确、数据完整无空项；所用物料、设备、工艺符合要求。 &
符合规定：%
% #VALUE_ID: LEX-P0008-V014
% #FIELD_VALUE: 批检验记录审核结果
\fieldvalue{\ensuremath{\checkmark}\,是}\quad $\square$ 否
\\ \cline{2-3}
&
检验报告单：检验项目完整、且与检验记录数据一致，报告结果经过质量控制部负责人批准。 &
符合规定：%
% #VALUE_ID: LEX-P0008-V015
% #FIELD_VALUE: 检验报告单审核结果
\fieldvalue{\ensuremath{\checkmark}\,是}\quad $\square$ 否
\\ \hline
&
质量保证部审核人签字：%
% #VALUE_ID: LEX-P0008-V016
% #FIELD_VALUE: 质量保证部审核人签字
% #TODO #HANDWRITTEN: [姓名] 2025.08.28；签名难以辨认
\fieldvalue{\handwritten{[姓名]\quad 2025.08.28}} &
\\ \cline{2-3}
&
质量保证部负责人签字：%
% #VALUE_ID: LEX-P0008-V017
% #FIELD_VALUE: 质量保证部负责人签字
% #TODO #HANDWRITTEN: 陈敏敏 2025.08.28；签名部分难以辨认
\fieldvalue{\handwritten{陈敏敏\quad 2025.08.28}} &
\\ \hline
\multicolumn{1}{|c|}{质量评价及放行结论} &
$\square$ 合格，满足放行标准，准予放行
\newline
$\square$ 不合格，不满足放行标准，拒绝放行
\newline
不合格原因：\underline{\hspace{5cm}}
\newline
% #VALUE_ID: LEX-P0008-V018
% #FIELD_VALUE: 不合格原因
% #HANDWRITTEN: 检验合格，2025.08.28；后续内容难以辨认
\fieldvalue{\handwritten{检验合格，2025.08.28\quad [字迹不清]}}
\newline
% #VALUE_ID: LEX-P0008-V019
% #FIELD_VALUE: 质量评价及放行结论
\fieldvalue{\ensuremath{\checkmark\,\square\text{合格，不满足放行标准，暂停放行}}} &
质量受权人批准/日期：

% #VALUE_ID: LEX-P0008-V020
% #FIELD_VALUE: 质量受权人批准/日期
% #TODO #HANDWRITTEN: [签名] 2025.08.28；签名及日期部分难以辨认
\fieldvalue{\handwritten{[签名]\quad 2025.08.28}}
\\ \hline
\end{tabularx}

\begin{center}
第 3 页 共 3 页
\end{center}
% LEXOID_PAGE_COMPLETED: 8/59

\newpage

\begin{center}
\textbf{\Large 艾米迈托赛注射液（Z2）}

\textbf{\Large 批生产记录}
\end{center}

\vspace{2cm}

\noindent
\begin{tabularx}{\textwidth}{@{}lX@{}}
产品代码：&
% #VALUE_ID: LEX-P0009-V0001
% #FIELD_VALUE: 产品代码
\fieldvalue{C304}
\\[1.2em]
产品名称：&
% #VALUE_ID: LEX-P0009-V0002
% #FIELD_VALUE: 产品名称
\fieldvalue{艾米迈托赛注射液}
\\[1.2em]
规\quad 格：&
% #VALUE_ID: LEX-P0009-V0003
% #FIELD_VALUE: 规格
\fieldvalue{12ml：$6.0\times10^{7}$个细胞}
\\[1.2em]
生产批号：&
% #VALUE_ID: LEX-P0009-V0004
% #FIELD_VALUE: 生产批号
% #TODO #HANDWRITTEN: A3?221?05?01; handwriting unclear
\fieldvalue{\handwritten{A3?221?05?01}}
\\[1.2em]
批\quad 量：&
% #VALUE_ID: LEX-P0009-V0005
% #FIELD_VALUE: 批量
% #HANDWRITTEN: 200--480
\fieldvalue{\handwritten{200--480}}\ \text{袋}
\\[1.2em]
生产日期：&
% #VALUE_ID: LEX-P0009-V0006
% #FIELD_VALUE: 生产日期（起始日期）
% #HANDWRITTEN: 2025年06月09日
\fieldvalue{\handwritten{2025年06月09日}}
\quad 至 \quad
% #VALUE_ID: LEX-P0009-V0007
% #FIELD_VALUE: 生产日期（结束日期）
% #HANDWRITTEN: 2025年06月18日
\fieldvalue{\handwritten{2025年06月18日}}
\end{tabularx}
% LEXOID_PAGE_COMPLETED: 9/59

\newpage

\begin{center}
\textbf{艾米迈托赛注射液（Z2）批生产记录}

记录编号：YZ-BPR-MF-PD-01-020-a \qquad 版本号：1.3
\end{center}

\begin{tabularx}{\textwidth}{|l|X|l|X|l|X|}
\hline
产品名称 & 艾米迈托赛注射液 & 产品代码 & C304 & 实际数量 &
% #VALUE_ID: LEX-P0010-V0001
% #FIELD_VALUE: 实际数量
% #HANDWRITTEN: 5
\fieldvalue{\handwritten{5}} 瓶 \\
\hline
产品批号 &
% #VALUE_ID: LEX-P0010-V0002
% #FIELD_VALUE: 产品批号
% TODO #HANDWRITTEN: A?3?0?0?6?00; handwritten batch number is unclear
\fieldvalue{\handwritten{A?3?0?0?6?00}} &
规格 & 12\,mL: 6.0$\times 10^{7}$个细胞 & & \\
\hline
\end{tabularx}

\begin{center}
\textbf{生产记录审核单}
\end{center}

\small
\begin{tabularx}{\textwidth}{|c|X|l|}
\hline
序号 & 审核内容 & 审核结果 \\
\hline
1 & 批生产记录封面 &
% #VALUE_ID: LEX-P0010-V0003
% #FIELD_VALUE: 序号1审核结果
\fieldvalue{[✓]符合 \quad [ ]不符合} \\
\hline
2 & 批生产记录审核表 &
% #VALUE_ID: LEX-P0010-V0004
% #FIELD_VALUE: 序号2审核结果
\fieldvalue{[✓]有 \quad [ ]无} \\
\hline
3 & 批准的批生产指令【此页放在《批生产记录审核单》后】 &
% #VALUE_ID: LEX-P0010-V0005
% #FIELD_VALUE: 序号3审核结果
\fieldvalue{[✓]有 \quad [ ]无} \\
\hline
\multicolumn{3}{|c|}{物料情况} \\
\hline
4 & P2 代细胞出库申请单 1 张 &
% #VALUE_ID: LEX-P0010-V0006
% #FIELD_VALUE: 序号4审核结果
\fieldvalue{[✓]有 \quad [ ]无} \\
\hline
5 & 艾米迈托赛 P2 代细胞放行复印件 1 张 &
% #VALUE_ID: LEX-P0010-V0007
% #FIELD_VALUE: 序号5审核结果
\fieldvalue{[✓]符合 \quad [ ]不符合} \\
\hline
6 & 细胞入库申请单 1 张 &
% #VALUE_ID: LEX-P0010-V0008
% #FIELD_VALUE: 序号6审核结果
\fieldvalue{[✓]符合 \quad [ ]不符合} \\
\hline
\multicolumn{3}{|c|}{修订历史} \\
\hline
7 & 修订历史 &
% #VALUE_ID: LEX-P0010-V0009
% #FIELD_VALUE: 序号7审核结果
\fieldvalue{[✓]有 \quad [ ]无} \\
\hline
\multicolumn{3}{|c|}{细胞生产记录} \\
\hline
8 & 艾米迈托赛 P2 代细胞库复苏：\ 
% #VALUE_ID: LEX-P0010-V0010
% #FIELD_VALUE: 序号8审核结果
\fieldvalue{[✓]符合 \quad [ ]不符合}
\quad
% #VALUE_ID: LEX-P0010-V0011
% #FIELD_VALUE: P2代细胞库复苏数量
% #HANDWRITTEN: 1
\fieldvalue{\handwritten{1}} 份 &
\\
\hline
9 & 一级生物反应器细胞培养：\ 
% #VALUE_ID: LEX-P0010-V0012
% #FIELD_VALUE: 序号9审核结果
\fieldvalue{[✓]符合 \quad [ ]不符合}
\quad
% #VALUE_ID: LEX-P0010-V0013
% #FIELD_VALUE: 一级生物反应器细胞培养数量
% #HANDWRITTEN: 1
\fieldvalue{\handwritten{1}} 份 &
\\
\hline
10 & 二级生物反应器细胞培养：\ 
% #VALUE_ID: LEX-P0010-V0014
% #FIELD_VALUE: 序号10审核结果
\fieldvalue{[✓]符合 \quad [ ]不符合}
\quad
% #VALUE_ID: LEX-P0010-V0015
% #FIELD_VALUE: 二级生物反应器细胞培养数量
% #HANDWRITTEN: 1
\fieldvalue{\handwritten{1}} 份 &
\\
\hline
11 & 三级生物反应器细胞培养：\ 
% #VALUE_ID: LEX-P0010-V0016
% #FIELD_VALUE: 序号11审核结果
\fieldvalue{[✓]符合 \quad [ ]不符合}
\quad
% #VALUE_ID: LEX-P0010-V0017
% #FIELD_VALUE: 三级生物反应器细胞培养数量
% #HANDWRITTEN: 1
\fieldvalue{\handwritten{1}} 份 &
\\
\hline
12 & 细胞分装：\ 
% #VALUE_ID: LEX-P0010-V0018
% #FIELD_VALUE: 序号12审核结果
\fieldvalue{[✓]符合 \quad [ ]不符合}
\quad
% #VALUE_ID: LEX-P0010-V0019
% #FIELD_VALUE: 细胞分装数量
% #HANDWRITTEN: 1
\fieldvalue{\handwritten{1}} 份 &
\\
\hline
13 & 程控降温：\ 
% #VALUE_ID: LEX-P0010-V0020
% #FIELD_VALUE: 序号13审核结果
\fieldvalue{[✓]符合 \quad [ ]不符合}
\quad
% #VALUE_ID: LEX-P0010-V0021
% #FIELD_VALUE: 程控降温数量
% #HANDWRITTEN: 1
\fieldvalue{\handwritten{1}} 份 &
\\
\hline
\multicolumn{3}{|c|}{偏差及附件} \\
\hline
14 &
本批有无偏差发生及偏差编号是否被记录：\\
\quad $\square$ 未发生\\
\quad
% #VALUE_ID: LEX-P0010-V0022
% #FIELD_VALUE: 序号14偏差发生情况
\fieldvalue{[✓]已发生且已记录}\\
偏差编号及对应页码：\ 
% #VALUE_ID: LEX-P0010-V0023
% #FIELD_VALUE: 偏差编号1
% TODO #HANDWRITTEN: DF? 2025085; handwritten deviation number is unclear
\fieldvalue{\handwritten{DF? 2025085}}
\quad
% #VALUE_ID: LEX-P0010-V0024
% #FIELD_VALUE: 偏差编号1对应页码
% TODO #HANDWRITTEN: 90页; handwritten page number is unclear
\fieldvalue{\handwritten{90页}}\\
\quad
% #VALUE_ID: LEX-P0010-V0025
% #FIELD_VALUE: 偏差编号2
% TODO #HANDWRITTEN: DF? 202508?; handwritten deviation number is unclear
\fieldvalue{\handwritten{DF? 202508?}}
\quad
% #VALUE_ID: LEX-P0010-V0026
% #FIELD_VALUE: 偏差编号2对应页码
% TODO #HANDWRITTEN: 59页; handwritten page number is unclear
\fieldvalue{\handwritten{59页}}
&
\\
\hline
15 &
（1）生产类设备在线电子记录或报告\\
（2）中控检测记录或报告\\
（3）产品通畅信息表：\ 
% #VALUE_ID: LEX-P0010-V0027
% #FIELD_VALUE: 产品通畅信息表数量
% #HANDWRITTEN: 1
\fieldvalue{\handwritten{1}} 份 &
\\
\hline
\end{tabularx}
\normalsize

% TODO: The footer page numbering is omitted as boilerplate/page-number content.
% LEXOID_PAGE_COMPLETED: 10/59

\newpage

\section*{艾米迈托萜注射液（Z2）批生产记录}

\begin{tabularx}{\textwidth}{|p{2.4cm}|X|p{2.4cm}|X|}
\hline
\multicolumn{4}{|c|}{记录编号：YZ-BPR-MF-PD-01-020-a \hspace{1em} 版本号：%
% #VALUE_ID: LEX-P0011-V0001
% #FIELD_VALUE: 版本号
\fieldvalue{1.3}}\\
\hline
产品名称 & 艾米迈托萜注射液 & 产品代码 & C304\\
\hline
产品批号 &
% #VALUE_ID: LEX-P0011-V0002
% #FIELD_VALUE: 产品批号
% #TODO #HANDWRITTEN: A?32?0?50?1; handwriting is unclear
\fieldvalue{\handwritten{A?32?0?50?1}} &
实际数量 &
% #VALUE_ID: LEX-P0011-V0003
% #FIELD_VALUE: 实际数量
% #HANDWRITTEN: 342
\fieldvalue{\handwritten{342}} 袋\\
\hline
\end{tabularx}

\vspace{0.4em}

\begin{tabularx}{\textwidth}{|p{2.4cm}|X|}
\hline
\multicolumn{2}{|l|}{（4）其它：$\square$ 无 \hspace{2em} $\square$ 有}\\
\hline
\multirow{2}{*}{\shortstack{生产\\工艺\\部核\\核}} &
% #VALUE_ID: LEX-P0011-V0004
% #FIELD_VALUE: 符合规定
% #HANDWRITTEN: ✓
\fieldvalue{\handwritten{\checkmark}} 符合规定 \hspace{1em} $\square$ 不符合规定，不符合原因：\underline{\hspace{5cm}}\\
&
生产部审核人/日期：
% #VALUE_ID: LEX-P0011-V0005
% #HANDWRITTEN: 林芳
\handwritten{林芳}
\hspace{2em}
% #VALUE_ID: LEX-P0011-V0006
% #HANDWRITTEN: 2025.06.28
\handwritten{2025.06.28}\\
\cline{2-2}
&
% #VALUE_ID: LEX-P0011-V0007
% #FIELD_VALUE: 符合规定
% #HANDWRITTEN: ✓
\fieldvalue{\handwritten{\checkmark}} 符合规定 \hspace{1em} $\square$ 不符合规定，不符合原因：\underline{\hspace{5cm}}\\
&
生产部审核人/日期：
% #VALUE_ID: LEX-P0011-V0008
% #HANDWRITTEN: 许海军
\handwritten{许海军}
\hspace{2em}
% #VALUE_ID: LEX-P0011-V0009
% #HANDWRITTEN: 2025.06.30
\handwritten{2025.06.30}\\
\cline{2-2}
&
% #VALUE_ID: LEX-P0011-V0010
% #FIELD_VALUE: 符合规定
% #HANDWRITTEN: ✓
\fieldvalue{\handwritten{\checkmark}} 符合规定 \hspace{1em} $\square$ 不符合规定，不符合原因：\underline{\hspace{5cm}}\\
&
生产部负责人/日期：
% #VALUE_ID: LEX-P0011-V0011
% #HANDWRITTEN: 李???
\handwritten{李???}
\hspace{2em}
% #VALUE_ID: LEX-P0011-V0012
% #HANDWRITTEN: 2025.7.02
\handwritten{2025.7.02}\\
\cline{2-2}
&
% #VALUE_ID: LEX-P0011-V0013
% #FIELD_VALUE: 符合规定
% #HANDWRITTEN: ✓
\fieldvalue{\handwritten{\checkmark}} 符合规定 \hspace{1em} $\square$ 不符合规定，不符合原因：\underline{\hspace{5cm}}\\
&
商业化生产负责人/日期：
% #VALUE_ID: LEX-P0011-V0014
% #HANDWRITTEN: 杨???
\handwritten{杨???}
\hspace{2em}
% #VALUE_ID: LEX-P0011-V0015
% #HANDWRITTEN: 2025.08.14
\handwritten{2025.08.14}\\
\hline
\multirow{2}{*}{\shortstack{QA\\审核}} &
% #VALUE_ID: LEX-P0011-V0016
% #FIELD_VALUE: 符合规定
% #HANDWRITTEN: ✓
\fieldvalue{\handwritten{\checkmark}} 符合规定 \hspace{1em} $\square$ 不符合规定，不符合原因：\underline{\hspace{5cm}}\\
&
QA审核人/日期：
% #VALUE_ID: LEX-P0011-V0017
% #HANDWRITTEN: ????
\handwritten{????}
\hspace{2em}
% #VALUE_ID: LEX-P0011-V0018
% #HANDWRITTEN: 2025.08.16
\handwritten{2025.08.16}\\
\cline{2-2}
&
% #VALUE_ID: LEX-P0011-V0019
% #FIELD_VALUE: 符合规定
% #HANDWRITTEN: ✓
\fieldvalue{\handwritten{\checkmark}} 符合规定 \hspace{1em} $\square$ 不符合规定，不符合原因：\underline{\hspace{5cm}}\\
&
质量授权人/日期：
% #VALUE_ID: LEX-P0011-V0020
% #HANDWRITTEN: 韩彦玲
\handwritten{韩彦玲}
\hspace{2em}
% #VALUE_ID: LEX-P0011-V0021
% #HANDWRITTEN: 2025.8.17
\handwritten{2025.8.17}\\
\hline
\end{tabularx}

\vfill

\begin{center}
第 3 页 共 222 页
\end{center}
% LEXOID_PAGE_COMPLETED: 11/59

\newpage

\begin{center}
华熙生物科技股份有限公司\\
\textit{PLANTABLE LIFE SCIENCE CO., LTD.}\\[2mm]
\textbf{艾米迈托赛注射液（Z2）批生产记录}
\end{center}

记录编号：YZ-BPR-MF-PD-01-020-a

\begin{center}
\begin{tabularx}{0.9\textwidth}{|>{\centering\arraybackslash}X|>{\centering\arraybackslash}X|>{\centering\arraybackslash}X|>{\centering\arraybackslash}X|}
\hline
品名 & 规格 & 生产批号 & 批量\\
\hline
% #VALUE_ID: LEX-P0012-V0001
% #FIELD_VALUE: 品名
\fieldvalue{艾米迈托赛注射液} &
% #VALUE_ID: LEX-P0012-V0002
% #FIELD_VALUE: 规格
\fieldvalue{12ml:6.0$\times10^{7}$个细胞} &
% #VALUE_ID: LEX-P0012-V0003
% #FIELD_VALUE: 生产批号
% #TODO #HANDWRITTEN: A3322?05?001; handwriting is partially unclear
\fieldvalue{\handwritten{A3322?05?001}} &
\fieldvalue{}\\
\hline
\end{tabularx}
\end{center}

\begin{center}
\textbf{递交单}
\end{center}

品名：
% #VALUE_ID: LEX-P0012-V0004
% #FIELD_VALUE: 品名
\fieldvalue{艾米迈托赛注射液}
\hfill
规格：
% #VALUE_ID: LEX-P0012-V0005
% #FIELD_VALUE: 规格
\fieldvalue{12ml:6.0$\times10^{7}$个细胞}

生产批号：
% #VALUE_ID: LEX-P0012-V0006
% #FIELD_VALUE: 生产批号
% #TODO #HANDWRITTEN: A3322?20?06/001; handwriting is partially unclear
\fieldvalue{\handwritten{A3322?20?06/001}}
\hfill
生产工序：
% #VALUE_ID: LEX-P0012-V0007
% #FIELD_VALUE: 生产工序
\fieldvalue{艾米迈托赛注射液P2代细胞培养}

\vspace{2mm}

\begin{center}
\begin{tabularx}{0.9\textwidth}{|>{\centering\arraybackslash}p{0.16\textwidth}|X|}
\hline
名称 &
% #VALUE_ID: LEX-P0012-V0008
% #FIELD_VALUE: 名称
\fieldvalue{一级生物反应器细胞接种转移液}\\
\hline
\multirow{3}{*}{信息} &
本次递交接种液信息：\newline
接种液体积：
% #VALUE_ID: LEX-P0012-V0009
% #FIELD_VALUE: 接种液体积
% #TODO #HANDWRITTEN: 10? ml; handwritten numeral is unclear
\fieldvalue{\handwritten{10?}}\ ml\\
&
活细胞总量：
% #VALUE_ID: LEX-P0012-V0010
% #FIELD_VALUE: 活细胞总量
% #HANDWRITTEN: 3.90$\times10^{7}$
\fieldvalue{\handwritten{3.90$\times10^{7}$}}\ 个细胞\\
&
细胞活率：
% #VALUE_ID: LEX-P0012-V0011
% #FIELD_VALUE: 细胞活率
% #HANDWRITTEN: 92.7
\fieldvalue{\handwritten{92.7}}\%\\
\hline
接收工序 &
接收工序：一级生物反应器细胞培养\\
\hline
递交时间 &
% #VALUE_ID: LEX-P0013-V0013
% #FIELD_VALUE: 递交时间
% #HANDWRITTEN: 2025年06月07日08时12分
\fieldvalue{\handwritten{2025年06月07日08时12分}}\\
\hline
递交人/日期 &
% #VALUE_ID: LEX-P0012-V0014
% #FIELD_VALUE: 递交人
% #TODO #HANDWRITTEN: 胡俊旭; handwritten name is partially unclear
\fieldvalue{\handwritten{胡俊旭}}\quad
% #VALUE_ID: LEX-P0012-V0015
% #FIELD_VALUE: 递交日期
% #HANDWRITTEN: 2025.06.09
\fieldvalue{\handwritten{2025.06.09}}\\
\hline
接收人/日期 &
% #VALUE_ID: LEX-P0012-V0016
% #FIELD_VALUE: 接收人
% #TODO #HANDWRITTEN: 张俊; handwritten name is partially unclear
\fieldvalue{\handwritten{张俊}}\quad
% #VALUE_ID: LEX-P0012-V0017
% #FIELD_VALUE: 接收日期
% #HANDWRITTEN: 2025.06.09
\fieldvalue{\handwritten{2025.06.09}}\\
\hline
\end{tabularx}
\end{center}

% TODO: The page contains a stamped area near the upper-right corner; stamp text is treated as boilerplate/stamp content.
% LEXOID_PAGE_COMPLETED: 12/59

\newpage

\begin{center}
朗生卓越生物科技（北京）有限公司\\
\textit{PLATINUM CELL EXCELLENCE BIOTECH}\\[0.4em]
记录编号：
% #VALUE_ID: LEX-P0013-V0001
% #FIELD_VALUE: 记录编号
\fieldvalue{YZ-BPR-MF-PD-01-020-a}
\qquad
版本号：
% #VALUE_ID: LEX-P0013-V0002
% #FIELD_VALUE: 版本号
\fieldvalue{1.3}\\[1.2em]

\section*{艾米迈托赛注射液（Z2）}
\subsection*{批生产记录}
\end{center}

\noindent
工序名称：
% #VALUE_ID: LEX-P0013-V0003
% #FIELD_VALUE: 工序名称
\underline{\fieldvalue{二级生物反应器细胞培养}}
\\[1.2em]

\noindent
规\quad 格：
% #VALUE_ID: LEX-P0013-V0004
% #FIELD_VALUE: 规格
\underline{\fieldvalue{12ml；$6.0\times 10^{7}$个细胞}}
\\[1.2em]

\noindent
生产批号：
% #VALUE_ID: LEX-P0013-V0005
% #FIELD_VALUE: 生产批号
% #TODO #HANDWRITTEN: 4333222050601; handwritten digits are somewhat unclear
\underline{\fieldvalue{\handwritten{4333222050601}}}
\\[1.2em]

\noindent
批\quad 量：
% #VALUE_ID: LEX-P0013-V0006
% #FIELD_VALUE: 批量
% #HANDWRITTEN: 200--480
\underline{\fieldvalue{\handwritten{200--480}}}\,袋
\\[1.2em]

\noindent
生产日期：
% #VALUE_ID: LEX-P0013-V0007
% #FIELD_VALUE: 生产日期（起）
% #HANDWRITTEN: 2025年06月09日
\fieldvalue{\handwritten{2025年06月09日}}
至
% #VALUE_ID: LEX-P0013-V0008
% #FIELD_VALUE: 生产日期（止）
% #HANDWRITTEN: 2025年06月12日
\fieldvalue{\handwritten{2025年06月12日}}
% LEXOID_PAGE_COMPLETED: 13/59

\newpage

\begin{center}
\textbf{艾米迈托塞注射液（Z2）批生产记录}
\end{center}

\noindent
记录编号：YZ-BPR-MF-PD-01-020-a \hfill 版本号：1.3

\begin{center}
\begin{tabularx}{\textwidth}{|X|X|X|X|}
\hline
品名 & 规格 & 生产批号 & 批量\\
\hline
艾米迈托塞注射液 & 12ml；$6.0\times10^{7}$个细胞 &
% #VALUE_ID: LEX-P0014-V0001
% #FIELD_VALUE: 生产批号
% #TODO #HANDWRITTEN: [批号，难辨]; 手写字符不清
\fieldvalue{\handwritten{[批号，难辨]}} &
200--480袋\\
\hline
\end{tabularx}
\end{center}

\noindent\textbf{生产工序二：一般生物反应器细胞培养}

\medskip
\noindent\textbf{1. 一般生物反应器培养前准备}

\medskip
\noindent\textbf{1.1 生产前确认：}

\small
\begin{tabularx}{\textwidth}{|p{0.55\textwidth}|X|}
\hline
\multicolumn{1}{|c|}{项目} & \multicolumn{1}{c|}{操作记录}\\
\hline
操作人员按《人员进出洁净生产区管理规程》要求进入洁净区生产车间：
&
操作员名称及编码：
% #VALUE_ID: LEX-P0014-V0002
% #FIELD_VALUE: 操作员名称及编码
% #TODO #HANDWRITTEN: [姓名及编码，难辨]; 手写内容不清
\fieldvalue{\handwritten{[姓名及编码，难辨]}}\\
\hline
确认生产区操作间在清场有效期内（C级日清洁周期3天）：
&
是否在清洁有效期内：

% #VALUE_ID: LEX-P0014-V0003
% #FIELD_VALUE: 是否在清洁有效期内
% #HANDWRITTEN: 是
\fieldvalue{\handwritten{是}}

清洁有效期：
% #VALUE_ID: LEX-P0014-V0004
% #FIELD_VALUE: 清洁有效期
% #HANDWRITTEN: 2025.06.10
\fieldvalue{\handwritten{2025.06.10}}\\
\hline
确认操作间温度18--26℃、湿度30\%--75\%、相对压差$>0$Pa（操作间对C级走廊）：
&
温度：
% #VALUE_ID: LEX-P0014-V0005
% #FIELD_VALUE: 温度
% #HANDWRITTEN: 21.1
\fieldvalue{\handwritten{21.1}}℃

湿度：
% #VALUE_ID: LEX-P0014-V0006
% #FIELD_VALUE: 湿度
% #HANDWRITTEN: 52.5
\fieldvalue{\handwritten{52.5}}\%

压差：
% #VALUE_ID: LEX-P0014-V0007
% #FIELD_VALUE: 压差
% #HANDWRITTEN: 12
\fieldvalue{\handwritten{12}} Pa

是否符合要求：

% #VALUE_ID: LEX-P0014-V0008
% #FIELD_VALUE: 是否符合要求
% #HANDWRITTEN: 是
\fieldvalue{\handwritten{是}} \quad 否\\
\hline
确认操作间内无上批产品遗留物、文件及与本产品生产无关物料：
&
% #VALUE_ID: LEX-P0014-V0009
% #FIELD_VALUE: 操作间内无遗留物料
% #HANDWRITTEN: 是
\fieldvalue{\handwritten{是}} \quad 否\\
\hline
检查存放的批记录及其它相关配套文件均与生产指令相对应，且均受控：
&
% #VALUE_ID: LEX-P0014-V0010
% #FIELD_VALUE: 批记录及配套文件受控
% #HANDWRITTEN: 是
\fieldvalue{\handwritten{是}} \quad 否\\
\hline
确认设备完好，且在验证有效期内，计量器具有“合格证”。且在有效期内，详见《主要设备/仪表确认表》：
&
% #VALUE_ID: LEX-P0014-V0011
% #FIELD_VALUE: 设备及计量器具确认
% #HANDWRITTEN: 是
\fieldvalue{\handwritten{是}} \quad 否\\
\hline
反应器氮源压力应控制在如下范围：压缩空气进气压力0.05--0.15MPa，二氧化碳进气压力0.05--0.15MPa，氧气进气压力为0.05--0.15MPa：
&
压缩空气：
% #VALUE_ID: LEX-P0014-V0012
% #FIELD_VALUE: 压缩空气压力
% #HANDWRITTEN: 0.1
\fieldvalue{\handwritten{0.1}} MPa

二氧化碳：
% #VALUE_ID: LEX-P0014-V0013
% #FIELD_VALUE: 二氧化碳压力
% #HANDWRITTEN: 0.7
\fieldvalue{\handwritten{0.7}} MPa

氧气：
% #VALUE_ID: LEX-P0014-V0014
% #FIELD_VALUE: 氧气压力
% #HANDWRITTEN: 0.1
\fieldvalue{\handwritten{0.1}} MPa

是否符合要求：

% #VALUE_ID: LEX-P0014-V0015
% #FIELD_VALUE: 气体压力是否符合要求
% #HANDWRITTEN: 是
\fieldvalue{\handwritten{是}} \quad 否\\
\hline
将操作间状态标识更换为“生产加工状态卡”，写明品名、规格、批号等，并确认与本批生产指令一致：
&
% #VALUE_ID: LEX-P0014-V0016
% #FIELD_VALUE: 生产加工状态卡确认
% #HANDWRITTEN: 是
\fieldvalue{\handwritten{是}} \quad 否\\
\hline
核对本次生产用物料耗材、溶液和识别的品名、规格、批号、数量，检查物料外包装的完整性以及是否在有效期/复验期内，并记录相关信息，详见《物料确认表》和《溶液确认表》：
&
% #VALUE_ID: LEX-P0014-V0017
% #FIELD_VALUE: 物料及溶液确认
% #HANDWRITTEN: 是
\fieldvalue{\handwritten{是}} \quad 否\\
\hline
备注：
&
% #VALUE_ID: LEX-P0014-V0018
% #FIELD_VALUE: 备注
% #HANDWRITTEN: 勾
\fieldvalue{\handwritten{勾}}\\
\hline
\end{tabularx}

\normalsize
\medskip

\noindent
操作人/日期：
% #VALUE_ID: LEX-P0014-V0019
% #FIELD_VALUE: 操作人签名
% #TODO #HANDWRITTEN: [签名，难辨]; 手写签名无法准确辨认
\fieldvalue{\handwritten{[签名，难辨]}}
\quad
% #VALUE_ID: LEX-P0014-V0020
% #FIELD_VALUE: 操作日期
% #HANDWRITTEN: 2025.06.09
\fieldvalue{\handwritten{2025.06.09}}

\hfill
复核人/日期：
% #VALUE_ID: LEX-P0014-V0021
% #FIELD_VALUE: 复核人签名
% #TODO #HANDWRITTEN: [签名，难辨]; 手写签名无法准确辨认
\fieldvalue{\handwritten{[签名，难辨]}}
\quad
% #VALUE_ID: LEX-P0014-V0022
% #FIELD_VALUE: 复核日期
% #HANDWRITTEN: 2025.06.09
\fieldvalue{\handwritten{2025.06.09}}

\begin{center}
第17页 共222页
\end{center}
% LEXOID_PAGE_COMPLETED: 14/59

\newpage

\section*{艾米迈托囊注射液（Z2）批生产记录}

\begin{tabularx}{\textwidth}{|c|c|c|c|}
\hline
品名 & 规格 & 生产批号 & 批量 \\
\hline
艾米迈托囊注射液 & 12ml；6.0×10$^{7}$个细胞 & 
% #VALUE_ID: LEX-P0015-V0001
% #FIELD_VALUE: 生产批号
% #HANDWRITTEN: 443221250100
\fieldvalue{\handwritten{443221250100}} 
& 200--480袋 \\
\hline
\end{tabularx}

\vspace{0.5em}

1.1.1 主要设备/仪表确认表

\small
\begin{tabularx}{\textwidth}{|c|X|c|c|X|c|}
\hline
序号 & 设备名称 & 品牌/型号 & 设备编号 & 检查项 & 确认结果 \\
\hline
1 & 生物反应器控制台（主机） & 华龛生物 &
% #VALUE_ID: LEX-P0015-V0002
% #FIELD_VALUE: 设备编号（1）
% #HANDWRITTEN: 1123610
\fieldvalue{\handwritten{1123610}} &
\multirow{7}{=}{\parbox{0.22\textwidth}{1.设备应有“完好”标识；\\
2.设备在验证有效期内；\\
3.计量器具有“合格证”且在有效期内。}} &
% #VALUE_ID: LEX-P0015-V0003
% #FIELD_VALUE: 确认结果（1）
\fieldvalue{是} \quad $\square$ 否 \\
\hline
2 & 生物反应器罐体 & 华龛生物 &
% #VALUE_ID: LEX-P0015-V0004
% #FIELD_VALUE: 设备编号（2）
% #HANDWRITTEN: V56T05202-220708
\fieldvalue{\handwritten{V56T05202-220708}} &
& 
% #VALUE_ID: LEX-P0015-V0005
% #FIELD_VALUE: 确认结果（2）
\fieldvalue{是} \quad $\square$ 否 \\
\hline
3 & 接管机 & 上海英启 &
% #VALUE_ID: LEX-P0015-V0006
% #FIELD_VALUE: 设备编号（3）
% #HANDWRITTEN: 1123610
\fieldvalue{\handwritten{1123610}} &
&
% #VALUE_ID: LEX-P0015-V0007
% #FIELD_VALUE: 确认结果（3）
\fieldvalue{是} \quad $\square$ 否 \\
\hline
4 & 封管机 & 上海英启 &
% #VALUE_ID: LEX-P0015-V0008
% #FIELD_VALUE: 设备编号（4）
% #HANDWRITTEN: 1123612
\fieldvalue{\handwritten{1123612}} &
&
% #VALUE_ID: LEX-P0015-V0009
% #FIELD_VALUE: 确认结果（4）
\fieldvalue{是} \quad $\square$ 否 \\
\hline
5 & & & & &
$\square$ 是 \quad $\square$ 否 \\
\hline
6 & & & & &
% #VALUE_ID: LEX-P0015-V0010
% #FIELD_VALUE: 确认结果（6）
\fieldvalue{是} \quad $\square$ 否 \\
\hline
7 & & & & &
% #VALUE_ID: LEX-P0015-V0011
% #FIELD_VALUE: 确认结果（7）
\fieldvalue{是} \quad $\square$ 否 \\
\hline
\end{tabularx}
\normalsize

% #VALUE_ID: LEX-P0015-V0012
% #HANDWRITTEN: 核对 2025.06.09
\handwritten{核对 2025.06.09}

备注：\rule{0.75\textwidth}{0.4pt}

确认人/日期：
% #VALUE_ID: LEX-P0015-V0013
% #FIELD_VALUE: 确认人
% #HANDWRITTEN: 秦梅
\fieldvalue{\handwritten{秦梅}}
\quad
% #VALUE_ID: LEX-P0015-V0014
% #FIELD_VALUE: 确认日期
% #HANDWRITTEN: 2025.06.09
\fieldvalue{\handwritten{2025.06.09}}

\hfill
复核人/日期：
% #VALUE_ID: LEX-P0015-V0015
% #FIELD_VALUE: 复核人
% #HANDWRITTEN: 李琦
\fieldvalue{\handwritten{李琦}}
\quad
% #VALUE_ID: LEX-P0015-V0016
% #FIELD_VALUE: 复核日期
% #HANDWRITTEN: 25-06-9
\fieldvalue{\handwritten{25-06-9}}

\vspace{0.5em}

1.1.2 物料确认表

\small
\begin{tabularx}{\textwidth}{|c|X|X|c|X|c|c|}
\hline
序号 & 名称 & 品牌 & 规格 & 批号 & 数量 & 是否符合要求 \\
\hline
1 & 一次性无菌输液袋 &
% #VALUE_ID: LEX-P0015-V0017
% #FIELD_VALUE: 品牌（1）
% #HANDWRITTEN: 黎华生物
\fieldvalue{\handwritten{黎华生物}} &
1个/包 &
% #VALUE_ID: LEX-P0015-V0018
% #FIELD_VALUE: 批号（1）
% #HANDWRITTEN: H040125050601
\fieldvalue{\handwritten{H040125050601}} &
% #VALUE_ID: LEX-P0015-V0019
% #FIELD_VALUE: 数量（1）
% #HANDWRITTEN: 1
\fieldvalue{\handwritten{1}} &
% #VALUE_ID: LEX-P0015-V0020
% #FIELD_VALUE: 是否符合要求（1）
\fieldvalue{是} \quad $\square$ 否 \\
\hline
2 &
% #VALUE_ID: LEX-P0015-V0021
% #FIELD_VALUE: 名称（2）
% #HANDWRITTEN: 3L一次性无菌输液袋
\fieldvalue{\handwritten{3L一次性无菌输液袋}} &
% #VALUE_ID: LEX-P0015-V0022
% #FIELD_VALUE: 品牌（2）
% #HANDWRITTEN: 华龛
\fieldvalue{\handwritten{华龛}} &
% #VALUE_ID: LEX-P0015-V0023
% #FIELD_VALUE: 规格（2）
% #HANDWRITTEN: 1个/包
\fieldvalue{\handwritten{1个/包}} &
% #VALUE_ID: LEX-P0015-V0024
% #FIELD_VALUE: 批号（2）
% #HANDWRITTEN: H040125050501
\fieldvalue{\handwritten{H040125050501}} &
% #VALUE_ID: LEX-P0015-V0025
% #FIELD_VALUE: 数量（2）
% #HANDWRITTEN: 5
\fieldvalue{\handwritten{5}} &
% #VALUE_ID: LEX-P0015-V0026
% #FIELD_VALUE: 是否符合要求（2）
\fieldvalue{是} \quad $\square$ 否 \\
\hline
3 & & & & & & $\square$ 是 \quad $\square$ 否 \\
\hline
\end{tabularx}
\normalsize

% #VALUE_ID: LEX-P0015-V0027
% #HANDWRITTEN: 核梅 2025.06.09
\handwritten{核梅 2025.06.09}

备注：
% #VALUE_ID: LEX-P0015-V0028
% #HANDWRITTEN: ✓
\handwritten{✓}

确认人/日期：
% #VALUE_ID: LEX-P0015-V0029
% #FIELD_VALUE: 确认人
% #HANDWRITTEN: 杜梅
\fieldvalue{\handwritten{杜梅}}
\quad
% #VALUE_ID: LEX-P0015-V0030
% #FIELD_VALUE: 确认日期
% #HANDWRITTEN: 2025.06.09
\fieldvalue{\handwritten{2025.06.09}}

\hfill
复核人/日期：
% #VALUE_ID: LEX-P0015-V0031
% #FIELD_VALUE: 复核人
% #HANDWRITTEN: 李琦
\fieldvalue{\handwritten{李琦}}
\quad
% #VALUE_ID: LEX-P0015-V0032
% #FIELD_VALUE: 复核日期
% #HANDWRITTEN: 25-06-9
\fieldvalue{\handwritten{25-06-9}}

\vspace{0.5em}

1.1.3 容器确认表

% TODO: The remainder of this section is not visible on this page.

\begin{center}
第 18 页 共 222 页
\end{center}
% LEXOID_PAGE_COMPLETED: 15/59

\newpage

\begin{center}
\textbf{艾米迈托囊注射液（Z2）批生产记录}
\end{center}

\noindent
原件\hfill 记录编号：YZ-BPR-MF-PD-01-020-a\hfill 版本号：13

\begin{tabularx}{\textwidth}{|X|X|X|X|}
\hline
品名 & 规格 & 生产批号 & 批量\\
\hline
艾米迈托囊注射液 & 12ml：6.0$\times 10^7$个细胞 &
% #VALUE_ID: LEX-P0016-V0001
% #FIELD_VALUE: 生产批号
% #TODO #HANDWRITTEN: 432?220?06; 字迹较模糊
\fieldvalue{\handwritten{432?220?06}} & 200--480 袋\\
\hline
\end{tabularx}

\begin{tabularx}{\textwidth}{|c|X|X|X|c|X|c|}
\hline
序号 & 名称 & 规格 & 批号 & 数量 & 有效期 & 是否符合要求\\
\hline
1 & 细胞培养基 & 1.2L/袋 &
% #VALUE_ID: LEX-P0016-V0002
% #FIELD_VALUE: 第1项批号
% #HANDWRITTEN: 2025060701
\fieldvalue{\handwritten{2025060701}} &
% #VALUE_ID: LEX-P0016-V0003
% #FIELD_VALUE: 第1项数量
% #HANDWRITTEN: 1
\fieldvalue{\handwritten{1}} &
% #VALUE_ID: LEX-P0016-V0004
% #FIELD_VALUE: 第1项有效期
% #HANDWRITTEN: 2025.06.07
\fieldvalue{\handwritten{2025.06.07}} &
% #VALUE_ID: LEX-P0016-V0005
% #FIELD_VALUE: 第1项是否符合要求
\fieldvalue{是（已勾选）}\\
\hline
2 & 微载体悬液 & 2.66g，$\geq$300ml/瓶 &
% #VALUE_ID: LEX-P0016-V0006
% #FIELD_VALUE: 第2项批号
% #HANDWRITTEN: 2025060701
\fieldvalue{\handwritten{2025060701}} &
% #VALUE_ID: LEX-P0016-V0007
% #FIELD_VALUE: 第2项数量
% #HANDWRITTEN: 1
\fieldvalue{\handwritten{1}} &
% #VALUE_ID: LEX-P0016-V0008
% #FIELD_VALUE: 第2项有效期
% #HANDWRITTEN: 2025.06.14
\fieldvalue{\handwritten{2025.06.14}} &
% #VALUE_ID: LEX-P0016-V0009
% #FIELD_VALUE: 第2项是否符合要求
\fieldvalue{是（已勾选）}\\
\hline
3 & & &
% #VALUE_ID: LEX-P0016-V0010
% #FIELD_VALUE: 第3项批号
% #HANDWRITTEN: 核销 2025.06.09
\fieldvalue{\handwritten{核销 2025.06.09}} & & & \\
\hline
\multicolumn{7}{|l|}{备注：\hspace{0.8\textwidth}}\\
\hline
\end{tabularx}

\noindent
确认人/日期：
% #VALUE_ID: LEX-P0016-V0011
% #FIELD_VALUE: 确认人
% #TODO #HANDWRITTEN: 张梅?; 姓名笔迹不清
\fieldvalue{\handwritten{张梅?}}
\hfill
% #VALUE_ID: LEX-P0016-V0012
% #FIELD_VALUE: 确认日期
% #HANDWRITTEN: 2025.06.09
\fieldvalue{\handwritten{2025.06.09}}

\noindent
复核人/日期：
% #VALUE_ID: LEX-P0016-V0013
% #FIELD_VALUE: 复核人
% #TODO #HANDWRITTEN: 李莉?; 姓名笔迹不清
\fieldvalue{\handwritten{李莉?}}
\hfill
% #VALUE_ID: LEX-P0016-V0014
% #FIELD_VALUE: 复核日期
% #HANDWRITTEN: 2025.06.09
\fieldvalue{\handwritten{2025.06.09}}

\medskip
\noindent
\textbf{1.2 一级生物反应器接种前准备：}

\begin{tabularx}{\textwidth}{|p{0.58\textwidth}|X|}
\hline
\multicolumn{1}{|c|}{项目} & \multicolumn{1}{c|}{操作记录}\\
\hline
1）罐体确认：

确认一级生物反应器罐体已处理完毕，灭菌合格，密闭完整，无菌菌霉滤器、查看控制主机上的历史曲线，确认准备的罐体相关参数控制稳定无异常，校准 DO 电极。 &
罐体是否符合要求：

% #VALUE_ID: LEX-P0016-V0015
% #FIELD_VALUE: 罐体是否符合要求
\fieldvalue{是（已勾选）}

DO 电极校准时间：

% #VALUE_ID: LEX-P0016-V0016
% #FIELD_VALUE: DO电极校准时间
% #HANDWRITTEN: 06：33
\fieldvalue{\handwritten{06：33}}\\
\hline
2）排出一级生物反应器罐中 0.9\%氯化钠注射液：

关闭搅拌和温度控制，使用接管机将一次性无菌储液袋（3L）袋子与一级生物反应器罐体主管路连接，降速向 0.9\%氯化钠注射液全部排出至储液袋内；然后使用封管机将储液袋热合封口，脱离储罐后，留存备用。 &
搅拌和温度控制关闭时间：

% #VALUE_ID: LEX-P0016-V0017
% #FIELD_VALUE: 搅拌和温度控制关闭时间
% #HANDWRITTEN: 06：33
\fieldvalue{\handwritten{06：33}}

排液截止时间：

% #VALUE_ID: LEX-P0016-V0018
% #FIELD_VALUE: 排液截止时间
% #HANDWRITTEN: 06：33--06：43
\fieldvalue{\handwritten{06：33--06：43}}\\
\hline
3）细胞培养基导入：

将细胞培养基（1.2L/袋）通过接管机连接到罐体主管路，按照补液方向将硅胶管安装到蠕动泵头上，打开，开启蠕动泵将培养基全部泵入罐体内。 &
培养基导入时间：

% #VALUE_ID: LEX-P0016-V0019
% #FIELD_VALUE: 培养基导入时间
% #HANDWRITTEN: 06：50--06：53
\fieldvalue{\handwritten{06：50--06：53}}\\
\hline
4）微载体导入：

将浸泡完成的微载体（2.66g 瓶）通过接管机连接到罐体进出液管路，按照补液方向将硅胶管安装到蠕动泵头上，确认管路上打开，并启动蠕动泵将微载体全部泵入罐体内。 &
导入时间：

% #VALUE_ID: LEX-P0016-V0020
% #FIELD_VALUE: 微载体导入时间
% #HANDWRITTEN: 06：54--07：03
\fieldvalue{\handwritten{06：54--07：03}}

瓶清洗次数：

% #VALUE_ID: LEX-P0016-V0021
% #FIELD_VALUE: 瓶清洗次数
% #HANDWRITTEN: 1
\fieldvalue{\handwritten{1}} 次，

% #VALUE_ID: LEX-P0016-V0022
% #FIELD_VALUE: 瓶清洗体积
% #HANDWRITTEN: 60 ml/次
\fieldvalue{\handwritten{60 ml/次}}\\
\hline
\end{tabularx}

\noindent\hfill 第 19 页 共 222 页
% LEXOID_PAGE_COMPLETED: 16/59

\newpage

\begin{center}
\textbf{艾米迈托赛注射液（Z2）批生产记录}
\end{center}

\begin{tabularx}{\textwidth}{|p{0.18\textwidth}|X|p{0.18\textwidth}|X|}
\hline
品名 & 艾米迈托赛注射液 & 规格 & $2\,\mathrm{ml}\colon 6.0\times10^{9}$ 个细胞 \\
\hline
生产批号 &
% #VALUE_ID: LEX-P0017-V0001
% #FIELD_VALUE: 生产批号
% #TODO #HANDWRITTEN: 143322……; 手写数字辨识不清
\fieldvalue{\handwritten{143322……}} &
批量 & 200--480 支 \\
\hline
\end{tabularx}

\vspace{0.2em}

\begin{tabularx}{\textwidth}{|p{0.56\textwidth}|X|}
\hline
\parbox[t]{\linewidth}{
路上管夹打开，开启蠕动泵将微载体全部泵入罐中；微载体\\
泵入后，通过进出液管路回少量培养基对微载体进行清\\
洗，将清洗后的培养基泵液一并导入罐体内，确认罐体最终\\
定容体积达到 $1.5\pm0.3\ \mathrm{L}$ 范围内。
}
&
\parbox[t]{\linewidth}{
容器体积：\rule{0.35\linewidth}{0.4pt}
}
\\
\hline
\parbox[t]{\linewidth}{
5）参数设置开启确认：\\[0.3em]
确认设置开启后如下参数：转速 35rpm，表通流量 20ml/min，\\
底流量 10ml/min，温度 37.0℃，pH 值 7.10；
}
&
\parbox[t]{\linewidth}{
参数开始时间：
% #VALUE_ID: LEX-P0017-V0002
% #FIELD_VALUE: 参数开始时间
% #HANDWRITTEN: 2：02
\fieldvalue{\handwritten{2：02}}\\[0.3em]
参数设置值：\\
转速：
% #VALUE_ID: LEX-P0017-V0003
% #FIELD_VALUE: 转速
% #HANDWRITTEN: 35
\fieldvalue{\handwritten{35}} rpm\\
温度：
% #VALUE_ID: LEX-P0017-V0004
% #FIELD_VALUE: 温度
% #HANDWRITTEN: 37.0
\fieldvalue{\handwritten{37.0}} $^\circ\mathrm{C}$\\
pH：
% #VALUE_ID: LEX-P0017-V0005
% #FIELD_VALUE: pH
% #HANDWRITTEN: 7.10
\fieldvalue{\handwritten{7.10}}\\
表通流量：
% #VALUE_ID: LEX-P0017-V0006
% #FIELD_VALUE: 表通流量
% #HANDWRITTEN: 20
\fieldvalue{\handwritten{20}} ml/min\\
底流量：
% #VALUE_ID: LEX-P0017-V0007
% #FIELD_VALUE: 底流量
% #HANDWRITTEN: 10
\fieldvalue{\handwritten{10}} ml/min
}
\\
\hline
\parbox[t]{\linewidth}{
6）温度达到 37.0$\pm$2.0℃范围内，取样 20ml 左右，进行离线\\
pH 检测，差异值$>$0.1 时，应对在线电极做单点校准。
}
&
\parbox[t]{\linewidth}{
取样时间：
% #VALUE_ID: LEX-P0017-V0008
% #FIELD_VALUE: 取样时间
% #HANDWRITTEN: 07：41
\fieldvalue{\handwritten{07：41}}\\
取样时温度：
% #VALUE_ID: LEX-P0017-V0009
% #FIELD_VALUE: 取样时温度
% #HANDWRITTEN: 37.1
\fieldvalue{\handwritten{37.1}}\\
在线 pH：
% #VALUE_ID: LEX-P0017-V0010
% #FIELD_VALUE: 在线 pH
% #HANDWRITTEN: 7.35
\fieldvalue{\handwritten{7.35}}\\
离线 pH：
% #VALUE_ID: LEX-P0017-V0011
% #FIELD_VALUE: 离线 pH
% #HANDWRITTEN: 7.17
\fieldvalue{\handwritten{7.17}}\\
电极是否校准：$\square$ 是
% #VALUE_ID: LEX-P0017-V0012
% #FIELD_VALUE: 电极是否校准
% #HANDWRITTEN: √ 是
\fieldvalue{\handwritten{√\ 是}}\\
\hspace{4em}$\square$ 否\\
校准后参数：
% #VALUE_ID: LEX-P0017-V0013
% #FIELD_VALUE: 校准后参数
% #HANDWRITTEN: 7.26
\fieldvalue{\handwritten{7.26}}
}
\\
\hline
备注：
% #VALUE_ID: LEX-P0017-V0014
% #FIELD_VALUE: 备注
% #TODO #HANDWRITTEN: ／; 备注栏中的手写符号含义不明确
\fieldvalue{\handwritten{／}}
&
\\
\hline
\end{tabularx}

\vspace{0.4em}

\noindent
确认人/日期：
% #VALUE_ID: LEX-P0017-V0015
% #FIELD_VALUE: 确认人
% #TODO #HANDWRITTEN: 方……; 签名辨识不清
\fieldvalue{\handwritten{方……}}
\quad
% #VALUE_ID: LEX-P0017-V0016
% #FIELD_VALUE: 确认日期
% #HANDWRITTEN: 2025.06.09
\fieldvalue{\handwritten{2025.06.09}}
\hfill
复核人/日期：
% #VALUE_ID: LEX-P0017-V0017
% #FIELD_VALUE: 复核人
% #TODO #HANDWRITTEN: 张梅; 手写签名辨识不完全确定
\fieldvalue{\handwritten{张梅}}
\quad
% #VALUE_ID: LEX-P0017-V0018
% #FIELD_VALUE: 复核日期
% #HANDWRITTEN: 2025.06.09
\fieldvalue{\handwritten{2025.06.09}}
% LEXOID_PAGE_COMPLETED: 17/59

\newpage

\begin{center}
\textbf{艾米迈托霉注射液批生产指令}
\end{center}

\begin{tabularx}{\textwidth}{|p{2.2cm}|p{2.7cm}|X|p{2.7cm}|X|}
\hline
\multicolumn{2}{|c|}{产品名称} &
% #VALUE_ID: LEX-P0018-V0001
% #FIELD_VALUE: 产品名称
\fieldvalue{艾米迈托霉注射液} &
\multicolumn{1}{c|}{} & \\
\hline
\multicolumn{2}{|c|}{生产批号} &
% #VALUE_ID: LEX-P0018-V0002
% #FIELD_VALUE: 生产批号
% #TODO #HANDWRITTEN: A3322?20?3601; handwriting unclear
\fieldvalue{\handwritten{A3322?20?3601}} &
生产代码 &
% #VALUE_ID: LEX-P0018-V0003
% #FIELD_VALUE: 产品代码
% #HANDWRITTEN: C304
\fieldvalue{\handwritten{C304}} \\
\hline
\multicolumn{2}{|c|}{生产线} &
% #VALUE_ID: LEX-P0018-V0004
% #FIELD_VALUE: 生产线
% #TODO #HANDWRITTEN: 脉腾?间示?产品?生产线; handwriting unclear
\fieldvalue{\handwritten{脉腾?间示?产品?生产线}} &
计划产量 &
% #VALUE_ID: LEX-P0018-V0005
% #FIELD_VALUE: 计划产量
% #TODO #HANDWRITTEN: 200--408?袋; handwriting unclear
\fieldvalue{\handwritten{200--408?袋}} \\
\hline
\multicolumn{2}{|c|}{生产起始日期} &
% #VALUE_ID: LEX-P0018-V0006
% #FIELD_VALUE: 生产起始日期
% #TODO #HANDWRITTEN: 20?年7月9日; year unclear
\fieldvalue{\handwritten{20?年7月9日}} &
\multicolumn{2}{c|}{} \\
\hline
\multicolumn{2}{|c|}{\multirow{2}{*}{基本信息}} &
\multicolumn{3}{l|}{报料艾米迈托霉注射液 P2 代细胞库信息：} \\
\cline{3-5}
\multicolumn{2}{|c|}{} &
批号：
% #VALUE_ID: LEX-P0018-V0007
% #FIELD_VALUE: P2代细胞库批号
% #TODO #HANDWRITTEN: A332?20?0909; handwriting unclear
\fieldvalue{\handwritten{A332?20?0909}}
&
代码：
% #VALUE_ID: LEX-P0018-V0008
% #FIELD_VALUE: P2代细胞库代码
% #HANDWRITTEN: C103
\fieldvalue{\handwritten{C103}}
\quad
数量：
% #VALUE_ID: LEX-P0018-V0009
% #FIELD_VALUE: P2代细胞库数量
% #HANDWRITTEN: 1袋
\fieldvalue{\handwritten{1袋}} \\
\hline
\multicolumn{2}{|c|}{执行工艺规程} &
编码：
% #VALUE_ID: LEX-P0018-V0010
% #FIELD_VALUE: 工艺规程编码
% #TODO #HANDWRITTEN: Z-MF-P?01-020; handwriting unclear
\fieldvalue{\handwritten{Z-MF-P?01-020}}
&
版本号：
% #VALUE_ID: LEX-P0018-V0011
% #FIELD_VALUE: 工艺规程版本号
% #TODO #HANDWRITTEN: 7?2?13; handwriting unclear
\fieldvalue{\handwritten{7?2?13}}
& \\
\cline{3-5}
\multicolumn{2}{|c|}{} &
制剂人签字：
% #VALUE_ID: LEX-P0018-V0012
% #FIELD_VALUE: 制剂人签字
% #TODO #HANDWRITTEN: 手写签名; 字迹难以辨认
\fieldvalue{\handwritten{手写签名}}
&
日期：
% #VALUE_ID: LEX-P0018-V0013
% #FIELD_VALUE: 制剂人签字日期
% #HANDWRITTEN: 2025.06.05
\fieldvalue{\handwritten{2025.06.05}}
& \\
\hline
\multicolumn{2}{|c|}{生产审核} &
审核意见：
% #VALUE_ID: LEX-P0014-V0014
% #FIELD_VALUE: 审核意见签字
% #TODO #HANDWRITTEN: 手写签名; 字迹难以辨认
\fieldvalue{\handwritten{手写签名}}
&
\multicolumn{2}{c|}{} \\
\cline{3-5}
\multicolumn{2}{|c|}{} &
生产工艺部负责人签字：
% #VALUE_ID: LEX-P0018-V0015
% #FIELD_VALUE: 生产工艺部负责人签字
% #TODO #HANDWRITTEN: 手写签名; 字迹难以辨认
\fieldvalue{\handwritten{手写签名}}
&
日期：
% #VALUE_ID: LEX-P0018-V0016
% #FIELD_VALUE: 生产工艺部负责人签字日期
% #HANDWRITTEN: 2025.06.05
\fieldvalue{\handwritten{2025.06.05}}
& \\
\hline
\multicolumn{2}{|c|}{QA审核} &
\multicolumn{3}{l|}{以上基本信息审核无误：是
% #VALUE_ID: LEX-P0018-V0017
% #FIELD_VALUE: 基本信息审核无误
\fieldvalue{已选中}
\quad 否} \\
\cline{3-5}
\multicolumn{2}{|c|}{} &
\multicolumn{3}{l|}{公用介质是否满足需求：是
% #VALUE_ID: LEX-P0018-V0018
% #FIELD_VALUE: 公用介质是否满足需求
\fieldvalue{已选中}
\quad 否} \\
\cline{3-5}
\multicolumn{2}{|c|}{} &
\multicolumn{3}{l|}{环境是否满足生产需求：是
% #VALUE_ID: LEX-P0018-V0019
% #FIELD_VALUE: 环境是否满足生产需求
\fieldvalue{已选中}
\quad 否} \\
\cline{3-5}
\multicolumn{2}{|c|}{} &
\multicolumn{3}{l|}{是否有影响生产的偏差：是 \quad 否
% #VALUE_ID: LEX-P0018-V0020
% #FIELD_VALUE: 是否有影响生产的偏差
\fieldvalue{已选中}} \\
\cline{3-5}
\multicolumn{2}{|c|}{} &
\multicolumn{3}{l|}{是否有影响生产的变更：是 \quad 否
% #VALUE_ID: LEX-P0018-V0021
% #FIELD_VALUE: 是否有影响生产的变更
\fieldvalue{已选中}} \\
\cline{3-5}
\multicolumn{2}{|c|}{} &
质量保证部负责人签字/日期：
% #VALUE_ID: LEX-P0018-V0022
% #FIELD_VALUE: 质量保证部负责人签字
% #TODO #HANDWRITTEN: 手写签名; 字迹难以辨认
\fieldvalue{\handwritten{手写签名}}
\quad
% #VALUE_ID: LEX-P0018-V0023
% #FIELD_VALUE: 质量保证部负责人签字日期
% #HANDWRITTEN: 2025.06.05
\fieldvalue{\handwritten{2025.06.05}}
& \\
\hline
\multicolumn{2}{|c|}{批准} &
% #VALUE_ID: LEX-P0018-V0024
% #FIELD_VALUE: 批准意见
\fieldvalue{同意} \\
\cline{3-5}
\multicolumn{2}{|c|}{} &
不同意，原因：\rule{5cm}{0.4pt}
& \multicolumn{2}{c|}{} \\
\cline{3-5}
\multicolumn{2}{|c|}{} &
生产负责人签字/日期：
% #VALUE_ID: LEX-P0018-V0025
% #FIELD_VALUE: 生产负责人签字
% #TODO #HANDWRITTEN: 手写签名; 字迹难以辨认
\fieldvalue{\handwritten{手写签名}}
\quad
% #VALUE_ID: LEX-P0018-V0026
% #FIELD_VALUE: 生产负责人签字日期
% #HANDWRITTEN: 2025.06.05
\fieldvalue{\handwritten{2025.06.05}}
& \\
\hline
\multicolumn{2}{|c|}{指令接收} &
\multicolumn{3}{l|}{指令接收人签字/日期：
% #VALUE_ID: LEX-P0018-V0027
% #FIELD_VALUE: 指令接收人签字
% #TODO #HANDWRITTEN: 林蓉; handwriting unclear
\fieldvalue{\handwritten{林蓉}}
\quad
% #VALUE_ID: LEX-P0018-V0028
% #FIELD_VALUE: 指令接收人签字日期
% #HANDWRITTEN: 2025.06.05
\fieldvalue{\handwritten{2025.06.05}}} \\
\hline
\end{tabularx}

\begin{center}
第 1 页 共 1 页
\end{center}
% LEXOID_PAGE_COMPLETED: 18/59

\newpage

\begin{center}
\textbf{P2代细胞出库申请单、放行证粘贴}
\end{center}

\noindent
品名：
% #VALUE_ID: LEX-P0019-V0001
% #FIELD_VALUE: 品名
\fieldvalue{艾米迦托美注射液}
\hfill
规格：12ml：$6.0\times 10^{7}$个细胞

\noindent
生产批号：
% #VALUE_ID: LEX-P0019-V0002
% #FIELD_VALUE: 生产批号
% #TODO #HANDWRITTEN: A332?L?600; handwriting is partially illegible
\fieldvalue{\handwritten{A332?L?600}}

\vspace{1em}

\begin{center}
\textbf{细胞出库申请单}
\end{center}

\noindent
文件编码：REC-YZ-SMP-MS-03-002-p

\vspace{0.5em}

\renewcommand{\arraystretch}{1.45}
\small
\begin{tabularx}{\textwidth}{|p{3.0cm}|X|p{2.5cm}|X|}
\hline
\multicolumn{4}{|c|}{\textbf{细胞出库申请单}}\\
\hline
出库产品名称
&
% #VALUE_ID: LEX-P0019-V0003
% #FIELD_VALUE: 出库产品名称
\fieldvalue{\ensuremath{\checkmark}\ 艾米迦托美注射液 P2代细胞库}
&
&
\\[-0.2em]
\cline{2-4}
&
$\square$ 艾米迦托美注射液
&
&
\\[-0.2em]
\cline{2-4}
&
$\square$ 其它：\hrulefill
&
&
\\
\hline
出库产品批号
&
% #VALUE_ID: LEX-P0019-V0004
% #FIELD_VALUE: 出库产品批号
% #TODO #HANDWRITTEN: A332?L?600; handwriting is partially illegible
\fieldvalue{\handwritten{A332?L?600}}
&
规格
&
12ml：
% #VALUE_ID: LEX-P0019-V0005
% #FIELD_VALUE: 规格
% #HANDWRITTEN: 5
\fieldvalue{\handwritten{5}}
$\times 10^{7}$个细胞
\\
\hline
出库产品数量
&
% #VALUE_ID: LEX-P0019-V0006
% #FIELD_VALUE: 出库产品数量
% #TODO #HANDWRITTEN: 1; handwritten quantity is faint
\fieldvalue{\handwritten{1}}
&
&
\\
\hline
细胞用途
&
% #VALUE_ID: LEX-P0019-V0007
% #FIELD_VALUE: 细胞用途
\fieldvalue{\ensuremath{\checkmark}\ 艾米迦托美注射液生产}
&
&
\\[-0.2em]
\cline{2-4}
&
$\square$ 入部库存充无干细胞注射液（已接说明书处理）生产
&
&
\\[-0.2em]
\cline{2-4}
&
$\square$ 其它
&
&
\\
\hline
申请人
&
% #VALUE_ID: LEX-P0019-V0008
% #FIELD_VALUE: 申请人
% #TODO #HANDWRITTEN: 林蓉; signature handwriting is uncertain
\fieldvalue{\handwritten{林蓉}}
&
申请日期
&
% #VALUE_ID: LEX-P0019-V0009
% #FIELD_VALUE: 申请日期
% #HANDWRITTEN: 2025-06-08
\fieldvalue{\handwritten{2025-06-08}}
\\
\hline
部门经理
&
% #VALUE_ID: LEX-P0019-V0010
% #FIELD_VALUE: 部门经理
% #TODO #HANDWRITTEN: 不清; signature is illegible
\fieldvalue{\handwritten{不清}}
&
批准日期
&
% #VALUE_ID: LEX-P0019-V0011
% #FIELD_VALUE: 批准日期
% #HANDWRITTEN: 2025-06-09
\fieldvalue{\handwritten{2025-06-09}}
\\
\hline
部门负责人
&
% #VALUE_ID: LEX-P0019-V0012
% #FIELD_VALUE: 部门负责人
% #TODO #HANDWRITTEN: 不清; signature is illegible
\fieldvalue{\handwritten{不清}}
&
批准日期
&
% #VALUE_ID: LEX-P0019-V0013
% #FIELD_VALUE: 批准日期
% #TODO #HANDWRITTEN: 2025-06-09; final digit is unclear
\fieldvalue{\handwritten{2025-06-09}}
\\
\hline
出库人
&
% #VALUE_ID: LEX-P0019-V0014
% #FIELD_VALUE: 出库人
% #TODO #HANDWRITTEN: 不清; handwriting is illegible
\fieldvalue{\handwritten{不清}}
&
出库日期
&
% #VALUE_ID: LEX-P0019-V0015
% #FIELD_VALUE: 出库日期
% #HANDWRITTEN: 2025-06-09
\fieldvalue{\handwritten{2025-06-09}}
\\
\hline
备注
&
% #VALUE_ID: LEX-P0019-V0016
% #FIELD_VALUE: 备注
% #TODO #HANDWRITTEN: /; handwritten mark is unclear
\fieldvalue{\handwritten{/}}
&
\multicolumn{2}{r|}{出库部门保存联}
\\
\hline
\end{tabularx}

\vspace{0.5em}

\noindent
\begin{minipage}[t][8cm][c]{\textwidth}
\centering
P2代细胞放行证粘贴
\end{minipage}

\begin{center}
第 4 页 共 222 页
\end{center}
% LEXOID_PAGE_COMPLETED: 19/59

\newpage

\begin{center}
\small
铂纳生物医药科技（北京）有限公司 \hfill 记录编号：YZ-BPR-MF-PD-01-020\\
\textit{PLATINUM FINE EXCELLENCY BIOTECH}\\[2mm]
\textbf{艾米迈托赛注射液（Z2）生产记录}\\[8mm]
\Large\textbf{艾米迈托赛注射液（Z2）}\\[4mm]
\Large\textbf{批生产记录}
\end{center}

\vspace{12mm}

\noindent
\begin{tabular}{@{}l@{\quad}l@{}}
工序名称：&
% #VALUE_ID: LEX-P0020-V0001
% #FIELD_VALUE: 工序名称
\fieldvalue{艾米迈托赛注射液 P2 代细胞库复苏}\\[7mm]
规格：&
% #VALUE_ID: LEX-P0020-V0002
% #FIELD_VALUE: 规格
\fieldvalue{12ml：$6.0\times10^{7}$个细胞}\\[7mm]
生产批号：&
% #VALUE_ID: LEX-P0020-V0003
% #FIELD_VALUE: 生产批号
% #TODO #HANDWRITTEN: AB322)0050601; handwritten batch number is partially unclear
\fieldvalue{\handwritten{AB322)0050601}}\\[7mm]
批量：&
% #VALUE_ID: LEX-P0020-V0004
% #FIELD_VALUE: 批量
% #HANDWRITTEN: 200--480
\fieldvalue{\handwritten{200--480}}\quad 袋
\end{tabular}

\vspace{11mm}

\noindent
生产日期：\quad
% #VALUE_ID: LEX-P0020-V0005
% #FIELD_VALUE: 生产日期（起始日期）
% #HANDWRITTEN: 2024年6月8日
\fieldvalue{\handwritten{2024年6月8日}}
至
% #VALUE_ID: LEX-P0020-V0006
% #FIELD_VALUE: 生产日期（终止日期）
% #TODO #HANDWRITTEN: 2025年6月?日; final handwritten day is unclear
\fieldvalue{\handwritten{2025年6月?日}}

\vfill

\begin{center}
\small 第 5 页 共 222 页
\end{center}
% LEXOID_PAGE_COMPLETED: 20/59

\newpage

\begin{center}
{\small 铂维生物制药科技（北京）有限公司}\\
{\Large 艾米迈托赛注射液（LZ）批生产记录}\\
记录编号：YZ-PBR-MF-PD-01-020-\ldots
\end{center}

\begin{center}
\begin{tabularx}{\textwidth}{|>{\centering\arraybackslash}X|>{\centering\arraybackslash}X|>{\centering\arraybackslash}X|>{\centering\arraybackslash}X|}
\hline
规格 & 生产批号 & 批量\\
\hline
艾米迈托赛注射液 &
12ml：$6.0\times 10^{9}$个细胞 &
% #VALUE_ID: LEX-P0021-V0001
% #FIELD_VALUE: 生产批号
% #HANDWRITTEN: （1422?2906?）
\fieldvalue{\handwritten{（1422?2906?）}} &
% #VALUE_ID: LEX-P0021-V0002
% #FIELD_VALUE: 批量
% #HANDWRITTEN: 200--480
\fieldvalue{\handwritten{200--480}}\\
\hline
\end{tabularx}
\end{center}

\section*{1.\ 生产前确认：}

\begin{tabularx}{\textwidth}{|>{\raggedright\arraybackslash}X|>{\raggedright\arraybackslash}X|}
\hline
\multicolumn{1}{|c|}{项目} & \multicolumn{1}{c|}{操作记录}\\
\hline
操作人员按《人员进出洁净生产区管理规程》要求进入洁净区生产车间。 &
操作间名称及编码：
% #VALUE_ID: LEX-P0021-V0003
% #FIELD_VALUE: 操作间名称及编码
% #TODO #HANDWRITTEN: 548?6; handwritten code is unclear
\fieldvalue{\handwritten{548?6}}\\
\hline
确认生产区操作间在清场有效期内（C级日清洁效期3天）； &
是否在清洁效期内：%
% #VALUE_ID: LEX-P0021-V0004
% #FIELD_VALUE: 是否在清洁效期内
\fieldvalue{\ensuremath{\boxtimes}\,是}\quad $\square$ 否

清洁效期：
% #VALUE_ID: LEX-P0021-V0005
% #FIELD_VALUE: 清洁效期
% #HANDWRITTEN: 2025.02.10
\fieldvalue{\handwritten{2025.02.10}}\\
\hline
确认操作间温度18--26℃、湿度30\%--75\%、相对压差$\geq 0$Pa（操作间对C级走廊）； &
温度：
% #VALUE_ID: LEX-P0021-V0006
% #FIELD_VALUE: 温度
% #HANDWRITTEN: 17? 
\fieldvalue{\handwritten{17?}}℃\quad
湿度：
% #VALUE_ID: LEX-P0021-V0007
% #FIELD_VALUE: 湿度
% #HANDWRITTEN: 46.7
\fieldvalue{\handwritten{46.7}}\%

压差：
% #VALUE_ID: LEX-P0021-V0008
% #FIELD_VALUE: 压差
% #HANDWRITTEN: 10
\fieldvalue{\handwritten{10}}Pa

是否符合要求：
% #VALUE_ID: LEX-P0021-V0009
% #FIELD_VALUE: 是否符合要求
\fieldvalue{\ensuremath{\boxtimes}\,是}\quad $\square$ 否\\
\hline
确认操作间内无上批产品遗留物、文件及与本产品生产无关物料； &
% #VALUE_ID: LEX-P0021-V0010
% #FIELD_VALUE: 操作间内无遗留物及无关物料
\fieldvalue{\ensuremath{\boxtimes}\,是}\quad $\square$ 否\\
\hline
检查发放的批记录及其它相关配套文件均与生产指令相对应，且均受控； &
% #VALUE_ID: LEX-P0021-V0011
% #FIELD_VALUE: 文件与生产指令相对应且受控
\fieldvalue{\ensuremath{\boxtimes}\,是}\quad $\square$ 否\\
\hline
确认设备好，且在验证有效期内，计量器具有“合格证”且在有效期内，详见《主要设备/仪表确认表》； &
% #VALUE_ID: LEX-P0021-V0012
% #FIELD_VALUE: 设备及计量器具符合要求
\fieldvalue{\ensuremath{\boxtimes}\,是}\quad $\square$ 否\\
\hline
将操作间状态标识更换为“生产加工状态卡”，写明品名、规格、批号等，并确认与本批生产指令一致； &
% #VALUE_ID: LEX-P0021-V0013
% #FIELD_VALUE: 操作间状态标识与生产指令一致
\fieldvalue{\ensuremath{\boxtimes}\,是}\quad $\square$ 否\\
\hline
核对本次生产用物料耗材、溶液和试剂的品名、规格、批号、数量，检查物料外包装的完整性以及是否在有效期/复验期内，并记录相关信息，详见《物料确认表》和《溶液确认表》； &
% #VALUE_ID: LEX-P0021-V0014
% #FIELD_VALUE: 物料、耗材、溶液和试剂符合要求
\fieldvalue{\ensuremath{\boxtimes}\,是}\quad $\square$ 否\\
\hline
备注： &
% #VALUE_ID: LEX-P0021-V0015
% #FIELD_VALUE: 备注
% #HANDWRITTEN: -
\fieldvalue{\handwritten{-}}\\
\hline
\end{tabularx}

\vspace{0.5em}

确认人/日期：
% #VALUE_ID: LEX-P0021-V0016
% #FIELD_VALUE: 确认人
% #TODO #HANDWRITTEN: Z?lb; handwritten name is unclear
\fieldvalue{\handwritten{Z?lb}}
\quad
% #VALUE_ID: LEX-P0021-V0017
% #FIELD_VALUE: 确认日期
% #HANDWRITTEN: 2025.06.09
\fieldvalue{\handwritten{2025.06.09}}
\hfill
复核人/日期：
% #VALUE_ID: LEX-P0021-V0018
% #FIELD_VALUE: 复核人
% #TODO #HANDWRITTEN: 赵?; handwritten name is unclear
\fieldvalue{\handwritten{赵?}}
\quad
% #VALUE_ID: LEX-P0021-V0019
% #FIELD_VALUE: 复核日期
% #HANDWRITTEN: 2025.06.09
\fieldvalue{\handwritten{2025.06.09}}

\subsection*{1.1\ 主要设备/仪表确认表：}

\begin{center}
\begin{tabularx}{\textwidth}{|c|X|X|X|X|X|}
\hline
序号 & 设备名称 & 设备型号 & 设备编号 & 检查项 & 确认是否符合要求\\
\hline
\end{tabularx}
\end{center}

% TODO: The equipment table continues beyond the visible portion of this page.

\begin{center}
{\small 第 6 页 共 22 页}
\end{center}
% LEXOID_PAGE_COMPLETED: 21/59

\newpage

\section*{艾米迈托基注射液（Z2）批生产记录}

\noindent
记录编号：YZ-BPR-MF-PD-01-020-a\hfill 批号：13

\begin{tabularx}{\textwidth}{|c|X|X|X|c|c|}
\hline
序号 & 品名 & 规格 & 生产批号 & 批量 & \\
\hline
1 &
电热恒温循环水槽 &
DK-8AB &
% #VALUE_ID: LEX-P0022-V0001
% #FIELD_VALUE: 生产批号
% #TODO #HANDWRITTEN: 112?02; handwriting unclear
\fieldvalue{\handwritten{112?02}} &
&
% #VALUE_ID: LEX-P0022-V0002
% #FIELD_VALUE: 设备确认
\fieldvalue{☑ 是}\quad □ 否
\\
\hline
2 &
洁净工作台 &
SW-CJ-2FD &
% #VALUE_ID: LEX-P0022-V0003
% #FIELD_VALUE: 生产批号
% #TODO #HANDWRITTEN: 114?02; handwriting unclear
\fieldvalue{\handwritten{114?02}} &
&
% #VALUE_ID: LEX-P0022-V0004
% #FIELD_VALUE: 设备确认
\fieldvalue{☑ 是}\quad □ 否
\\
\hline
3 &
3D FloTrix\textsuperscript{\textregistered} vivaPREP 细胞收获系统（简称 VP-管） &
3D FloTrix\textsuperscript{\textregistered} vivaPREP &
% #VALUE_ID: LEX-P0022-V0005
% #FIELD_VALUE: 生产批号
% #TODO #HANDWRITTEN: 112?70; handwriting unclear
\fieldvalue{\handwritten{112?70}} &
&
% #VALUE_ID: LEX-P0022-V0006
% #FIELD_VALUE: 设备确认
\fieldvalue{☑ 是}\quad □ 否
\\
\hline
4 &
接管机 &
泰尔茂 TSCD II &
% #VALUE_ID: LEX-P0022-V0007
% #FIELD_VALUE: 生产批号
% #HANDWRITTEN: 1120068
\fieldvalue{\handwritten{1120068}} &
&
% #VALUE_ID: LEX-P0022-V0008
% #FIELD_VALUE: 设备确认
\fieldvalue{☑ 是}\quad □ 否
\\
\hline
5 &
封管机 &
Biosealer CR 6AA &
% #VALUE_ID: LEX-P0022-V0009
% #FIELD_VALUE: 生产批号
% #HANDWRITTEN: 1120901
\fieldvalue{\handwritten{1120901}} &
&
% #VALUE_ID: LEX-P0022-V0010
% #FIELD_VALUE: 设备确认
\fieldvalue{☑ 是}\quad □ 否
\\
\hline
6 & & & & & □ 是\quad □ 否\\
\hline
7 & & & & & □ 是\quad □ 否\\
\hline
\end{tabularx}

\noindent
\parbox{\textwidth}{%
\textbf{设备确认要求：}
1.\ 设备应有“完好”标识；
2.\ 设备在验证有效期内；
3.\ 计量器具在有效期内。%
}

% #VALUE_ID: LEX-P0022-V0011
% #TODO #HANDWRITTEN: 2025-06-09（前部文字不清）; diagonal handwriting across rows 6--7
\noindent\handwritten{2025-06-09（前部文字不清）}

% #VALUE_ID: LEX-P0022-V0012
% #TODO #HANDWRITTEN: 不清; vertical handwriting beside the right-side checkboxes
\noindent\handwritten{不清}

\noindent
备注：
% #VALUE_ID: LEX-P0022-V0013
% #TODO #HANDWRITTEN: 短划线; handwritten mark in the remarks field
\handwritten{短划线}

\medskip

\noindent
确认人/日期：
% #VALUE_ID: LEX-P0022-V0014
% #TODO #HANDWRITTEN: 弘峰; handwritten name is unclear
% #FIELD_VALUE: 确认人
% #HANDWRITTEN: 弘峰
\fieldvalue{\handwritten{弘峰}}
\quad
% #VALUE_ID: LEX-P0022-V0015
% #FIELD_VALUE: 确认日期
% #HANDWRITTEN: 2025.06.09
\fieldvalue{\handwritten{2025.06.09}}
\qquad
复核人/日期：
% #VALUE_ID: LEX-P0022-V0016
% #TODO #HANDWRITTEN: 赵永刚; handwritten name is somewhat unclear
% #FIELD_VALUE: 复核人
% #HANDWRITTEN: 赵永刚
\fieldvalue{\handwritten{赵永刚}}
\quad
% #VALUE_ID: LEX-P0022-V0017
% #FIELD_VALUE: 复核日期
% #HANDWRITTEN: 2025.06.09
\fieldvalue{\handwritten{2025.06.09}}

\subsection*{1.2\ 物料确认表}

\begin{tabularx}{\textwidth}{|c|X|c|c|X|c|c|}
\hline
序号 & 名称 & 品牌 & 规格 & 批号 & 数量 & 是否符合要求 \\
\hline
1 &
3D FloTrix\textsuperscript{\textregistered} vivaPREP 细胞收获系统\textsuperscript{\textregistered} -\textsuperscript{简}容包（简称 VP-管包） &
华泰生物 &
1 个/包 &
% #VALUE_ID: LEX-P0022-V0018
% #FIELD_VALUE: 批号
% #TODO #HANDWRITTEN: 105?30/2025?06; handwriting unclear
\fieldvalue{\handwritten{105?30/2025?06}} &
% #VALUE_ID: LEX-P0022-V0019
% #FIELD_VALUE: 数量
% #HANDWRITTEN: 1
\fieldvalue{\handwritten{1}} &
% #VALUE_ID: LEX-P0022-V0020
% #FIELD_VALUE: 是否符合要求
\fieldvalue{☑ 是}\quad □ 否
\\
\hline
2 &
一次性使用无菌注射器 &
康德莱 &
1 个/包 &
% #VALUE_ID: LEX-P0022-V0021
% #FIELD_VALUE: 批号
% #TODO #HANDWRITTEN: H?0628/2025?0129; handwriting unclear
\fieldvalue{\handwritten{H?0628/2025?0129}} &
% #VALUE_ID: LEX-P0022-V0022
% #FIELD_VALUE: 数量
% #HANDWRITTEN: 1
\fieldvalue{\handwritten{1}} &
% #VALUE_ID: LEX-P0022-V0023
% #FIELD_VALUE: 是否符合要求
\fieldvalue{☑ 是}\quad □ 否
\\
\hline
3 & & & & & & □ 是\quad □ 否\\
\hline
4 & & & & & & □ 是\quad □ 否\\
\hline
\end{tabularx}

% #VALUE_ID: LEX-P0022-V0024
% #TODO #HANDWRITTEN: 2025-06-09（前部姓名不清）; diagonal handwriting across the lower blank rows
\noindent\handwritten{2025-06-09（前部姓名不清）}

\noindent
备注：

\medskip

\noindent
确认人/日期：
% #VALUE_ID: LEX-P0022-V0025
% #TODO #HANDWRITTEN: 弘峰; handwritten name is unclear
% #FIELD_VALUE: 确认人
% #HANDWRITTEN: 弘峰
\fieldvalue{\handwritten{弘峰}}
\quad
% #VALUE_ID: LEX-P0022-V0026
% #FIELD_VALUE: 确认日期
% #HANDWRITTEN: 2025.06.09
\fieldvalue{\handwritten{2025.06.09}}
\qquad
复核人/日期：
% #VALUE_ID: LEX-P0022-V0027
% #TODO #HANDWRITTEN: 赵永刚; handwritten name is unclear
% #FIELD_VALUE: 复核人
% #HANDWRITTEN: 赵永刚
\fieldvalue{\handwritten{赵永刚}}
\quad
% #VALUE_ID: LEX-P0022-V0028
% #FIELD_VALUE: 复核日期
% #HANDWRITTEN: 2025.06.09
\fieldvalue{\handwritten{2025.06.09}}

\subsection*{1.3\ 溶液确认表}

\noindent
% TODO: The solution-confirmation table begins at the bottom of this page and continues on the following page; no table rows are visible here.

\begin{center}
第 7 页 共 22 页
\end{center}
% LEXOID_PAGE_COMPLETED: 22/59

\newpage

\section{艾米迈托赛注射液（Z2）批生产记录}

\noindent
记录编号：YZ-BPR-MF-PD-01-020-a\hfill 版本号：1.3

\begin{tabularx}{\textwidth}{|c|X|X|X|X|}
\hline
品名 & 规格 & 生产批号 & 批量 \\
\hline
艾米迈托赛注射液 & 12ml：$6.0\times10^{7}$个细胞 &
% #VALUE_ID: LEX-P0023-V0001
% #FIELD_VALUE: 生产批号
% #TODO #HANDWRITTEN: 430?202506?01; 字迹不清
\fieldvalue{\handwritten{430?202506?01}} &
% #VALUE_ID: LEX-P0023-V0002
% #FIELD_VALUE: 批量
% #TODO #HANDWRITTEN: 20袋-?; 字迹不清
\fieldvalue{\handwritten{20袋-?}} \\
\hline
\end{tabularx}

\begin{tabularx}{\textwidth}{|c|X|X|X|X|X|}
\hline
序号 & 名称 & 规格 & 批号 & 数量 & 有效期 \\
\hline
1 &
洗涤液（2\%人白蛋白溶液） &
500ml/袋 &
% #VALUE_ID: LEX-P0023-V0003
% #FIELD_VALUE: 第1项批号
% #TODO #HANDWRITTEN: 202506?01; 字迹不清
\fieldvalue{\handwritten{202506?01}} &
% #VALUE_ID: LEX-P0023-V0004
% #FIELD_VALUE: 第1项数量
\fieldvalue{\handwritten{1}} &
% #VALUE_ID: LEX-P0023-V0005
% #FIELD_VALUE: 第1项有效期
% #HANDWRITTEN: 2025.07.06
\fieldvalue{\handwritten{2025.07.06}} \\
\hline
\multicolumn{1}{|c|}{} &
\multicolumn{1}{X|}{} &
\multicolumn{1}{X|}{} &
\multicolumn{1}{X|}{} &
\multicolumn{1}{X|}{} &
\multicolumn{1}{X|}{是否符合要求：%
% #VALUE_ID: LEX-P0023-V0006
% #FIELD_VALUE: 第1项是否符合要求
% #HANDWRITTEN: 是（勾选）
\fieldvalue{\handwritten{是（勾选）}} \quad 否} \\
\hline
2 &
重悬液（完全培养基） &
500ml/袋 &
% #VALUE_ID: LEX-P0023-V0007
% #FIELD_VALUE: 第2项批号
% #TODO #HANDWRITTEN: 202506?01; 字迹不清
\fieldvalue{\handwritten{202506?01}} &
% #VALUE_ID: LEX-P0023-V0008
% #FIELD_VALUE: 第2项数量
\fieldvalue{\handwritten{1}} &
% #VALUE_ID: LEX-P0023-V0009
% #FIELD_VALUE: 第2项有效期
% #HANDWRITTEN: 2025.06.07
\fieldvalue{\handwritten{2025.06.07}} \\
\hline
\multicolumn{1}{|c|}{} &
\multicolumn{1}{X|}{} &
\multicolumn{1}{X|}{} &
\multicolumn{1}{X|}{} &
\multicolumn{1}{X|}{} &
\multicolumn{1}{X|}{是否符合要求：%
% #VALUE_ID: LEX-P0023-V0010
% #FIELD_VALUE: 第2项是否符合要求
% #HANDWRITTEN: 是（勾选）
\fieldvalue{\handwritten{是（勾选）}} \quad 否} \\
\hline
3 & & & & &
% #VALUE_ID: LEX-P0023-V0011
% #FIELD_VALUE: 第3项有效期
% #TODO #HANDWRITTEN: 2025.??.??; 字迹不清
\fieldvalue{\handwritten{2025.??.??}} \\
\hline
\end{tabularx}

备注：%
% #VALUE_ID: LEX-P0023-V0012
% #FIELD_VALUE: 备注
% #HANDWRITTEN: /
\fieldvalue{\handwritten{/}}

\noindent
确认入厂日期：%
% #VALUE_ID: LEX-P0023-V0013
% #FIELD_VALUE: 确认人
% #TODO #HANDWRITTEN: 张?坤; 姓名字迹不清
\fieldvalue{\handwritten{张?坤}}\quad
% #VALUE_ID: LEX-P0023-V0014
% #FIELD_VALUE: 确认日期
% #HANDWRITTEN: 2025.06.07
\fieldvalue{\handwritten{2025.06.07}}
\hfill
复核人/日期：%
% #VALUE_ID: LEX-P0023-V0015
% #FIELD_VALUE: 复核人
% #TODO #HANDWRITTEN: 张?坤; 姓名字迹不清
\fieldvalue{\handwritten{张?坤}}\quad
% #VALUE_ID: LEX-P0023-V0016
% #FIELD_VALUE: 复核日期
% #HANDWRITTEN: 2025.06.09
\fieldvalue{\handwritten{2025.06.09}}

\subsection{P2 低温质粒冷冻复苏}

\subsubsection{复苏前准备}

复苏前准备（开始时间：%
% #VALUE_ID: LEX-P0023-V0017
% #FIELD_VALUE: 开始时间
% #HANDWRITTEN: 01：45
\fieldvalue{\handwritten{01：45}}）（以电热恒温水槽加热时间为准）

\begin{tabularx}{\textwidth}{|X|X|}
\hline
项目 & 操作记录 \\
\hline
1）在电热恒温循环水槽内加入注射用水至刻度线以上，并用量尺测量水槽内水位高度（$\geq 10\mathrm{cm}$），等待加热至水槽温度显示设定温度（$39^\circ\mathrm{C}$）。 &
水位高度：%
% #VALUE_ID: LEX-P0023-V0018
% #FIELD_VALUE: 水位高度
% #HANDWRITTEN: 10
\fieldvalue{\handwritten{10}} cm

设定温度：%
% #VALUE_ID: LEX-P0023-V0019
% #FIELD_VALUE: 设定温度
% #HANDWRITTEN: 39
\fieldvalue{\handwritten{39}} $^\circ$C \\
\hline
\end{tabularx}

备注：%
% #VALUE_ID: LEX-P0023-V0020
% #FIELD_VALUE: 复苏前准备备注
% #HANDWRITTEN: /
\fieldvalue{\handwritten{/}}

\noindent
操作人/日期：%
% #VALUE_ID: LEX-P0023-V0021
% #FIELD_VALUE: 操作人
% #TODO #HANDWRITTEN: 张?坤; 姓名字迹不清
\fieldvalue{\handwritten{张?坤}}\quad
% #VALUE_ID: LEX-P0023-V0022
% #FIELD_VALUE: 操作日期
% #HANDWRITTEN: 2025.06.07
\fieldvalue{\handwritten{2025.06.07}}
\hfill
复核人/日期：%
% #VALUE_ID: LEX-P0023-V0023
% #FIELD_VALUE: 复核人
% #TODO #HANDWRITTEN: 张?坤; 姓名字迹不清
\fieldvalue{\handwritten{张?坤}}\quad
% #VALUE_ID: LEX-P0023-V0024
% #FIELD_VALUE: 复核日期
% #HANDWRITTEN: 2025.06.09
\fieldvalue{\handwritten{2025.06.09}}

\begin{tabularx}{\textwidth}{|X|X|}
\hline
项目 & 操作记录 \\
\hline
2）按生产指令要求填写《细胞出库申请单》，从细胞库内领取艾米迈托赛P2代细胞1袋，并按照《物料出入洁净区管理规程》将细胞经液氮转运至细胞处理间，并将细胞信息标签粘贴于右侧； &
细胞经液氮转运至操作间时间：%
% #VALUE_ID: LEX-P0023-V0025
% #FIELD_VALUE: 细胞经液氮转运至操作间时间
% #HANDWRITTEN: 01：25
\fieldvalue{\handwritten{01：25}}

细胞信息标签粘贴处：

% TODO: 此处为页面中可见的细胞信息标签图片，未提供图像文件名。

% #VALUE_ID: LEX-P0023-V0026
% #FIELD_VALUE: 标签处签名
% #TODO #HANDWRITTEN: 张?坤; 图片内签名字迹不清
\fieldvalue{\handwritten{张?坤}}\quad
% #VALUE_ID: LEX-P0023-V0027
% #FIELD_VALUE: 标签处日期
% #HANDWRITTEN: 2025.06.09
\fieldvalue{\handwritten{2025.06.09}} \\
\hline
\end{tabularx}

备注：%
% #VALUE_ID: LEX-P0023-V0028
% #FIELD_VALUE: 操作记录备注
% #HANDWRITTEN: /
\fieldvalue{\handwritten{/}}

\noindent
操作人/日期：%
% #VALUE_ID: LEX-P0023-V0029
% #FIELD_VALUE: 操作人
% #TODO #HANDWRITTEN: 张?坤; 姓名字迹不清
\fieldvalue{\handwritten{张?坤}}\quad
% #VALUE_ID: LEX-P0023-V0030
% #FIELD_VALUE: 操作日期
% #HANDWRITTEN: 2025.06.09
\fieldvalue{\handwritten{2025.06.09}}
\hfill
复核人/日期：%
% #VALUE_ID: LEX-P0023-V0031
% #FIELD_VALUE: 复核人
% #TODO #HANDWRITTEN: 张?坤; 姓名字迹不清
\fieldvalue{\handwritten{张?坤}}\quad
% #VALUE_ID: LEX-P0023-V0032
% #FIELD_VALUE: 复核日期
% #HANDWRITTEN: 2025.06.09
\fieldvalue{\handwritten{2025.06.09}}
% LEXOID_PAGE_COMPLETED: 23/59

\newpage

\begin{center}
\textbf{艾米迈托赛针液（Z2）批生产记录}
\end{center}

\begin{tabularx}{\textwidth}{|>{\centering\arraybackslash}X|>{\centering\arraybackslash}X|>{\centering\arraybackslash}X|>{\centering\arraybackslash}X|}
\hline
产品名称 & 规格 & 生产批号 & 批量 \\
\hline
艾米迈托赛针液 & 12ml：$6.0\times10^{7}$个细胞 &
% #VALUE_ID: LEX-P0024-V0001
% #FIELD_VALUE: 生产批号
% #TODO #HANDWRITTEN: A332?045?01; 手写内容部分不清晰
\fieldvalue{\handwritten{A332?045?01}} &
% #VALUE_ID: LEX-P0024-V0002
% #FIELD_VALUE: 批量
% #TODO #HANDWRITTEN: 20.10?80; 手写内容部分不清晰
\fieldvalue{\handwritten{20.10?80}} \\
\hline
\end{tabularx}

\medskip

\textbf{2.2 VP管路包装：}

\begin{tabularx}{\textwidth}{|>{\raggedright\arraybackslash}X|>{\raggedright\arraybackslash}p{0.25\textwidth}|}
\hline
\centering 项目 & \centering 操作记录 \arraybackslash \\
\hline
1）洁净工作台开启紫外照射至少30min，开启风机，运行自净10分钟以上，将耗材包和相关器械摆清洁后传入洁净工作台。 &
紫外照射时间：

% #VALUE_ID: LEX-P0024-V0003
% #FIELD_VALUE: 紫外照射时间
% #HANDWRITTEN: 06:27--07:07
\fieldvalue{\handwritten{06:27--07:07}}

自净时间：

% #VALUE_ID: LEX-P0024-V0004
% #FIELD_VALUE: 自净时间
% #HANDWRITTEN: 07:07--07:17
\fieldvalue{\handwritten{07:07--07:17}}

是否符合要求：

% #VALUE_ID: LEX-P0024-V0005
% #FIELD_VALUE: 是否符合要求
% #HANDWRITTEN: 是（勾选）
\fieldvalue{\handwritten{\checkmark\ 是}}\quad 否
\\
\hline
2）耗材准备：首先检查包装的完整性，然后在洁净工作台内打开VP-管路包，确保管路包上的无菌指示标识为蓝色，打开管路包装，取出耗材，并关闭管路上所有管夹（12个白色和1个蓝色）和3个蓝色夹片；海绵其中一段耗材包内的穿刺针，使用时留有足够长度的PVC管路，以便接下来接管使用），保留待用； &
无菌指示标识颜色：

% #VALUE_ID: LEX-P0024-V0006
% #FIELD_VALUE: 无菌指示标识颜色
% #HANDWRITTEN: 蓝色（勾选）
\fieldvalue{\handwritten{\checkmark\ 蓝色}}\quad 其它\underline{\hspace{1.5cm}}

管夹是否已关闭：

% #VALUE_ID: LEX-P0024-V0007
% #FIELD_VALUE: 管夹是否已关闭
% #HANDWRITTEN: 是（勾选）
\fieldvalue{\handwritten{\checkmark\ 是}}\quad 否

夹片是否已关闭：

% #VALUE_ID: LEX-P0024-V0008
% #FIELD_VALUE: 夹片是否已关闭
% #HANDWRITTEN: 是（勾选）
\fieldvalue{\handwritten{\checkmark\ 是}}\quad 否

是否取下穿刺针：

% #VALUE_ID: LEX-P0024-V0009
% #FIELD_VALUE: 是否取下穿刺针
% #HANDWRITTEN: 是（勾选）
\fieldvalue{\handwritten{\checkmark\ 是}}\quad 否
\\
\hline
3）将耗材包从洁净工作台内拿出，将1袋渗透液（500ml/袋）连接至管路包的2号管路，1袋洗涤液（2\%人血白蛋白，500ml/袋）连接至3号管路上； &
滤液是否连接完毕：

% #VALUE_ID: LEX-P0024-V0010
% #FIELD_VALUE: 滤液是否连接完毕
% #HANDWRITTEN: 是（勾选）
\fieldvalue{\handwritten{\checkmark\ 是}}\quad 否
\\
\hline
\end{tabularx}

\vfill

\begin{center}
第 9 页 共 222 页
\end{center}
% LEXOID_PAGE_COMPLETED: 24/59

\newpage

\section*{艾米迈托赛注射液（Z2）批生产记录}

\begin{center}
\begin{tabularx}{\textwidth}{|>{\centering\arraybackslash}X|>{\centering\arraybackslash}X|>{\centering\arraybackslash}X|>{\centering\arraybackslash}X|}
\hline
品名 & 规格 & 生产批号 & 批量 \\
\hline
艾米迈托赛注射液 & 12ml；$6.0\times10^{7}$个细胞 &
% #VALUE_ID: LEX-P0025-V0001
% #FIELD_VALUE: 生产批号
% #HANDWRITTEN: A3?2) DXS 0600
\fieldvalue{\handwritten{A3?2) DXS 0600}} &
% #VALUE_ID: LEX-P0025-V0002
% #FIELD_VALUE: 批量
% #HANDWRITTEN: 20--?98 接
\fieldvalue{\handwritten{20--?98 接}} \\
\hline
\end{tabularx}
\end{center}

记录编号：YZ-BPR-MF-PD-01-020-483

\subsection*{4）VP设备开机、自检通过后，登录选择“复苏细胞洗涤”工艺并加载程序，输入批次信息及耗材包批号，VP管路包连接安装系统，重复浸泡和洗涤液袋分别悬挂在对应支架上；}

\begin{flushright}
\begin{tabular}{|p{0.22\textwidth}|p{0.18\textwidth}|}
\hline
开机时间： &
% #VALUE_ID: LEX-P0025-V0003
% #FIELD_VALUE: 开机时间
% #HANDWRITTEN: 07:1?
\fieldvalue{\handwritten{07:1?}} \\ \hline
管路包安装完成时间： &
% #VALUE_ID: LEX-P0025-V0004
% #FIELD_VALUE: 管路包安装完成时间
% #HANDWRITTEN: 07:27
\fieldvalue{\handwritten{07:27}} \\ \hline
\end{tabular}
\end{flushright}

\subsection*{5）设定参数：选择“复苏细胞洗涤”工艺，点击编辑好的操作程序，检查参数是否符合如下要求：}

\begin{center}
\small
\begin{tabularx}{0.72\textwidth}{|c|c|c|c|c|c|}
\hline
项目 & 步骤 & 具体参数 & 项目 & 步骤 & 具体参数 \\
\hline
\multirow{4}{*}{1、设备预设值} &
空样品设置量（g） & 20 &
\multirow{4}{*}{4、样品洗涤} &
样品洗涤 & 1* \\
\cline{2-3}\cline{5-6}
&样品输入量（g）&11&洗液次数（次）&1*\\
\cline{2-3}\cline{5-6}
&样品原工艺前安装&0*&洗液使用体积（g）&220\\
\cline{2-3}\cline{5-6}
&样品导管&1*&洗液滴加气（kg）&200\\
\hline
\multirow{8}{*}{2、样品混配} &
剪样使用体积（g）&40&
\multirow{8}{*}{}
洗液滴加时间（min）&5\\
\cline{2-3}\cline{5-6}
&自动稀释体积（g）&1* &洗液相分料（g）&200\\
\cline{2-3}\cline{5-6}
&样品解冻速度（g/s）&0.5&洗液添加加速度（rpm/s）&120\\
\cline{2-3}\cline{5-6}
&降解血浆包相体积（g）&0&洗液混匀加速度（rpm/s）&1000\\
\cline{2-3}\cline{5-6}
&离心精度误差修正（g）&1*&洗液混匀时间（min）&120\\
\cline{2-3}\cline{5-6}
&预热洗液洗涤液（g）&180&洗液混匀加减时间（s）&5\\
\cline{2-3}\cline{5-6}
&洗涤样品设置量（g）&0&洗液分次次数（次）&4\\
\cline{2-3}\cline{5-6}
&样品混配浓度&1*&洗液分次加速度（rpm/s）&3\\
\hline
\multirow{11}{*}{3、样品准备} &
样品混配总初样（g）&0&
\multirow{11}{*}{5、产品输出}&
终产液（g）&50\\
\cline{2-3}\cline{5-6}
&浓缩离心力（xg）&200&手动收集&1*\\
\cline{2-3}\cline{5-6}
&浓缩离心时间（min）&5&重悬混匀加速度（rpm/s）&120\\
\cline{2-3}\cline{5-6}
&浓缩离心偏差&20&重悬混匀加速度（rpm/s）&1000\\
\cline{2-3}\cline{5-6}
&样品离心体积（g）&0&重悬混匀加速度（rpm/s）&120\\
\cline{2-3}\cline{5-6}
&混匀加速度（rpm）&1000&重悬混匀时间（s）&5\\
\cline{2-3}\cline{5-6}
&混匀速度（rpm）&120&重悬混匀次数（次）&2\\
\cline{2-3}\cline{5-6}
&混匀时间（s）&5&重悬混匀加速度（rpm/s）&3\\
\cline{2-3}\cline{5-6}
&混匀次数（次）&3&&\\
\cline{2-3}
&混匀方式（°）&3&&\\
\hline
\multicolumn{6}{|l|}{6、“*”表示是否执行该步骤，1代表设备运行此步骤，0代表设备不运行此步骤。}\\
\hline
\end{tabularx}
\end{center}

\begin{flushright}
\begin{tabular}{|p{0.28\textwidth}|p{0.15\textwidth}|}
\hline
参数设定是否符合要求： &
\begin{tabular}{l}
% #VALUE_ID: LEX-P0025-V0005
% #FIELD_VALUE: 参数设定是否符合要求
% #HANDWRITTEN: ✓ 是
\fieldvalue{\handwritten{✓ 是}}\quad 否
\end{tabular}
\\ \hline
设定完成时间： &
% #VALUE_ID: LEX-P0025-V0006
% #FIELD_VALUE: 设定完成时间
% #HANDWRITTEN: 07:32
\fieldvalue{\handwritten{07:32}} \\ \hline
\end{tabular}
\end{flushright}

\vfill

\begin{center}
第 10 页\quad 共 222 页
\end{center}
% LEXOID_PAGE_COMPLETED: 25/59

\newpage

\begin{center}
\textbf{艾米迈托赛注射液（Z2）批生产记录}
\end{center}

记录编号：YZ-BPR-MF-PD-01-020-

\begin{tabularx}{\textwidth}{|X|X|X|X|}
\hline
品名 & 规格 & 生产批号 & 批量 \\
\hline
艾米迈托赛注射液 & 12ml；$6.0\times10^{7}$个细胞 &
% #VALUE_ID: LEX-P0026-V0001
% #FIELD_VALUE: 生产批号
% #TODO #HANDWRITTEN: A332D?200?; handwriting is unclear
\fieldvalue{\handwritten{A332D?200?}} &
% #VALUE_ID: LEX-P0026-V0002
% #FIELD_VALUE: 批量
% #TODO #HANDWRITTEN: 20-?0-?袋; handwriting is unclear
\fieldvalue{\handwritten{20-?0-?袋}} \\
\hline
\end{tabularx}

\vspace{0.5em}

5）参数设定完成后，根据设备提示开始耗材包自检及洗涤液预填充，当设备调阅预准备细胞样品时，可开始细胞水槽复苏；

实际完成时间：
% #VALUE_ID: LEX-P0026-V0003
% #FIELD_VALUE: 实际完成时间
% #HANDWRITTEN: 07：38
\fieldvalue{\handwritten{07：38}}

\begin{tabularx}{\textwidth}{|X|}
\hline
备注：\\[1.5em]
\hline
\end{tabularx}

操作人/日期：
% #VALUE_ID: LEX-P0026-V0004
% #FIELD_VALUE: 操作人
% #TODO #HANDWRITTEN: 张坤; handwriting is somewhat unclear
\fieldvalue{\handwritten{张坤}}
\quad
% #VALUE_ID: LEX-P0026-V0005
% #FIELD_VALUE: 日期
% #HANDWRITTEN: 2025.06.09
\fieldvalue{\handwritten{2025.06.09}}
\hfill
复核人/日期：
% #VALUE_ID: LEX-P0026-V0006
% #FIELD_VALUE: 复核人
% #TODO #HANDWRITTEN: 赵小?; handwriting is unclear
\fieldvalue{\handwritten{赵小?}}
\quad
% #VALUE_ID: LEX-P0026-V0007
% #FIELD_VALUE: 日期
% #HANDWRITTEN: 2025.06.09
\fieldvalue{\handwritten{2025.06.09}}

\section*{3. P2 代细胞复苏及 VP 处理}

\subsection*{3.1 细胞复苏（开始时间：
% #VALUE_ID: LEX-P0026-V0008
% #FIELD_VALUE: 开始时间
% #HANDWRITTEN: 07：41
\fieldvalue{\handwritten{07：41}}
）（以细胞离开液氮为准）}

\begin{tabularx}{\textwidth}{|p{0.58\textwidth}|X|}
\hline
\centering 项目 & \centering 操作记录 \tabularnewline
\hline
1）复苏准备工作完成，待电热恒温水槽中注射用水加热至目标温度范围（38.5$\pm$1.5℃）后，用止血钳从液氮中夹起冻存细胞，核对项目信息后，将冻存细胞放入电热恒温水槽内，位置以水波没过冻存细胞为宜。 &
电热恒温水槽显示温度：
% #VALUE_ID: LEX-P0026-V0009
% #FIELD_VALUE: 电热恒温水槽显示温度
% #HANDWRITTEN: 22.9
\fieldvalue{\handwritten{22.9}} ℃\\
\hline
2）记录复苏开始时间，并开启计时器，复苏过程中不断轻轻摇晃，观察细胞融化进度，待冻存袋细胞融化、无可见冰晶后，立即将所有冻存细胞从水浴中取出，记录复苏时长（20s$\sim$70s）； &
复苏时间：
% #VALUE_ID: LEX-P0026-V0010
% #FIELD_VALUE: 复苏时间
% #HANDWRITTEN: 07：41--07：41
\fieldvalue{\handwritten{07：41--07：41}}\\
&
复苏时长：
% #VALUE_ID: LEX-P0026-V0011
% #FIELD_VALUE: 复苏时长
% #HANDWRITTEN: 0 min 40 s
\fieldvalue{\handwritten{0 min 40 s}}\\
\hline
3）复苏完成后，检查电热恒温水槽温度； &
电热恒温水槽显示温度：
% #VALUE_ID: LEX-P0026-V0012
% #FIELD_VALUE: 电热恒温水槽显示温度
% #HANDWRITTEN: 38.7
\fieldvalue{\handwritten{38.7}} ℃\\
\hline
\end{tabularx}

操作人/日期：
% #VALUE_ID: LEX-P0026-V0013
% #FIELD_VALUE: 操作人
% #TODO #HANDWRITTEN: 张坤; handwriting is unclear
\fieldvalue{\handwritten{张坤}}
\quad
% #VALUE_ID: LEX-P0026-V0014
% #FIELD_VALUE: 日期
% #HANDWRITTEN: 2025.06.09
\fieldvalue{\handwritten{2025.06.09}}
\hfill
复核人/日期：
% #VALUE_ID: LEX-P0026-V0015
% #FIELD_VALUE: 复核人
% #TODO #HANDWRITTEN: 赵小?; handwriting is unclear
\fieldvalue{\handwritten{赵小?}}
\quad
% #VALUE_ID: LEX-P0026-V0016
% #FIELD_VALUE: 日期
% #HANDWRITTEN: 2025.06.09
\fieldvalue{\handwritten{2025.06.09}}

\subsection*{3.2 VP 处理}

\begin{tabularx}{\textwidth}{|p{0.58\textwidth}|X|}
\hline
1）细胞连接

将复苏的细胞冻存袋用75\%酒精喷洒表面消毒，传入洁净工作台内，打开细胞冻存袋螺口，将穿刺孔与细胞冻存袋螺口进行连接。 &
细胞冻存袋穿刺连接时间：

% #VALUE_ID: LEX-P0026-V0017
% #FIELD_VALUE: 细胞冻存袋穿刺连接时间
% #TODO #HANDWRITTEN: 07：45; handwriting is unclear
\fieldvalue{\handwritten{07：45}}\\
\hline
\end{tabularx}

% TODO：本节内容在后续页面继续。
% LEXOID_PAGE_COMPLETED: 26/59

\newpage

\begin{center}
{\Large \textbf{艾米近无菌注射液（Z2）批生产记录}}
\end{center}

\noindent
记录编号：
% #VALUE_ID: LEX-P0027-V0001
% #FIELD_VALUE: 记录编号
\fieldvalue{YZ-BPR-MF-PD-01-020-a}
\hfill
批号：
% #VALUE_ID: LEX-P0027-V0002
% #FIELD_VALUE: 批号
\fieldvalue{1.3}

\begin{tabularx}{\textwidth}{|>{\centering\arraybackslash}p{0.18\textwidth}|X|>{\centering\arraybackslash}p{0.18\textwidth}|X|}
\hline
品名 & 规格 & 生产批号 & 批量 \\
\hline
艾米近无菌注射液 & 12ml；$6.0\times10^7$个细胞 &
% #VALUE_ID: LEX-P0027-V0003
% #FIELD_VALUE: 生产批号
% TODO #HANDWRITTEN: 2025.06.09/01; handwriting partially unclear
\fieldvalue{\handwritten{2025.06.09/01}} &
% #VALUE_ID: LEX-P0027-V0004
% #FIELD_VALUE: 批量
% TODO #HANDWRITTEN: 2号-袋; handwriting unclear
\fieldvalue{\handwritten{2号-袋}} \\
\hline
\end{tabularx}

\medskip

\noindent
2）系统运行：使用无菌离心管将细胞冻存液与管路$\mathrm{L1}$号管路相连，并将连接好的细胞液存袋轻轻压下；导管上端；导管上的所有管夹，启动运行程序，开始进行细胞处理，直至中间产品重量完成。

\begin{tabularx}{\textwidth}{|X|>{\raggedright\arraybackslash}p{0.36\textwidth}|}
\hline
& 细胞收获时间：%
% #VALUE_ID: LEX-P0027-V0005
% #FIELD_VALUE: 细胞收获时间
% #HANDWRITTEN: 07：45
\fieldvalue{\handwritten{07：45}} \\[2pt]
& 中间品量完成时间：%
% #VALUE_ID: LEX-P0027-V0006
% #FIELD_VALUE: 中间品量完成时间
% #HANDWRITTEN: 08：20
\fieldvalue{\handwritten{08：20}} \\
\hline
\multicolumn{2}{|l|}{备注：} \\
\hline
\end{tabularx}

\noindent
操作人/日期：
% #VALUE_ID: LEX-P0027-V0007
% #FIELD_VALUE: 操作人
% #HANDWRITTEN: 张娜
\fieldvalue{\handwritten{张娜}}
\quad
% #VALUE_ID: LEX-P0027-V0008
% #FIELD_VALUE: 操作日期
% #HANDWRITTEN: 2025.06.09
\fieldvalue{\handwritten{2025.06.09}}
\hfill
复核人/日期：
% #VALUE_ID: LEX-P0027-V0009
% #FIELD_VALUE: 复核人
% #HANDWRITTEN: 赵晓阳
\fieldvalue{\handwritten{赵晓阳}}
\quad
% #VALUE_ID: LEX-P0027-V0010
% #FIELD_VALUE: 复核日期
% #HANDWRITTEN: 2025.06.09
\fieldvalue{\handwritten{2025.06.09}}

\subsection*{3.3 细胞计数}

\begin{tabularx}{\textwidth}{|>{\centering\arraybackslash}p{0.58\textwidth}|X|}
\hline
\textbf{项目} & \textbf{操作记录} \\
\hline
1）VP设备运行完毕后，记录产品袋内细胞悬液体积（M），使用计算机将产品袋从VP-管路包上融合，使其脱离管路包； &
细胞悬液体积（M）：
% #VALUE_ID: LEX-P0027-V0011
% #FIELD_VALUE: 细胞悬液体积（M）
% #HANDWRITTEN: 48.0
\fieldvalue{\handwritten{48.0}} ml（g） \\
\hline
2）将产品袋经75\%酒精表面消毒后传入洁净工作台，用止血钳把产品袋分成接头取样口的保护帽，拧一次性使用无菌注射器与中间产品袋的鲁尔接头取样并连接，此过程应注意无菌操作；细胞产品袋取样进行检测计数，记录细胞密度（$\rho$）和活率； &
取样时间：
% #VALUE_ID: LEX-P0027-V0012
% #FIELD_VALUE: 取样时间
% #HANDWRITTEN: 07：36
\fieldvalue{\handwritten{07：36}} \\[3pt]
& 细胞密度（$\rho$）：
% #VALUE_ID: LEX-P0027-V0013
% #FIELD_VALUE: 细胞密度（$\rho$）
% #TODO #HANDWRITTEN: 1.23$\times10^7$; exponent partially unclear
\fieldvalue{\handwritten{1.23$\times10^7$}}个细胞/ml; \\[3pt]
& 细胞活率：
% #VALUE_ID: LEX-P0027-V0014
% #FIELD_VALUE: 细胞活率
% #HANDWRITTEN: 97.1
\fieldvalue{\handwritten{97.1}}\% \\
\hline
\multicolumn{2}{|l|}{备注：} \\
\hline
\end{tabularx}

\noindent
操作人/日期：
% #VALUE_ID: LEX-P0027-V0015
% #FIELD_VALUE: 操作人
% #HANDWRITTEN: 张娜
\fieldvalue{\handwritten{张娜}}
\quad
% #VALUE_ID: LEX-P0027-V0016
% #FIELD_VALUE: 操作日期
% #HANDWRITTEN: 2025.06.09
\fieldvalue{\handwritten{2025.06.09}}
\hfill
复核人/日期：
% #VALUE_ID: LEX-P0027-V0017
% #FIELD_VALUE: 复核人
% #HANDWRITTEN: 赵晓阳
\fieldvalue{\handwritten{赵晓阳}}
\quad
% #VALUE_ID: LEX-P0027-V0018
% #FIELD_VALUE: 复核日期
% #HANDWRITTEN: 2025.06.09
\fieldvalue{\handwritten{2025.06.09}}

\subsection*{3.4 一级生物反应器接种制备：}

\begin{tabularx}{\textwidth}{|>{\centering\arraybackslash}X|>{\centering\arraybackslash}X|}
\hline
项目 & 操作记录 \\
\hline
\end{tabularx}

% TODO: This section continues on the following page.
% LEXOID_PAGE_COMPLETED: 27/59

\newpage

\begin{center}
\textbf{艾米近托赛注射液（L2）批生产记录}
\end{center}

\begin{tabularx}{\textwidth}{|p{0.18\textwidth}|X|p{0.18\textwidth}|p{0.18\textwidth}|}
\hline
产品名 & 规格 & 生产批号 & 批量 \\
\hline
% #VALUE_ID: LEX-P0028-V0001
% #FIELD_VALUE: 产品名
\fieldvalue{艾米近托赛注射液}
&
% #VALUE_ID: LEX-P0028-V0002
% #FIELD_VALUE: 规格
\fieldvalue{12ml：$6.0\times10^{7}$个细胞}
&
% #VALUE_ID: LEX-P0028-V0003
% #FIELD_VALUE: 生产批号
% #TODO #HANDWRITTEN: [批号不清]; 蓝色手写内容难以辨认
\fieldvalue{\handwritten{[批号不清]}}
&
% #VALUE_ID: LEX-P0028-V0004
% #FIELD_VALUE: 批量
% #TODO #HANDWRITTEN: 10瓶; 手写字迹不完全清晰
\fieldvalue{\handwritten{10瓶}}
\\
\hline
\end{tabularx}

\vspace{2mm}

\begin{tabularx}{\textwidth}{|p{0.50\textwidth}|X|}
\hline
\textbf{1）}确认产品袋内细胞悬液的活细胞总数在
$(4.0\mathord{-}9.2)\times10^{7}$范围内，即为一级生物反应器细胞接种母液（体积约为
\underline{\hspace{1cm}}\,ml）。

注：活细胞总数（Q）=活细胞密度（$\rho$）$\times$细胞悬液体积（m） &
一级生物反应器细胞密度换液细胞总数（Q）\\
& =活细胞密度（$\rho$）$\times$细胞换母液体积（m）\\
& =
% #VALUE_ID: LEX-P0028-V0005
% #FIELD_VALUE: 活细胞密度
% #TODO #HANDWRITTEN: [约4.0--9.2]$\times10^{7}$个细胞/ml; 原图手写公式数值难以完全辨认
\fieldvalue{\handwritten{[约4.0--9.2]$\times10^{7}$个细胞/ml}}
$\times$
% #VALUE_ID: LEX-P0028-V0006
% #FIELD_VALUE: 细胞换液体积
% #TODO #HANDWRITTEN: 12.0 ml; 手写数值不完全清晰
\fieldvalue{\handwritten{12.0 ml}}
$=$
% #VALUE_ID: LEX-P0028-V0007
% #FIELD_VALUE: 活细胞总数
% #TODO #HANDWRITTEN: [总数不清]; 原图手写结果难以辨认
\fieldvalue{\handwritten{[总数不清]}} \\
& 是 否符合要求：\quad
% #VALUE_ID: LEX-P0028-V0008
% #FIELD_VALUE: 是否符合要求
% #HANDWRITTEN: 是（勾选）
\fieldvalue{\handwritten{是（勾选）}}\quad $\square$否
\\
\hline
\textbf{2）}将重悬液（完全培养基）接种合载断脱胶 VP 管路内，通过接管充接一缓生物反应器细胞换液母液，将接种母液灌进重悬液，总体积
$500$ ml左右，混合均匀，热合封口，即为一级生物反应器细胞接种液。 &
混合时间：
% #VALUE_ID: LEX-P0028-V0009
% #FIELD_VALUE: 混合时间
% #HANDWRITTEN: 08:50
\fieldvalue{\handwritten{08:50}}

细胞换液数量：
% #VALUE_ID: LEX-P0028-V0010
% #FIELD_VALUE: 细胞换液数量
% #HANDWRITTEN: 100 ml
\fieldvalue{\handwritten{100 ml}}
\\
\hline
\textbf{3）}填写《递交单》，将一级生物反应器细胞接种液传入下一步骤； &
转移时间：
% #VALUE_ID: LEX-P0028-V0011
% #FIELD_VALUE: 转移时间
% #HANDWRITTEN: 08:12
\fieldvalue{\handwritten{08:12}}
\\
\hline
备注：
% #VALUE_ID: LEX-P0028-V0012
% #FIELD_VALUE: 备注
% #HANDWRITTEN: /
\fieldvalue{\handwritten{/}}
&
\\
\hline
\end{tabularx}

\vspace{3mm}

\noindent
操作人/日期：
% #VALUE_ID: LEX-P0028-V0013
% #FIELD_VALUE: 操作人
% #TODO #HANDWRITTEN: [姓名不清]; 蓝色手写签名难以辨认
\fieldvalue{\handwritten{[姓名不清]}}
\quad
% #VALUE_ID: LEX-P0028-V0014
% #FIELD_VALUE: 操作日期
% #HANDWRITTEN: 2025.06.09
\fieldvalue{\handwritten{2025.06.09}}
\hfill
复核人/日期：
% #VALUE_ID: LEX-P0028-V0015
% #FIELD_VALUE: 复核人
% #TODO #HANDWRITTEN: [姓名不清]; 蓝色手写签名难以辨认
\fieldvalue{\handwritten{[姓名不清]}}
\quad
% #VALUE_ID: LEX-P0028-V0016
% #FIELD_VALUE: 复核日期
% #HANDWRITTEN: 2025.06.09
\fieldvalue{\handwritten{2025.06.09}}

% TODO: 页眉中的公司标识、记录编号及版本号为固定印刷内容；页面底部页码未重复录入。
% LEXOID_PAGE_COMPLETED: 28/59

\newpage

\begin{center}
\textbf{艾米迈托赛注射液（LZ2）批生产记录}
\end{center}

\noindent
记录编号：
% #VALUE_ID: LEX-P0029-V0001
% #FIELD_VALUE: 记录编号
\fieldvalue{YZ-BPR-MF-PD-01-020-02}
\hfill
版本号：
% #VALUE_ID: LEX-P0029-V0002
% #FIELD_VALUE: 版本号
\fieldvalue{1.3}

\begin{tabularx}{\textwidth}{|>{\centering\arraybackslash}X|>{\centering\arraybackslash}X|>{\centering\arraybackslash}X|>{\centering\arraybackslash}X|}
\hline
品名 & 规格 & 生产批号 & 批量 \\
\hline
艾米迈托赛注射液 & 12ml：$6.0\times10^{7}$个细胞 &
% #VALUE_ID: LEX-P0029-V0003
% #FIELD_VALUE: 生产批号
% #TODO #HANDWRITTEN: A?32?b?0109; handwriting unclear
\fieldvalue{\handwritten{A?32?b?0109}} &
% #VALUE_ID: LEX-P0029-V0004
% #FIELD_VALUE: 批量
% #TODO #HANDWRITTEN: 20-48; handwriting unclear
\fieldvalue{\handwritten{20-48}} \\
\hline
\end{tabularx}

\medskip
\noindent\textbf{4. 清场：}

\medskip
\noindent 操作依据：

\begin{enumerate}
\item 《生产区清场管理规程》 \hfill YZ-SMP-PD-03-001
\item 《生产、检验状态标识管理规程》 \hfill YZ-SMP-QA-04-014
\item 生产工序：艾米迈托赛注射液 P2 代细胞库复苏
\end{enumerate}

\noindent
清场开始时间：
% #VALUE_ID: LEX-P0029-V0005
% #FIELD_VALUE: 清场开始时间
% #HANDWRITTEN: 08时03分
\fieldvalue{\handwritten{08时03分}}
\hfill
清场房间编码：
% #VALUE_ID: LEX-P0029-V0006
% #FIELD_VALUE: 清场房间编码
% #TODO #HANDWRITTEN: 5?8C; handwriting unclear
\fieldvalue{\handwritten{5?8C}}

\noindent
75\%酒精料号：
% #VALUE_ID: LEX-P0029-V0007
% #FIELD_VALUE: 75\%酒精料号
% #TODO #HANDWRITTEN: 20?060701; handwriting unclear
\fieldvalue{\handwritten{20?060701}}

\medskip

\small
\begin{tabularx}{\textwidth}{|X|>{\centering\arraybackslash}p{0.22\textwidth}|}
\hline
\multicolumn{1}{|c|}{清场项目及标准} & \multicolumn{1}{c|}{执行情况} \\
\hline
将本班次使用的各类工具配件恢复原位放置； &
% #VALUE_ID: LEX-P0029-V0008
% #FIELD_VALUE: 执行情况
\fieldvalue{\ensuremath{\boxtimes}\ 已执行\quad $\square$否} \\
\hline
将本班次生产有关的其它物品品出生产现场； &
% #VALUE_ID: LEX-P0029-V0009
% #FIELD_VALUE: 执行情况
\fieldvalue{\ensuremath{\boxtimes}\ 已执行\quad $\square$否} \\
\hline
使用75\%酒精喷洒设备、操作台面、环境领巾、转移窗、门把手等； &
% #VALUE_ID: LEX-P0029-V0010
% #FIELD_VALUE: 执行情况
\fieldvalue{\ensuremath{\boxtimes}\ 已执行\quad $\square$否} \\
\hline
使用75\%酒精喷洒地面，然后使用拖把擦拭地面； &
% #VALUE_ID: LEX-P0029-V0011
% #FIELD_VALUE: 执行情况
\fieldvalue{\ensuremath{\boxtimes}\ 已执行\quad $\square$否} \\
\hline
按照《生产、检验状态标识管理规程》更换房间、设备状态标识； &
% #VALUE_ID: LEX-P0029-V0012
% #FIELD_VALUE: 执行情况
\fieldvalue{\ensuremath{\boxtimes}\ 已执行\quad $\square$否} \\
\hline
将本班次操作产生的废弃物移出洁净生产区； &
% #VALUE_ID: LEX-P0029-V0013
% #FIELD_VALUE: 执行情况
\fieldvalue{\ensuremath{\boxtimes}\ 已执行\quad $\square$否} \\
\hline
\end{tabularx}
\normalsize

\medskip
\noindent
清场结束时间：
% #VALUE_ID: LEX-P0029-V0014
% #FIELD_VALUE: 清场结束时间
% #HANDWRITTEN: 09时05分
\fieldvalue{\handwritten{09时05分}}

\medskip
\noindent
清场人/日期：
% #VALUE_ID: LEX-P0029-V0015
% #FIELD_VALUE: 清场人
% #TODO #HANDWRITTEN: 胡晓贝; handwriting unclear
\fieldvalue{\handwritten{胡晓贝}}
\hspace{1cm}
% #VALUE_ID: LEX-P0029-V0016
% #FIELD_VALUE: 清场日期
% #HANDWRITTEN: 2025.06.09
\fieldvalue{\handwritten{2025.06.09}}

\medskip

\begin{tabularx}{\textwidth}{|X|}
\hline
清场检查 \\
\hline
% #VALUE_ID: LEX-P0029-V0017
% #FIELD_VALUE: 清场检查结果
\fieldvalue{\ensuremath{\boxtimes}\ 合格} \\
\hline
$\square$不合格，具体内容为： \\
\hline
\hspace{1cm}整改情况及整改人： \\
\hline
\end{tabularx}

\medskip
\noindent
确认人/日期：
% #VALUE_ID: LEX-P0029-V0018
% #FIELD_VALUE: 确认人
% #TODO #HANDWRITTEN: 赵姗姗; handwriting unclear
\fieldvalue{\handwritten{赵姗姗}}
\hspace{1cm}
% #VALUE_ID: LEX-P0029-V0019
% #FIELD_VALUE: 确认日期
% #HANDWRITTEN: 2025.06.09
\fieldvalue{\handwritten{2025.06.09}}

\vfill
\begin{center}
第 14 页 共 222 页
\end{center}
% LEXOID_PAGE_COMPLETED: 29/59

\newpage

\begin{center}
\textbf{计数样本交接表}
\end{center}

\noindent
\begin{tabularx}{\textwidth}{lX lX}
品名：&
% #VALUE_ID: LEX-P0030-V0001
% #FIELD_VALUE: 品名
% #TODO #HANDWRITTEN: 盐水片?注射液; handwriting unclear
\fieldvalue{\handwritten{盐水片?注射液}}&
规格：&
% #VALUE_ID: LEX-P0030-V0002
% #FIELD_VALUE: 规格
% #TODO #HANDWRITTEN: 12 ml，6.0$\times$10^{8} cells; handwriting unclear
\fieldvalue{\handwritten{12 ml，$6.0\times10^{8}$ cells}}\\
生产批号：&
% #VALUE_ID: LEX-P0030-V0003
% #FIELD_VALUE: 生产批号
% #TODO #HANDWRITTEN: A?202?0601; handwriting unclear
\fieldvalue{\handwritten{A?202?0601}}&
生产工序：&
% #VALUE_ID: LEX-P0030-V0004
% #FIELD_VALUE: 生产工序
% #TODO #HANDWRITTEN: 支原体检测?及细胞制备; handwriting unclear
\fieldvalue{\handwritten{支原体检测?及细胞制备}}\\
\end{tabularx}

\vspace{2mm}

\begin{tabularx}{\textwidth}{p{0.22\textwidth}|X}
\hline
样品编码 &
% #VALUE_ID: LEX-P0030-V0005
% #FIELD_VALUE: 样品编码
% #TODO #HANDWRITTEN: ZS?061-P?A?202?0601-P3-?; alphanumeric handwriting unclear
\fieldvalue{\handwritten{ZS?061-P?A?202?0601-P3-?}}\\
\hline
样品类型 &
EP管\quad 一管，每管$\geq 0.3$ml\\
\cline{2-2}
& 口取样壶，一壶，每壶$\geq 0.3$ml\\
\hline
送样时间 &
20
% #VALUE_ID: LEX-P0030-V0006
% #FIELD_VALUE: 送样时间
% #TODO #HANDWRITTEN: 25年06月01日09时27分; date handwriting partly unclear
\fieldvalue{\handwritten{25}}年
\fieldvalue{\handwritten{06}}月
\fieldvalue{\handwritten{01}}日
\fieldvalue{\handwritten{09}}时
\fieldvalue{\handwritten{27}}分
\quad
% #VALUE_ID: LEX-P0030-V0007
% #TODO #HANDWRITTEN: Z? 2025.06.09; handwritten note beside date
\handwritten{Z? 2025.06.09}\\
\hline
样本递交人/日期 &
% #VALUE_ID: LEX-P0030-V0008
% #FIELD_VALUE: 样本递交人/日期
% #TODO #HANDWRITTEN: Z?杨中 2025.06.09; handwriting unclear
\fieldvalue{\handwritten{Z?杨中\quad 2025.06.09}}\\
\hline
样本接收人/日期 &
% #VALUE_ID: LEX-P0030-V0009
% #FIELD_VALUE: 样本接收人/日期
% #TODO #HANDWRITTEN: 李? 2025.06.09; handwriting unclear
\fieldvalue{\handwritten{李?\quad 2025.06.09}}\\
\hline
样本接收人/日期 &
% #VALUE_ID: LEX-P0030-V0010
% #FIELD_VALUE: 样本接收人/日期
% #TODO #HANDWRITTEN: 李? 2025.06.09; handwriting unclear
\fieldvalue{\handwritten{李?\quad 2025.06.09}}\\
\hline
样本计数结果 &
活细胞密度：
% #VALUE_ID: LEX-P0030-V0011
% #FIELD_VALUE: 活细胞密度
% #HANDWRITTEN: 1.23
\fieldvalue{\handwritten{1.23}}
$\times 10^{6}$ cells/ml
\\
\cline{2-2}
&
细胞活率：
% #VALUE_ID: LEX-P0030-V0012
% #FIELD_VALUE: 细胞活率
% #HANDWRITTEN: 97
\fieldvalue{\handwritten{97}}\%\\
\hline
计数人/日期 &
% #VALUE_ID: LEX-P0030-V0013
% #FIELD_VALUE: 计数人/日期
% #TODO #HANDWRITTEN: 赵? 2025.06.09; handwriting unclear
\fieldvalue{\handwritten{赵?\quad 2025.06.09}}\\
\hline
递交/复核人/日期 &
% #VALUE_ID: LEX-P0030-V0014
% #FIELD_VALUE: 递交/复核人/日期
% #TODO #HANDWRITTEN: 王? 2025.06.09; handwriting unclear
\fieldvalue{\handwritten{王?\quad 2025.06.09}}\\
\hline
\end{tabularx}

\vspace{5mm}

\begin{center}
\textbf{细胞计数记录}
\end{center}

\begin{tabularx}{\textwidth}{p{0.58\textwidth}|X}
\hline
AO/PI染液批号：&
% #VALUE_ID: LEX-P0030-V0015
% #FIELD_VALUE: AO/PI染液批号
% #TODO #HANDWRITTEN: 20?1/2023?01; alphanumeric handwriting unclear
\fieldvalue{\handwritten{20?1/2023?01}}\\
\hline
细胞计数仪编号：&
% #VALUE_ID: LEX-P0030-V0016
% #FIELD_VALUE: 细胞计数仪编号
% #TODO #HANDWRITTEN: 112?00; handwriting unclear
\fieldvalue{\handwritten{112?00}}\\
\hline
细胞稀释倍数（Dilution Factor）：&
% #VALUE_ID: LEX-P0030-V0017
% #FIELD_VALUE: 细胞稀释倍数
\fieldvalue{%
\raisebox{0.1ex}{\(\checkmark\)}\,2
}\quad
$\square$4\quad $\square$其它\\
\hline
\end{tabularx}

\begin{tabularx}{\textwidth}{|X|c|c|c|c|}
\hline
\multicolumn{1}{|c|}{细胞计数仪内编号码} &
\multicolumn{1}{c|}{活细胞密度\newline(cells/ml)} &
\multicolumn{1}{c|}{细胞\newline 直径} &
\multicolumn{1}{c|}{活细胞直\newline 径（$\mu$m）} &
\multicolumn{1}{c|}{备注}\\
\hline
% #VALUE_ID: LEX-P0030-V0018
% #FIELD_VALUE: 细胞计数仪内编号码
% #TODO #HANDWRITTEN: 1?202?0601-P3-?; alphanumeric handwriting unclear
\fieldvalue{\handwritten{1?202?0601-P3-?}}
&
% #VALUE_ID: LEX-P0030-V0019
% #FIELD_VALUE: 活细胞密度
% #HANDWRITTEN: 1.11
\fieldvalue{\handwritten{1.11}}$\times10^{6}$
&
% #VALUE_ID: LEX-P0030-V0020
% #FIELD_VALUE: 细胞活率
% #HANDWRITTEN: 97.4
\fieldvalue{\handwritten{97.4}}\%
&
% #VALUE_ID: LEX-P0030-V0021
% #FIELD_VALUE: 活细胞直径
% #HANDWRITTEN: 20.1
\fieldvalue{\handwritten{20.1}}
&\\
\hline
% #VALUE_ID: LEX-P0030-V0022
% #FIELD_VALUE: 细胞计数仪内编号码
% #TODO #HANDWRITTEN: 2?202?0601-P3-?; alphanumeric handwriting unclear
\fieldvalue{\handwritten{2?202?0601-P3-?}}
&
% #VALUE_ID: LEX-P0030-V0023
% #FIELD_VALUE: 活细胞密度
% #HANDWRITTEN: 1.21
\fieldvalue{\handwritten{1.21}}$\times10^{6}$
&
% #VALUE_ID: LEX-P0030-V0024
% #FIELD_VALUE: 细胞活率
% #HANDWRITTEN: 97.3
\fieldvalue{\handwritten{97.3}}\%
&
% #VALUE_ID: LEX-P0030-V0025
% #FIELD_VALUE: 活细胞直径
% #HANDWRITTEN: 18.7
\fieldvalue{\handwritten{18.7}}
&\\
\hline
% #VALUE_ID: LEX-P0030-V0026
% #FIELD_VALUE: 细胞计数仪内编号码
% #TODO #HANDWRITTEN: 3?202?0601-P3-?; alphanumeric handwriting unclear
\fieldvalue{\handwritten{3?202?0601-P3-?}}
&
% #VALUE_ID: LEX-P0030-V0027
% #FIELD_VALUE: 活细胞密度
% #TODO #HANDWRITTEN: 1.?; handwriting unclear
\fieldvalue{\handwritten{1.?}}$\times10^{6}$
&
% #VALUE_ID: LEX-P0030-V0028
% #FIELD_VALUE: 细胞活率
% #HANDWRITTEN: 97.3
\fieldvalue{\handwritten{97.3}}\%
&
% #VALUE_ID: LEX-P0030-V0029
% #FIELD_VALUE: 活细胞直径
% #HANDWRITTEN: 20.0
\fieldvalue{\handwritten{20.0}}
&\\
\hline
& $\times10^{6}$ & \% & &\\
\hline
样本计数填写值 &
活细胞密度：
% #VALUE_ID: LEX-P0030-V0030
% #FIELD_VALUE: 活细胞密度
% #HANDWRITTEN: 1.22
\fieldvalue{\handwritten{1.22}}
$\times10^{6}$cells/ml；
细胞活率：
% #VALUE_ID: LEX-P0030-V0031
% #FIELD_VALUE: 细胞活率
% #HANDWRITTEN: 96.78
\fieldvalue{\handwritten{96.78}}\%
& & &\\
\hline
\end{tabularx}

\noindent
如需增加计数款数，可增加附件。

\vspace{2mm}

\begin{tabularx}{\textwidth}{lX lX}
计数人/日期：&
% #VALUE_ID: LEX-P0030-V0032
% #FIELD_VALUE: 计数人/日期
% #TODO #HANDWRITTEN: 赵? 2025.06.09; handwriting unclear
\fieldvalue{\handwritten{赵?\quad 2025.06.09}}&
复核人/日期：&
% #VALUE_ID: LEX-P0030-V0033
% #FIELD_VALUE: 复核人/日期
% #TODO #HANDWRITTEN: 王? 2025.06.09; handwriting unclear
\fieldvalue{\handwritten{王?\quad 2025.06.09}}\\
\end{tabularx}

\begin{center}
第 1 页 共 1 页
\end{center}
% LEXOID_PAGE_COMPLETED: 30/59

\newpage

\begin{center}
\textbf{CellSep\textsuperscript{\textregistered} PRO-3D FloTrix\textsuperscript{\textregistered} vIVAPE-REP细胞收获系统}\\
\textbf{\Large Recovery Cell Wash 工艺}\\
\textbf{\Large 运行报告}
\end{center}

\begin{tabularx}{\textwidth}{@{}lX lX@{}}
设备ID &
% #VALUE_ID: LEX-P0031-V0001
% #FIELD_VALUE: 设备ID
\fieldvalue{SYA1Y4MX03} &
报告日期 &
% #VALUE_ID: LEX-P0031-V0002
% #FIELD_VALUE: 报告日期
\fieldvalue{2025-06-09 08:30:20} \\

设备程序版本 &
% #VALUE_ID: LEX-P0031-V0003
% #FIELD_VALUE: 设备程序版本
\fieldvalue{1.4.87} &
系统界面ID &
% #VALUE_ID: LEX-P0031-V0004
% #FIELD_VALUE: 系统界面ID
\fieldvalue{8595ba-23d9e5} \\

控制程序版本 &
% #VALUE_ID: LEX-P0031-V0005
% #FIELD_VALUE: 控制程序版本
\fieldvalue{1.4.87} &
控制程序ID &
% #VALUE_ID: LEX-P0031-V0006
% #FIELD_VALUE: 控制程序ID
\fieldvalue{2d2643-21bea3} \\

实验编号(\#) &
% #VALUE_ID: LEX-P0031-V0007
% #FIELD_VALUE: 实验编号(\#)
\fieldvalue{A3Z22202506001} &
耗材编号(\#) &
% #VALUE_ID: LEX-P0031-V0008
% #FIELD_VALUE: 耗材编号(\#)
\fieldvalue{H053012025030601} \\

样品编号(\#) &
% #VALUE_ID: LEX-P0031-V0009
% #FIELD_VALUE: 样品编号(\#)
\fieldvalue{NA} &
缓冲液编号(\#) &
% #VALUE_ID: LEX-P0031-V0010
% #FIELD_VALUE: 缓冲液编号(\#)
\fieldvalue{NA}
\end{tabularx}

\vspace{0.5em}

\begin{tabularx}{\textwidth}{@{}lX lX@{}}
工艺名称 &
% #VALUE_ID: LEX-P0031-V0011
% #FIELD_VALUE: 工艺名称
\fieldvalue{Recovery Cell Wash} &
当前用户 &
% #VALUE_ID: LEX-P0031-V0012
% #FIELD_VALUE: 当前用户
\fieldvalue{PC} \\

开始时间 &
% #VALUE_ID: LEX-P0031-V0013
% #FIELD_VALUE: 开始时间
\fieldvalue{2025-06-09 07:36:28} &
结束时间 &
% #VALUE_ID: LEX-P0031-V0014
% #FIELD_VALUE: 结束时间
\fieldvalue{2025-06-09 08:30:13} \\

工艺耗时 &
% #VALUE_ID: LEX-P0031-V0015
% #FIELD_VALUE: 工艺耗时
\fieldvalue{03:54:45} &
工艺状态 &
% #VALUE_ID: LEX-P0031-V0016
% #FIELD_VALUE: 工艺状态
\fieldvalue{工艺结束}
\end{tabularx}

\vspace{0.5em}

\begin{tabularx}{\textwidth}{@{}X c c@{}}
\textbf{参数设置} & \textbf{输入参数值} & \textbf{校验参数值} \\ \hline
空样品重量(g) &
% #VALUE_ID: LEX-P0031-V0017
% #FIELD_VALUE: 空样品重量(g)，输入参数值
\fieldvalue{20} &
% #VALUE_ID: LEX-P0031-V0018
% #FIELD_VALUE: 空样品重量(g)，校验参数值
\fieldvalue{20} \\

样品输入(g) &
% #VALUE_ID: LEX-P0031-V0019
% #FIELD_VALUE: 样品输入(g)，输入参数值
\fieldvalue{11} &
% #VALUE_ID: LEX-P0031-V0020
% #FIELD_VALUE: 样品输入(g)，校验参数值
\fieldvalue{0} \\

样品袋工艺前空袋 &
% #VALUE_ID: LEX-P0031-V0021
% #FIELD_VALUE: 样品袋工艺前空袋，输入参数值
\fieldvalue{0} &
% #VALUE_ID: LEX-P0031-V0022
% #FIELD_VALUE: 样品袋工艺前空袋，校验参数值
\fieldvalue{0} \\

样品稀释 &
% #VALUE_ID: LEX-P0031-V0023
% #FIELD_VALUE: 样品稀释，输入参数值
\fieldvalue{1} &
% #VALUE_ID: LEX-P0031-V0024
% #FIELD_VALUE: 样品稀释，校验参数值
\fieldvalue{1} \\

稀释使用体积(g) &
% #VALUE_ID: LEX-P0031-V0025
% #FIELD_VALUE: 稀释使用体积(g)，输入参数值
\fieldvalue{40} &
% #VALUE_ID: LEX-P0031-V0026
% #FIELD_VALUE: 稀释使用体积(g)，校验参数值
\fieldvalue{40} \\

自动稀释 &
% #VALUE_ID: LEX-P0031-V0027
% #FIELD_VALUE: 自动稀释，输入参数值
\fieldvalue{1} &
% #VALUE_ID: LEX-P0031-V0028
% #FIELD_VALUE: 自动稀释，校验参数值
\fieldvalue{1} \\

样品稀释速度(g/s) &
% #VALUE_ID: LEX-P0031-V0029
% #FIELD_VALUE: 样品稀释速度(g/s)，输入参数值
\fieldvalue{0.5} &
% #VALUE_ID: LEX-P0031-V0030
% #FIELD_VALUE: 样品稀释速度(g/s)，校验参数值
\fieldvalue{0.5} \\

稀释后抽样品 &
% #VALUE_ID: LEX-P0031-V0031
% #FIELD_VALUE: 稀释后抽样品，输入参数值
\fieldvalue{0} &
% #VALUE_ID: LEX-P0031-V0032
% #FIELD_VALUE: 稀释后抽样品，校验参数值
\fieldvalue{0} \\

离心桶预填充滤芯 &
% #VALUE_ID: LEX-P0031-V0033
% #FIELD_VALUE: 离心桶预填充滤芯，输入参数值
\fieldvalue{1} &
% #VALUE_ID: LEX-P0031-V0034
% #FIELD_VALUE: 离心桶预填充滤芯，校验参数值
\fieldvalue{1} \\

预填充洗涤溶液(g) &
% #VALUE_ID: LEX-P0031-V0035
% #FIELD_VALUE: 预填充洗涤溶液(g)，输入参数值
\fieldvalue{180} &
% #VALUE_ID: LEX-P0031-V0036
% #FIELD_VALUE: 预填充洗涤溶液(g)，校验参数值
\fieldvalue{180} \\

清洗样品袋重量(g) &
% #VALUE_ID: LEX-P0031-V0037
% #FIELD_VALUE: 清洗样品袋重量(g)，输入参数值
\fieldvalue{0} &
% #VALUE_ID: LEX-P0031-V0038
% #FIELD_VALUE: 清洗样品袋重量(g)，校验参数值
\fieldvalue{0} \\

样品浓缩次数 &
% #VALUE_ID: LEX-P0031-V0039
% #FIELD_VALUE: 样品浓缩次数，输入参数值
\fieldvalue{1} &
% #VALUE_ID: LEX-P0031-V0040
% #FIELD_VALUE: 样品浓缩次数，校验参数值
\fieldvalue{1} \\

样品最终混匀频率(\#) &
% #VALUE_ID: LEX-P0031-V0041
% #FIELD_VALUE: 样品最终混匀频率(\#)，输入参数值
\fieldvalue{0} &
% #VALUE_ID: LEX-P0031-V0042
% #FIELD_VALUE: 样品最终混匀频率(\#)，校验参数值
\fieldvalue{0} \\

浓缩离心力(g) &
% #VALUE_ID: LEX-P0031-V0043
% #FIELD_VALUE: 浓缩离心力(g)，输入参数值
\fieldvalue{200} &
% #VALUE_ID: LEX-P0031-V0044
% #FIELD_VALUE: 浓缩离心力(g)，校验参数值
\fieldvalue{200} \\

浓缩离心总时间(min) &
% #VALUE_ID: LEX-P0031-V0045
% #FIELD_VALUE: 浓缩离心总时间(min)，输入参数值
\fieldvalue{5} &
% #VALUE_ID: LEX-P0031-V0046
% #FIELD_VALUE: 浓缩离心总时间(min)，校验参数值
\fieldvalue{5}
\end{tabularx}

\vfill

签名：

\begin{tabularx}{\textwidth}{@{}X X X@{}}
\centering
% #VALUE_ID: LEX-P0031-V0047
% #FIELD_VALUE: 生成人
% #TODO #HANDWRITTEN: 杨?; 签名笔迹难以辨认
\fieldvalue{\handwritten{杨?}}
&
\centering
% #VALUE_ID: LEX-P0031-V0048
% #FIELD_VALUE: 审核人
% #TODO #HANDWRITTEN: 赵?; 签名笔迹难以辨认
\fieldvalue{\handwritten{赵?}}
&
\centering
% #VALUE_ID: LEX-P0031-V0049
% #FIELD_VALUE: 日期
% #HANDWRITTEN: 2025.06.09
\fieldvalue{\handwritten{2025.06.09}}
\arraybackslash \\[-0.25em]
\centering 生成人 & \centering 审核人 & \centering 日期 \arraybackslash \\
\end{tabularx}

\begin{flushright}
1-5
\end{flushright}
% LEXOID_PAGE_COMPLETED: 31/59

\newpage

\begin{center}
\textbf{CellSep® PRO-3D FloTrix® vivaPREP细胞收获系统}\\
\textbf{Recovery Cell Wash 工艺}\\
\section*{运行报告}
\end{center}

\begin{tabularx}{\textwidth}{@{}>{\raggedright\arraybackslash}Xcc@{}}
浓缩离心桶内剩余(g) &
% #VALUE_ID: LEX-P0032-V0001
% #FIELD_VALUE: 浓缩离心桶内剩余(g)
\fieldvalue{20} &
% #VALUE_ID: LEX-P0032-V0002
% #FIELD_VALUE: 浓缩离心桶内剩余(g)
\fieldvalue{20} \\

样品复溶中体积(g) &
% #VALUE_ID: LEX-P0032-V0003
% #FIELD_VALUE: 样品复溶中体积(g)
\fieldvalue{0} &
% #VALUE_ID: LEX-P0032-V0004
% #FIELD_VALUE: 样品复溶中体积(g)
\fieldvalue{0} \\

混匀加速(rpm/s) &
% #VALUE_ID: LEX-P0032-V0005
% #FIELD_VALUE: 混匀加速(rpm/s)
\fieldvalue{120} &
% #VALUE_ID: LEX-P0032-V0006
% #FIELD_VALUE: 混匀加速(rpm/s)
\fieldvalue{120} \\

混匀高速速度(rpm) &
% #VALUE_ID: LEX-P0032-V0007
% #FIELD_VALUE: 混匀高速速度(rpm)
\fieldvalue{1000} &
% #VALUE_ID: LEX-P0032-V0008
% #FIELD_VALUE: 混匀高速速度(rpm)
\fieldvalue{1000} \\

混匀减速速度(rpm/s) &
% #VALUE_ID: LEX-P0032-V0009
% #FIELD_VALUE: 混匀减速速度(rpm/s)
\fieldvalue{120} &
% #VALUE_ID: LEX-P0032-V0010
% #FIELD_VALUE: 混匀减速速度(rpm/s)
\fieldvalue{120} \\

混匀间歇时间(s) &
% #VALUE_ID: LEX-P0032-V0011
% #FIELD_VALUE: 混匀间歇时间(s)
\fieldvalue{5} &
% #VALUE_ID: LEX-P0032-V0012
% #FIELD_VALUE: 混匀间歇时间(s)
\fieldvalue{5} \\

混匀次数(\#) &
% #VALUE_ID: LEX-P0032-V0013
% #FIELD_VALUE: 混匀次数(\#)
\fieldvalue{3} &
% #VALUE_ID: LEX-P0032-V0014
% #FIELD_VALUE: 混匀次数(\#)
\fieldvalue{3} \\

混匀方式(\#) &
% #VALUE_ID: LEX-P0032-V0015
% #FIELD_VALUE: 混匀方式(\#)
\fieldvalue{3} &
% #VALUE_ID: LEX-P0032-V0016
% #FIELD_VALUE: 混匀方式(\#)
\fieldvalue{3} \\

样品流速 &
% #VALUE_ID: LEX-P0032-V0017
% #FIELD_VALUE: 样品流速
\fieldvalue{1} &
% #VALUE_ID: LEX-P0032-V0018
% #FIELD_VALUE: 样品流速
\fieldvalue{1} \\

洗涤次数(\#) &
% #VALUE_ID: LEX-P0032-V0019
% #FIELD_VALUE: 洗涤次数(\#)
\fieldvalue{1} &
% #VALUE_ID: LEX-P0032-V0020
% #FIELD_VALUE: 洗涤次数(\#)
\fieldvalue{1} \\

洗涤使用体积(g) &
% #VALUE_ID: LEX-P0032-V0021
% #FIELD_VALUE: 洗涤使用体积(g)
\fieldvalue{220} &
% #VALUE_ID: LEX-P0032-V0022
% #FIELD_VALUE: 洗涤使用体积(g)
\fieldvalue{220} \\

洗涤离心力(xg) &
% #VALUE_ID: LEX-P0032-V0023
% #FIELD_VALUE: 洗涤离心力(xg)
\fieldvalue{200} &
% #VALUE_ID: LEX-P0032-V0024
% #FIELD_VALUE: 洗涤离心力(xg)
\fieldvalue{200} \\

洗涤离心时间(min) &
% #VALUE_ID: LEX-P0032-V0025
% #FIELD_VALUE: 洗涤离心时间(min)
\fieldvalue{5} &
% #VALUE_ID: LEX-P0032-V0026
% #FIELD_VALUE: 洗涤离心时间(min)
\fieldvalue{5} \\

洗涤桶内剩余(g) &
% #VALUE_ID: LEX-P0032-V0027
% #FIELD_VALUE: 洗涤桶内剩余(g)
\fieldvalue{20} &
% #VALUE_ID: LEX-P0032-V0028
% #FIELD_VALUE: 洗涤桶内剩余(g)
\fieldvalue{20} \\

洗涤混匀加速度(rpm/s) &
% #VALUE_ID: LEX-P0032-V0029
% #FIELD_VALUE: 洗涤混匀加速度(rpm/s)
\fieldvalue{120} &
% #VALUE_ID: LEX-P0032-V0030
% #FIELD_VALUE: 洗涤混匀加速度(rpm/s)
\fieldvalue{120} \\

洗涤混匀高速速度(rpm) &
% #VALUE_ID: LEX-P0032-V0031
% #FIELD_VALUE: 洗涤混匀高速速度(rpm)
\fieldvalue{1000} &
% #VALUE_ID: LEX-P0032-V0032
% #FIELD_VALUE: 洗涤混匀高速速度(rpm)
\fieldvalue{1000} \\

洗涤混匀减速速度(rpm/s) &
% #VALUE_ID: LEX-P0032-V0033
% #FIELD_VALUE: 洗涤混匀减速速度(rpm/s)
\fieldvalue{120} &
% #VALUE_ID: LEX-P0032-V0034
% #FIELD_VALUE: 洗涤混匀减速速度(rpm/s)
\fieldvalue{120} \\

洗涤混匀间歇时间(s) &
% #VALUE_ID: LEX-P0032-V0035
% #FIELD_VALUE: 洗涤混匀间歇时间(s)
\fieldvalue{5} &
% #VALUE_ID: LEX-P0032-V0036
% #FIELD_VALUE: 洗涤混匀间歇时间(s)
\fieldvalue{5} \\

洗涤混匀次数(\#) &
% #VALUE_ID: LEX-P0032-V0037
% #FIELD_VALUE: 洗涤混匀次数(\#)
\fieldvalue{4} &
% #VALUE_ID: LEX-P0032-V0038
% #FIELD_VALUE: 洗涤混匀次数(\#)
\fieldvalue{4} \\

洗涤混匀方式(\#) &
% #VALUE_ID: LEX-P0032-V0039
% #FIELD_VALUE: 洗涤混匀方式(\#)
\fieldvalue{3} &
% #VALUE_ID: LEX-P0032-V0040
% #FIELD_VALUE: 洗涤混匀方式(\#)
\fieldvalue{3} \\

使用制剂 &
% #VALUE_ID: LEX-P0032-V0041
% #FIELD_VALUE: 使用制剂
\fieldvalue{1} &
% #VALUE_ID: LEX-P0032-V0042
% #FIELD_VALUE: 使用制剂
\fieldvalue{1} \\

制剂洗涤次数(\#) &
% #VALUE_ID: LEX-P0032-V0043
% #FIELD_VALUE: 制剂洗涤次数(\#)
\fieldvalue{0} &
% #VALUE_ID: LEX-P0032-V0044
% #FIELD_VALUE: 制剂洗涤次数(\#)
\fieldvalue{0} \\

制剂洗涤使用体积(g) &
% #VALUE_ID: LEX-P0032-V0045
% #FIELD_VALUE: 制剂洗涤使用体积(g)
\fieldvalue{100} &
% #VALUE_ID: LEX-P0032-V0046
% #FIELD_VALUE: 制剂洗涤使用体积(g)
\fieldvalue{100} \\

制剂洗涤离心力(xg) &
% #VALUE_ID: LEX-P0032-V0047
% #FIELD_VALUE: 制剂洗涤离心力(xg)
\fieldvalue{200} &
% #VALUE_ID: LEX-P0032-V0048
% #FIELD_VALUE: 制剂洗涤离心力(xg)
\fieldvalue{200} \\

制剂洗涤离心时间(min) &
% #VALUE_ID: LEX-P0032-V0049
% #FIELD_VALUE: 制剂洗涤离心时间(min)
\fieldvalue{2} &
% #VALUE_ID: LEX-P0032-V0050
% #FIELD_VALUE: 制剂洗涤离心时间(min)
\fieldvalue{2} \\

制剂洗涤桶内剩余(g) &
% #VALUE_ID: LEX-P0032-V0051
% #FIELD_VALUE: 制剂洗涤桶内剩余(g)
\fieldvalue{20} &
% #VALUE_ID: LEX-P0032-V0052
% #FIELD_VALUE: 制剂洗涤桶内剩余(g)
\fieldvalue{20}
\end{tabularx}

\vspace{1.5em}

签名：

\begin{tabularx}{\textwidth}{@{}XXXX@{}}
\centering
% #VALUE_ID: LEX-P0032-V0053
% #FIELD_VALUE: 生产人
% #HANDWRITTEN: 张帅
\fieldvalue{\handwritten{张帅}}
&
\centering
% #VALUE_ID: LEX-P0032-V0054
% #FIELD_VALUE: 生产日期
% #HANDWRITTEN: 2025.06.09
\fieldvalue{\handwritten{2025.06.09}}
&
\centering
% #VALUE_ID: LEX-P0032-V0055
% #FIELD_VALUE: 审核人
% #HANDWRITTEN: 赵热旭
\fieldvalue{\handwritten{赵热旭}}
&
\centering
% #VALUE_ID: LEX-P0032-V0056
% #FIELD_VALUE: 审核日期
% #HANDWRITTEN: 2025.06.09
\fieldvalue{\handwritten{2025.06.09}}
\arraybackslash \\[-0.3em]
\cline{1-4}
\multicolumn{1}{c}{生产人} &
\multicolumn{1}{c}{生产日期} &
\multicolumn{1}{c}{审核人} &
\multicolumn{1}{c}{日期} \\
\end{tabularx}

\vfill

\begin{center}
\small 2025-06-09 07:31-Recovery Cell Wash.pdf \hfill 2-5
\end{center}
% LEXOID_PAGE_COMPLETED: 32/59

\newpage

\begin{center}
CellSep\textsuperscript{\textregistered} PRO-3D FloTrix\textsuperscript{\textregistered} vivaPREP细胞收获系统\\
\textbf{Recovery Cell Wash 工艺}\\
\textbf{运行报告}
\end{center}

\begin{tabularx}{\textwidth}{@{}>{\raggedright\arraybackslash}Xcc@{}}
制剂洗涤混匀加速度(rpm/s) &
% #VALUE_ID: LEX-P0033-V0001
% #FIELD_VALUE: 制剂洗涤混匀加速度(rpm/s)
\fieldvalue{500} &
% #VALUE_ID: LEX-P0033-V0002
% #FIELD_VALUE: 制剂洗涤混匀加速度(rpm/s)
\fieldvalue{500} \\
制剂洗涤混匀离心速度(rpm) &
% #VALUE_ID: LEX-P0033-V0003
% #FIELD_VALUE: 制剂洗涤混匀离心速度(rpm)
\fieldvalue{1000} &
% #VALUE_ID: LEX-P0033-V0004
% #FIELD_VALUE: 制剂洗涤混匀离心速度(rpm)
\fieldvalue{1000} \\
制剂洗涤混匀减速度(rpm/s) &
% #VALUE_ID: LEX-P0033-V0005
% #FIELD_VALUE: 制剂洗涤混匀减速度(rpm/s)
\fieldvalue{500} &
% #VALUE_ID: LEX-P0033-V0006
% #FIELD_VALUE: 制剂洗涤混匀减速度(rpm/s)
\fieldvalue{500} \\
制剂洗涤混匀间歇时间(s) &
% #VALUE_ID: LEX-P0033-V0007
% #FIELD_VALUE: 制剂洗涤混匀间歇时间(s)
\fieldvalue{5} &
% #VALUE_ID: LEX-P0033-V0008
% #FIELD_VALUE: 制剂洗涤混匀间歇时间(s)
\fieldvalue{5} \\
制剂洗涤混匀次数(\#) &
% #VALUE_ID: LEX-P0033-V0009
% #FIELD_VALUE: 制剂洗涤混匀次数(\#)
\fieldvalue{4} &
% #VALUE_ID: LEX-P0033-V0010
% #FIELD_VALUE: 制剂洗涤混匀次数(\#)
\fieldvalue{4} \\
制剂洗涤混匀方式(\#) &
% #VALUE_ID: LEX-P0033-V0011
% #FIELD_VALUE: 制剂洗涤混匀方式(\#)
\fieldvalue{3} &
% #VALUE_ID: LEX-P0033-V0012
% #FIELD_VALUE: 制剂洗涤混匀方式(\#)
\fieldvalue{3} \\
终产物(g) &
% #VALUE_ID: LEX-P0033-V0013
% #FIELD_VALUE: 终产物(g)
\fieldvalue{50} &
% #VALUE_ID: LEX-P0033-V0014
% #FIELD_VALUE: 终产物(g)
\fieldvalue{50} \\
手动装载 &
% #VALUE_ID: LEX-P0033-V0015
% #FIELD_VALUE: 手动装载
\fieldvalue{1} &
% #VALUE_ID: LEX-P0033-V0016
% #FIELD_VALUE: 手动装载
\fieldvalue{1} \\
重悬混匀加速度(rpm/s) &
% #VALUE_ID: LEX-P0033-V0017
% #FIELD_VALUE: 重悬混匀加速度(rpm/s)
\fieldvalue{120} &
% #VALUE_ID: LEX-P0033-V0018
% #FIELD_VALUE: 重悬混匀加速度(rpm/s)
\fieldvalue{120} \\
重悬混匀离心速度(rpm) &
% #VALUE_ID: LEX-P0033-V0019
% #FIELD_VALUE: 重悬混匀离心速度(rpm)
\fieldvalue{1000} &
% #VALUE_ID: LEX-P0033-V0020
% #FIELD_VALUE: 重悬混匀离心速度(rpm)
\fieldvalue{1000} \\
重悬混匀减速度(rpm/s) &
% #VALUE_ID: LEX-P0033-V0021
% #FIELD_VALUE: 重悬混匀减速度(rpm/s)
\fieldvalue{120} &
% #VALUE_ID: LEX-P0033-V0022
% #FIELD_VALUE: 重悬混匀减速度(rpm/s)
\fieldvalue{120} \\
重悬混匀间歇时间(s) &
% #VALUE_ID: LEX-P0033-V0023
% #FIELD_VALUE: 重悬混匀间歇时间(s)
\fieldvalue{5} &
% #VALUE_ID: LEX-P0033-V0024
% #FIELD_VALUE: 重悬混匀间歇时间(s)
\fieldvalue{5} \\
重悬混匀次数(\#) &
% #VALUE_ID: LEX-P0033-V0025
% #FIELD_VALUE: 重悬混匀次数(\#)
\fieldvalue{2} &
% #VALUE_ID: LEX-P0033-V0026
% #FIELD_VALUE: 重悬混匀次数(\#)
\fieldvalue{2} \\
重悬混匀方式(\#) &
% #VALUE_ID: LEX-P0033-V0027
% #FIELD_VALUE: 重悬混匀方式(\#)
\fieldvalue{3} &
% #VALUE_ID: LEX-P0033-V0028
% #FIELD_VALUE: 重悬混匀方式(\#)
\fieldvalue{3}
\end{tabularx}

\vspace{1em}

\begin{tabularx}{\textwidth}{@{}>{\raggedright\arraybackslash}Xl@{}}
\textbf{运行结果} & \textbf{参数值} \\
样品输入(g) &
% #VALUE_ID: LEX-P0033-V0029
% #FIELD_VALUE: 样品输入(g)
\fieldvalue{46.1} \\
预填充洗涤液浓度(g) &
% #VALUE_ID: LEX-P0033-V0030
% #FIELD_VALUE: 预填充洗涤液浓度(g)
\fieldvalue{179.5} \\
样品稀释使用体积(g) &
% #VALUE_ID: LEX-P0033-V0031
% #FIELD_VALUE: 样品稀释使用体积(g)
\fieldvalue{39.6} \\
第1次洗涤液体积(g) &
% #VALUE_ID: LEX-P0033-V0032
% #FIELD_VALUE: 第1次洗涤液体积(g)
\fieldvalue{212.3} \\
第1次洗涤液弃体积(g) &
% #VALUE_ID: LEX-P0033-V0033
% #FIELD_VALUE: 第1次洗涤液弃体积(g)
\fieldvalue{211.7} \\
终产物(g) &
% #VALUE_ID: LEX-P0033-V0034
% #FIELD_VALUE: 终产物(g)
\fieldvalue{48.5} \\
第1次重悬体积(g) &
% #VALUE_ID: LEX-P0033-V0035
% #FIELD_VALUE: 第1次重悬体积(g)
\fieldvalue{21.5} \\
第2次重悬体积(g) &
% #VALUE_ID: LEX-P0033-V0036
% #FIELD_VALUE: 第2次重悬体积(g)
\fieldvalue{27.0}
\end{tabularx}

\vspace{1em}

\begin{tabularx}{\textwidth}{@{}>{\raggedright\arraybackslash}Xl@{}}
\textbf{步骤名} & \textbf{耗时} \\
\end{tabularx}

\vfill

签名：

\begin{tabularx}{\textwidth}{@{}XXX@{}}
\centering
% #VALUE_ID: LEX-P0033-V0037
% #FIELD_VALUE: 生产人
% TODO #HANDWRITTEN: 张旭; handwritten signature is unclear
\fieldvalue{\handwritten{张旭}}
&
\centering
% #VALUE_ID: LEX-P0033-V0038
% #FIELD_VALUE: 审核人
% TODO #HANDWRITTEN: 赵某; handwritten signature is unclear
\fieldvalue{\handwritten{赵某}}
&
\centering
% #VALUE_ID: LEX-P0033-V0039
% #FIELD_VALUE: 日期
% #HANDWRITTEN: 2025.06.09
\fieldvalue{\handwritten{2025.06.09}} \arraybackslash \\[0.3em]
\centering 生产人 & \centering 审核人 & \centering 日期 \arraybackslash \\
\end{tabularx}
% LEXOID_PAGE_COMPLETED: 33/59

\newpage

\begin{center}
{\small CellSep\textsuperscript{\textregistered} PRO-3D FloTrix\textsuperscript{\textregistered} vivaPREP\textsuperscript{\textregistered}细胞收集系统}\\
\textbf{Recovery Cell Wash 工艺}\\
\textbf{运行报告}
\end{center}

\begin{tabularx}{\textwidth}{@{}lX@{}}
系统排气 &
% #VALUE_ID: LEX-P0034-V0001
% #FIELD_VALUE: 系统排气
\fieldvalue{1min 45s}\\
\hline
样品预处理 &
% #VALUE_ID: LEX-P0034-V0002
% #FIELD_VALUE: 样品预处理
\fieldvalue{10min 15s}\\
\hline
样品浓缩 &
% #VALUE_ID: LEX-P0034-V0003
% #FIELD_VALUE: 样品浓缩
\fieldvalue{16min 39s}\\
\hline
样品洗涤 &
% #VALUE_ID: LEX-P0034-V0004
% #FIELD_VALUE: 样品洗涤
\fieldvalue{19min 39s}\\
\hline
制剂液洗涤 &\\
\hline
产品重悬 &
% #VALUE_ID: LEX-P0034-V0005
% #FIELD_VALUE: 产品重悬
\fieldvalue{5min 22s}
\end{tabularx}

\vfill

签名：

\begin{tabularx}{0.85\textwidth}{@{}>{\centering\arraybackslash}X
>{\centering\arraybackslash}X
>{\centering\arraybackslash}X@{}}
\hrulefill & \hrulefill & \hrulefill\\[-1.2ex]
% #VALUE_ID: LEX-P0034-V0006
% #FIELD_VALUE: 生产人签名
% #TODO #HANDWRITTEN: 难辨；蓝色手写签名笔迹不清
\fieldvalue{\handwritten{难辨}} &
% #VALUE_ID: LEX-P0034-V0008
% #FIELD_VALUE: 审核人签名
% #TODO #HANDWRITTEN: 难辨；蓝色手写签名笔迹不清
\fieldvalue{\handwritten{难辨}} &
\\[1ex]
生产人 &
% #VALUE_ID: LEX-P0034-V0007
% #FIELD_VALUE: 生产人日期
% #TODO #HANDWRITTEN: 2025.06.07；日期笔迹部分不清
\fieldvalue{\handwritten{2025.06.07}} &
审核人\\
&
&
% #VALUE_ID: LEX-P0034-V0009
% #FIELD_VALUE: 审核人日期
% #TODO #HANDWRITTEN: 2025.06.09；日期笔迹部分不清
\fieldvalue{\handwritten{2025.06.09}}
\end{tabularx}

\noindent\footnotesize 2025-05-09 07:31-Recovery Cell Wash.pdf \hfill 4-5
% LEXOID_PAGE_COMPLETED: 34/59

\newpage

\begin{center}
\textbf{CellSep\textsuperscript{\textregistered} PRO-3D FloTrix\textsuperscript{\textregistered} vivaPREP细胞收获系统}\\
\textbf{Recovery Cell Wash 工艺}\\[2pt]
{\Large\textbf{运行报告}}
\end{center}

\vspace{0.5em}

\begin{center}
\textbf{称重传感器信息}
\end{center}

% TODO: The chart is visible on this page, but no source image filename is available.
\begin{center}
\fbox{\parbox{0.86\textwidth}{\centering
\textit{称重传感器信息图表占位符}\\[2pt]
图例：柱体、制剂液袋、缓冲液袋、废液袋、产品袋、样品液体\\
横轴刻度：15:00、30:00、45:00
}}
\end{center}

\vspace{0.8em}

\begin{center}
\textbf{离心桶压力}
\end{center}

% TODO: The chart is visible on this page, but no source image filename is available.
\begin{center}
\fbox{\parbox{0.86\textwidth}{\centering
\textit{离心桶压力图表占位符}\\[2pt]
纵轴：压力（kPa）\qquad 横轴刻度：15:00、30:00、45:00
}}
\end{center}

\vspace{0.8em}

\begin{center}
\textbf{离心力}
\end{center}

% TODO: The chart is visible on this page, but no source image filename is available.
\begin{center}
\fbox{\parbox{0.86\textwidth}{\centering
\textit{离心力图表占位符}\\[2pt]
横轴刻度：15:00、30:00、45:00
}}
\end{center}

\vfill

签名：

\begin{center}
\begin{tabular}{@{}c@{\hspace{1.5em}}c@{\hspace{1.5em}}c@{\hspace{1.5em}}c@{}}
\begin{minipage}{0.18\textwidth}\centering
% #VALUE_ID: LEX-P0035-V0001
% #FIELD_VALUE: 生产人签名
% #TODO #HANDWRITTEN: 张帅; 签名笔迹难以清晰辨认
\fieldvalue{\handwritten{张帅}}\\[-2pt]
\rule{0.18\textwidth}{0.4pt}\\[-1pt]
生产人
\end{minipage}
&
\begin{minipage}{0.18\textwidth}\centering
% #VALUE_ID: LEX-P0035-V0002
% #FIELD_VALUE: 生产人日期
% #HANDWRITTEN: 2025.06.09
\fieldvalue{\handwritten{2025.06.09}}\\[-2pt]
\rule{0.18\textwidth}{0.4pt}\\[-1pt]
日期
\end{minipage}
&
\begin{minipage}{0.18\textwidth}\centering
% #VALUE_ID: LEX-P0035-V0003
% #FIELD_VALUE: 审核人签名
% #TODO #HANDWRITTEN: 赵立; 签名笔迹难以清晰辨认
\fieldvalue{\handwritten{赵立}}\\[-2pt]
\rule{0.18\textwidth}{0.4pt}\\[-1pt]
审核人
\end{minipage}
&
\begin{minipage}{0.18\textwidth}\centering
% #VALUE_ID: LEX-P0035-V0004
% #FIELD_VALUE: 审核人日期
% #HANDWRITTEN: 2025.06.01
\fieldvalue{\handwritten{2025.06.01}}\\[-2pt]
\rule{0.18\textwidth}{0.4pt}\\[-1pt]
日期
\end{minipage}
\end{tabular}
\end{center}
% LEXOID_PAGE_COMPLETED: 35/59

\newpage

\noindent Assay:
% #VALUE_ID: LEX-P0036-V0001
% #FIELD_VALUE: Assay
\fieldvalue{MSC}
\hfill
% #VALUE_ID: LEX-P0036-V0002
% #FIELD_VALUE: Date and time
\fieldvalue{06/09/2025 08:42:14}

\medskip

\noindent Cell Type F1:
% #VALUE_ID: LEX-P0036-V0003
% #FIELD_VALUE: Cell Type F1
\fieldvalue{MSC (AO)}

\noindent Cell Type F2:
% #VALUE_ID: LEX-P0036-V0004
% #FIELD_VALUE: Cell Type F2
\fieldvalue{MSC (PI)}

\medskip

\noindent Sample ID:
% #VALUE_ID: LEX-P0036-V0005
% #FIELD_VALUE: Sample ID
\fieldvalue{2506090-A33220250601-P3-jiezhong-48. 1ml-1-1-1}

\noindent Dilution:
% #VALUE_ID: LEX-P0036-V0006
% #FIELD_VALUE: Dilution
\fieldvalue{2.00}

\bigskip
\noindent\textbf{Results:}

\medskip

\begin{tabularx}{\textwidth}{@{}lXX@{}}
\textbf{Count} & \textbf{Concentration} & \textbf{Mean Diameter} \\[2pt]
\hline
Total:
% #VALUE_ID: LEX-P0036-V0007
% #FIELD_VALUE: Total count
\fieldvalue{358}
&
% #VALUE_ID: LEX-P0036-V0008
% #FIELD_VALUE: Total concentration
\fieldvalue{$1.29\times 10^{6}$ cells/mL}
&
% #VALUE_ID: LEX-P0036-V0009
% #FIELD_VALUE: Total mean diameter
\fieldvalue{19.5 micron}
\\
Live:
% #VALUE_ID: LEX-P0036-V0010
% #FIELD_VALUE: Live count
\fieldvalue{330}
&
% #VALUE_ID: LEX-P0036-V0011
% #FIELD_VALUE: Live concentration
\fieldvalue{$1.19\times 10^{6}$ cells/mL}
&
% #VALUE_ID: LEX-P0036-V0012
% #FIELD_VALUE: Live mean diameter
\fieldvalue{20.1 micron}
\\
Dead:
% #VALUE_ID: LEX-P0036-V0013
% #FIELD_VALUE: Dead count
\fieldvalue{28}
&
% #VALUE_ID: LEX-P0036-V0014
% #FIELD_VALUE: Dead concentration
\fieldvalue{$9.80\times 10^{4}$ cells/mL}
&
% #VALUE_ID: LEX-P0036-V0015
% #FIELD_VALUE: Dead mean diameter
\fieldvalue{11.6 micron}
\end{tabularx}

\medskip

Viability:
% #VALUE_ID: LEX-P0036-V0016
% #FIELD_VALUE: Viability
\fieldvalue{92.4\%}

\bigskip

\begin{center}
\begin{tabular}{cc}
% TODO: Image visible on the page; source filename unavailable.
\fbox{\parbox[c][3.2cm][c]{0.38\textwidth}{\centering Image placeholder}}
&
% TODO: Image visible on the page; source filename unavailable.
\fbox{\parbox[c][3.2cm][c]{0.38\textwidth}{\centering Image placeholder}}
\\[1.5cm]
% TODO: Image visible on the page; source filename unavailable.
\fbox{\parbox[c][3.2cm][c]{0.38\textwidth}{\centering Image placeholder}}
&
% TODO: Image visible on the page; source filename unavailable.
\fbox{\parbox[c][3.2cm][c]{0.38\textwidth}{\centering Image placeholder}}
\end{tabular}
\end{center}

\vfill

\noindent 操作人：
% #VALUE_ID: LEX-P0036-V0017
% #FIELD_VALUE: 操作人
% #TODO #HANDWRITTEN: 赵立帅; cursive handwriting is partly unclear
\fieldvalue{\handwritten{赵立帅}}
\hfill
日期：
% #VALUE_ID: LEX-P0036-V0019
% #FIELD_VALUE: 日期
% #HANDWRITTEN: 2025.06.09
\fieldvalue{\handwritten{2025.06.09}}

\noindent 复核人：
% #VALUE_ID: LEX-P0036-V0018
% #FIELD_VALUE: 复核人
% #TODO #HANDWRITTEN: 张华; cursive handwriting is partly unclear
\fieldvalue{\handwritten{张华}}
\hfill
日期：
% #VALUE_ID: LEX-P0036-V0020
% #FIELD_VALUE: 日期
% #HANDWRITTEN: 2025.06.09
\fieldvalue{\handwritten{2025.06.09}}
% LEXOID_PAGE_COMPLETED: 36/59

\newpage

\noindent
\begin{tabularx}{\textwidth}{@{}Xr@{}}
Assay: MSC & 
% #VALUE_ID: LEX-P0037-V0001
% #FIELD_VALUE: Date and time
\fieldvalue{06/09/2025 08:45:00}
\end{tabularx}

\noindent
% #VALUE_ID: LEX-P0037-V0002
% #FIELD_VALUE: Assay
Assay: \fieldvalue{MSC}

\noindent
% #VALUE_ID: LEX-P0037-V0003
% #FIELD_VALUE: Cell Type F1
Cell Type F1: \fieldvalue{MSC (AO)}

\noindent
% #VALUE_ID: LEX-P0037-V0004
% #FIELD_VALUE: Cell Type F2
Cell Type F2: \fieldvalue{MSC (PI)}

\noindent
% #VALUE_ID: LEX-P0037-V0005
% #FIELD_VALUE: Sample ID
Sample ID: \fieldvalue{\texttt{250609-A3322202506001-P3-jiezhong-48.5ml-1-2-1}}

\noindent
% #VALUE_ID: LEX-P0037-V0006
% #FIELD_VALUE: Dilution
Dilution: \fieldvalue{2.00}

\vspace{0.5em}
\hrule
\vspace{0.8em}

\noindent Results:

\vspace{0.4em}

\noindent
\begin{tabularx}{\textwidth}{@{}XXX@{}}
\textbf{Count} & \textbf{Concentration} & \textbf{Mean Diameter} \\
\hline
Total:
% #VALUE_ID: LEX-P0037-V0007
% #FIELD_VALUE: Total count
\fieldvalue{392}
&
% #VALUE_ID: LEX-P0037-V0008
% #FIELD_VALUE: Total concentration
\fieldvalue{$1.40\times10^{6}$ cells/mL}
&
% #VALUE_ID: LEX-P0037-V0009
% #FIELD_VALUE: Total mean diameter
\fieldvalue{17.6 micron}
\\
Live:
% #VALUE_ID: LEX-P0037-V0010
% #FIELD_VALUE: Live count
\fieldvalue{361}
&
% #VALUE_ID: LEX-P0037-V0011
% #FIELD_VALUE: Live concentration
\fieldvalue{$1.29\times10^{6}$ cells/mL}
&
% #VALUE_ID: LEX-P0037-V0012
% #FIELD_VALUE: Live mean diameter
\fieldvalue{18.2 micron}
\\
Dead:
% #VALUE_ID: LEX-P0037-V0013
% #FIELD_VALUE: Dead count
\fieldvalue{31}
&
% #VALUE_ID: LEX-P0037-V0014
% #FIELD_VALUE: Dead concentration
\fieldvalue{$1.08\times10^{5}$ cells/mL}
&
% #VALUE_ID: LEX-P0037-V0015
% #FIELD_VALUE: Dead mean diameter
\fieldvalue{11.4 micron}
\end{tabularx}

\vspace{0.6em}

% #VALUE_ID: LEX-P0037-V0016
% #FIELD_VALUE: Viability
Viability: \fieldvalue{92.3\%}

\vspace{1em}

\noindent
\begin{tabular}{@{}c@{\hspace{5em}}c@{}}
% TODO: Image visible on page; source filename unavailable.
\fbox{\parbox[c][3.3cm][c]{0.29\textwidth}{\centering Image placeholder}}
&
% TODO: Image visible on page; source filename unavailable.
\fbox{\parbox[c][3.3cm][c]{0.29\textwidth}{\centering Image placeholder}}
\\[2em]
% TODO: Image visible on page; source filename unavailable.
\fbox{\parbox[c][3.3cm][c]{0.29\textwidth}{\centering Image placeholder}}
&
% TODO: Image visible on page; source filename unavailable.
\fbox{\parbox[c][3.3cm][c]{0.29\textwidth}{\centering Image placeholder}}
\end{tabular}

\vfill

\noindent
\begin{tabularx}{\textwidth}{@{}X X@{}}
操作人：
% #VALUE_ID: LEX-P0037-V0017
% #FIELD_VALUE: 操作人
% TODO #HANDWRITTEN: 赵小峰; handwritten Chinese name is partially illegible
\fieldvalue{\handwritten{赵小峰}}
&
日期：
% #VALUE_ID: LEX-P0037-V0018
% #FIELD_VALUE: 操作人日期
% #HANDWRITTEN: 2025.06.09
\fieldvalue{\handwritten{2025.06.09}}
\\[0.8em]
复核人：
% #VALUE_ID: LEX-P0037-V0019
% #FIELD_VALUE: 复核人
% TODO #HANDWRITTEN: 张小中; handwritten Chinese name is partially illegible
\fieldvalue{\handwritten{张小中}}
&
日期：
% #VALUE_ID: LEX-P0037-V0020
% #FIELD_VALUE: 复核人日期
% #HANDWRITTEN: 2025.06.09
\fieldvalue{\handwritten{2025.06.09}}
\end{tabularx}
% LEXOID_PAGE_COMPLETED: 37/59

\newpage

\noindent
Assay: \\
% #VALUE_ID: LEX-P0038-V0001
% #FIELD_VALUE: Assay
\fieldvalue{MSC}

\hfill
% #VALUE_ID: LEX-P0038-V0002
% #FIELD_VALUE: Report date and time
\fieldvalue{06/09/2025 08:48:00}

\medskip
\noindent
Cell Type F1: 
% #VALUE_ID: LEX-P0038-V0003
% #FIELD_VALUE: Cell Type F1
\fieldvalue{MSC (AO)}

\medskip
\noindent
Cell Type F2: 
% #VALUE_ID: LEX-P0038-V0004
% #FIELD_VALUE: Cell Type F2
\fieldvalue{MSC (PI)}

\medskip
\noindent
Sample ID: 
% #VALUE_ID: LEX-P0038-V0005
% #FIELD_VALUE: Sample ID
\fieldvalue{250609-A3372220506001-P3-jiezhong-48.sml-1-3-1}

\medskip
\noindent
Dilution: 
% #VALUE_ID: LEX-P0038-V0006
% #FIELD_VALUE: Dilution
\fieldvalue{2.00}

\bigskip
\noindent
Results:

\medskip
\renewcommand{\arraystretch}{1.15}
\begin{tabularx}{\textwidth}{@{}l X X@{}}
\textbf{Count} & \textbf{Concentration} & \textbf{Mean Diameter} \\
\hline
Total:
% #VALUE_ID: LEX-P0038-V0007
% #FIELD_VALUE: Total count
\fieldvalue{361}
&
% #VALUE_ID: LEX-P0038-V0008
% #FIELD_VALUE: Total concentration
\fieldvalue{$1.30\times10^{6}$ cells/mL}
&
% #VALUE_ID: LEX-P0038-V0009
% #FIELD_VALUE: Total mean diameter
\fieldvalue{19.4 micron}
\\
Live:
% #VALUE_ID: LEX-P0038-V0010
% #FIELD_VALUE: Live count
\fieldvalue{336}
&
% #VALUE_ID: LEX-P0038-V0011
% #FIELD_VALUE: Live concentration
\fieldvalue{$1.21\times10^{6}$ cells/mL}
&
% #VALUE_ID: LEX-P0038-V0012
% #FIELD_VALUE: Live mean diameter
\fieldvalue{20.0 micron}
\\
Dead:
% #VALUE_ID: LEX-P0038-V0013
% #FIELD_VALUE: Dead count
\fieldvalue{25}
&
% #VALUE_ID: LEX-P0038-V0014
% #FIELD_VALUE: Dead concentration
\fieldvalue{$8.70\times10^{4}$ cells/mL}
&
% #VALUE_ID: LEX-P0038-V0015
% #FIELD_VALUE: Dead mean diameter
\fieldvalue{11.5 micron}
\end{tabularx}

\medskip
\noindent
Viability:
% #VALUE_ID: LEX-P0038-V0016
% #FIELD_VALUE: Viability
\fieldvalue{93.3\%}

\bigskip
\begin{center}
\begin{minipage}[t]{0.43\textwidth}
\centering
% TODO: Image visible on this page; no source filename is available.
\fbox{\rule{0pt}{0.19\textheight}\rule{0.95\linewidth}{0pt}}
\end{minipage}
\hfill
\begin{minipage}[t]{0.43\textwidth}
\centering
% TODO: Image visible on this page; no source filename is available.
\fbox{\rule{0pt}{0.19\textheight}\rule{0.95\linewidth}{0pt}}
\end{minipage}

\vspace{1.5cm}

\begin{minipage}[t]{0.43\textwidth}
\centering
% TODO: Image visible on this page; no source filename is available.
\fbox{\rule{0pt}{0.19\textheight}\rule{0.95\linewidth}{0pt}}
\end{minipage}
\hfill
\begin{minipage}[t]{0.43\textwidth}
\centering
% TODO: Image visible on this page; no source filename is available.
\fbox{\rule{0pt}{0.19\textheight}\rule{0.95\linewidth}{0pt}}
\end{minipage}
\end{center}

\vfill

\noindent
操作人：%
% #VALUE_ID: LEX-P0038-V0017
% #FIELD_VALUE: 操作人
% #TODO #HANDWRITTEN: 赵永山; handwriting is partially unclear
\fieldvalue{\handwritten{赵永山}}

\hfill
日期：%
% #VALUE_ID: LEX-P0038-V0018
% #FIELD_VALUE: 操作人日期
% #HANDWRITTEN: 2025.06.09
\fieldvalue{\handwritten{2025.06.09}}

\medskip
\noindent
复核人：%
% #VALUE_ID: LEX-P0038-V0019
% #FIELD_VALUE: 复核人
% #TODO #HANDWRITTEN: 张仲; handwriting is partially unclear
\fieldvalue{\handwritten{张仲}}

\hfill
日期：%
% #VALUE_ID: LEX-P0038-V0020
% #FIELD_VALUE: 复核人日期
% #HANDWRITTEN: 2025.06.09
\fieldvalue{\handwritten{2025.06.09}}
% LEXOID_PAGE_COMPLETED: 38/59

\newpage

\begin{center}
{\large \textbf{艾米迈托赛注射液（Z2）批生产记录}}
\end{center}

\noindent
华中卓越生物科技（北京）有限公司\hfill 记录编号：YZ-BPR-MF-PD-01-020-a

\begin{tabularx}{\textwidth}{|>{\centering\arraybackslash}X|>{\centering\arraybackslash}X|>{\centering\arraybackslash}X|>{\centering\arraybackslash}X|}
\hline
品名 & 规格 & 生产批号 & 批量\\
\hline
艾米迈托赛注射液 & 12ml：$6.0\times10^7$个细胞 &
% #VALUE_ID: LEX-P0039-V0001
% #FIELD_VALUE: 生产批号
% #TODO #HANDWRITTEN: B3220226?60?; handwriting unclear
\fieldvalue{\handwritten{B3220226?60?}} &
200--480袋\\
\hline
\end{tabularx}

\medskip
\noindent\textbf{1.3 一级生物反应器接种}

\begin{tabularx}{\textwidth}{|p{0.52\textwidth}|X|}
\hline
\multicolumn{1}{|c|}{项目} & \multicolumn{1}{c|}{操作记录}\\
\hline
1）确认罐体参数达到如下范围且稳定30min以上：转速35rpm、表面通气20$\pm$5ml/min，底部通气10$\pm$3ml/min，温度37.0$\pm$2.0℃，溶氧80\%--110\%，pH 6.70--7.60； &
参数检查时间：
% #VALUE_ID: LEX-P0039-V0002
% #FIELD_VALUE: 参数检查时间
% #HANDWRITTEN: 08 : 53
\fieldvalue{\handwritten{08 : 53}}\\
& 转速：
% #VALUE_ID: LEX-P0039-V0003
% #FIELD_VALUE: 转速
% #HANDWRITTEN: 35
\fieldvalue{\handwritten{35}} rpm\\
& 温度：
% #VALUE_ID: LEX-P0039-V0004
% #FIELD_VALUE: 温度
% #HANDWRITTEN: 36.9
\fieldvalue{\handwritten{36.9}} ℃\\
& pH：
% #VALUE_ID: LEX-P0039-V0005
% #FIELD_VALUE: pH
% #HANDWRITTEN: 7.26
\fieldvalue{\handwritten{7.26}}\\
& 表观流量：
% #VALUE_ID: LEX-P0039-V0006
% #FIELD_VALUE: 表观流量
% #HANDWRITTEN: 20
\fieldvalue{\handwritten{20}} ml/min\\
& 底流流量：
% #VALUE_ID: LEX-P0039-V0007
% #FIELD_VALUE: 底流流量
% #HANDWRITTEN: 10
\fieldvalue{\handwritten{10}} ml/min\\
& DO：
% #VALUE_ID: LEX-P0039-V0008
% #FIELD_VALUE: DO
% #HANDWRITTEN: 93.5
\fieldvalue{\handwritten{93.5}}\%\\
& 是否符合要求：是$\square$ 否$\square$\\[-2pt]
& 
% #VALUE_ID: LEX-P0039-V0009
% #FIELD_VALUE: 是否符合要求
% #HANDWRITTEN: 是（勾选）
\fieldvalue{\handwritten{是（勾选）}}\\
& 稳定时长是否$>$30min：是$\square$ 否$\square$\\[-2pt]
&
% #VALUE_ID: LEX-P0039-V0010
% #FIELD_VALUE: 稳定时长是否大于30min
% #HANDWRITTEN: 是（勾选）
\fieldvalue{\handwritten{是（勾选）}}\\
\hline
2）接收一级生物反应器细胞接种液，确认接种液活细胞总数在（4.0--9.2）$\times10^7$范围内： &
接收时间：
% #VALUE_ID: LEX-P0039-V0011
% #FIELD_VALUE: 接收时间
% #HANDWRITTEN: 08 : 52
\fieldvalue{\handwritten{08 : 52}}\\
& 细胞接种液体积：
% #VALUE_ID: LEX-P0039-V0012
% #FIELD_VALUE: 细胞接种液体积
% #HANDWRITTEN: 500mL
\fieldvalue{\handwritten{500mL}}\\
& 活细胞总数：
% #VALUE_ID: LEX-P0039-V0013
% #FIELD_VALUE: 活细胞总数
% #HANDWRITTEN: 5.87×10^7
\fieldvalue{\handwritten{$5.87\times10^7$}} cells\\
& 细胞活率：
% #VALUE_ID: LEX-P0039-V0014
% #FIELD_VALUE: 细胞活率
% #HANDWRITTEN: 92.7
\fieldvalue{\handwritten{92.7}}\%\\
& 是否符合要求：是$\square$ 否$\square$\\[-2pt]
&
% #VALUE_ID: LEX-P0039-V0015
% #FIELD_VALUE: 是否符合要求
% #HANDWRITTEN: 是（勾选）
\fieldvalue{\handwritten{是（勾选）}}\\
\hline
3）使用接种机将接种液缓慢与反应器连接，泵入接种液，待接种液释放导入接种液混合后开始投入反应器内，初始培养体积在2.0$\pm$0.3L范围，接种完成； &
管路连接是否正确：是$\square$ 否$\square$\\[-2pt]
&
% #VALUE_ID: LEX-P0039-V0016
% #FIELD_VALUE: 管路连接是否正确
% #HANDWRITTEN: 是（勾选）
\fieldvalue{\handwritten{是（勾选）}}\\
& 初始培养体积：
% #VALUE_ID: LEX-P0039-V0017
% #FIELD_VALUE: 初始培养体积
% #HANDWRITTEN: 2.1L
\fieldvalue{\handwritten{2.1L}}\\
& 接种完成时间：
% #VALUE_ID: LEX-P0039-V0018
% #FIELD_VALUE: 接种完成时间
% #HANDWRITTEN: 08 : 57
\fieldvalue{\handwritten{08 : 57}}\\
\hline
4）设置溶氧值为50\%，并开启控制； &
开启溶氧控制时间：
% #VALUE_ID: LEX-P0039-V0019
% #FIELD_VALUE: 开启溶氧控制时间
% #HANDWRITTEN: 08 : 57
\fieldvalue{\handwritten{08 : 57}}\\
\hline
5）接种完成5min后，取样2袋，20ml/袋，检测液体pH、葡萄糖浓度、细胞密度、细胞活率、细胞生长状态，记录检测结果；\\
注：样品需与细胞密度和活率检测分装检测，且pH为取样管开启后言要检测项目；\\
当离心pH与在线pH差值$>$0.2时，应对在线电极单点校准。 &
取样时间：
% #VALUE_ID: LEX-P0039-V0020
% #FIELD_VALUE: 取样时间
% #HANDWRITTEN: 09 : 02
\fieldvalue{\handwritten{09 : 02}}\\
& 细胞密度：
% #VALUE_ID: LEX-P0039-V0021
% #FIELD_VALUE: 细胞密度
% #TODO #HANDWRITTEN: 9.44×10^?; exponent unclear
\fieldvalue{\handwritten{$9.44\times10^{?}$}} 个细胞/ml\\
& 活细胞率：
% #VALUE_ID: LEX-P0039-V0022
% #FIELD_VALUE: 活细胞率
% #HANDWRITTEN: 91
\fieldvalue{\handwritten{91}}\%\\
& 葡萄糖浓度：
% #VALUE_ID: LEX-P0039-V0023
% #FIELD_VALUE: 葡萄糖浓度
% #HANDWRITTEN: 22.7
\fieldvalue{\handwritten{22.7}} mmol/L\\
& 离线 pH：
% #VALUE_ID: LEX-P0039-V0024
% #FIELD_VALUE: 离线pH
% #HANDWRITTEN: 7.52
\fieldvalue{\handwritten{7.52}}\\
& 样品是否污染：是$\square$ 否$\square$\\[-2pt]
&
% #VALUE_ID: LEX-P0039-V0025
% #FIELD_VALUE: 样品是否污染
% #HANDWRITTEN: 是（勾选）
\fieldvalue{\handwritten{是（勾选）}}\\
& 镜检正常：是$\square$ 否$\square$\\[-2pt]
&
% #VALUE_ID: LEX-P0039-V0026
% #FIELD_VALUE: 镜检正常
% #HANDWRITTEN: 是（勾选）
\fieldvalue{\handwritten{是（勾选）}}\\
&
% #VALUE_ID: LEX-P0039-V0027
% #TODO #HANDWRITTEN: 疑似签名，无法辨认; standalone handwriting below the microscopy result
\handwritten{[签名无法辨认]}\\
\hline
\end{tabularx}

\begin{center}
第 21 页 共 222 页
\end{center}
% LEXOID_PAGE_COMPLETED: 39/59

\newpage

\begin{center}
\small
\begin{tabular}{@{}p{0.16\textwidth}p{0.25\textwidth}p{0.34\textwidth}p{0.18\textwidth}@{}}
\multicolumn{2}{l}{华兰生物基因工程有限公司} &
\multicolumn{2}{r}{记录编号：YZ-BPR-MF-PD-01-020-a} \\
\multicolumn{2}{l}{PLATYNUM LIFE SCIENCE BIOTECH} &
\multicolumn{2}{r}{%
% #VALUE_ID: LEX-P0040-V0001
% #FIELD_VALUE: 文件类型
\fieldvalue{原件}}\\
\end{tabular}

\vspace{2pt}

{\large\textbf{艾米迈托赛注射液（Z2）批生产记录}}

\vspace{4pt}

\renewcommand{\arraystretch}{1.35}
\begin{tabularx}{\textwidth}{|>{\centering\arraybackslash}p{0.18\textwidth}|X|>{\centering\arraybackslash}p{0.18\textwidth}|>{\centering\arraybackslash}p{0.17\textwidth}|}
\hline
品名 & 规格 & 生产批号 & 批量 \\
\hline
% #VALUE_ID: LEX-P0040-V0002
% #FIELD_VALUE: 品名
\fieldvalue{艾米迈托赛注射液}
&
% #VALUE_ID: LEX-P0040-V0003
% #FIELD_VALUE: 规格
\fieldvalue{12ml：$6.0\times10^{7}$个细胞}
&
% #VALUE_ID: LEX-P0040-V0004
% #FIELD_VALUE: 生产批号
% #HANDWRITTEN: 2025\ldots
\fieldvalue{\handwritten{2025\ldots}}
&
% #VALUE_ID: LEX-P0040-V0005
% #FIELD_VALUE: 批量
\fieldvalue{200--480袋}
\\
\hline
\end{tabularx}

\vspace{4pt}

\begin{tabularx}{\textwidth}{|p{0.49\textwidth}|X|}
\hline
\multicolumn{2}{l}{6）开启贴壁培养程序：} \\
\hline
\begin{minipage}[t]{\linewidth}
步骤1 分转速35rpm$\times$5min，步骤2 为转速0rpm$\times$25min，

步骤1 和步骤2 交替循环44次，间速完成后恒定转速

35rpm；
\end{minipage}
&
\begin{minipage}[t]{\linewidth}
贴壁培养开始时间：
% #VALUE_ID: LEX-P0040-V0006
% #FIELD_VALUE: 贴壁培养开始时间
% #HANDWRITTEN: 9:15
\fieldvalue{\handwritten{9:15}}

\vspace{8pt}

程序是否符合要求：是\(\Box\) 否\(\Box\)
\end{minipage}
\\
\hline
备注：
% #VALUE_ID: LEX-P0040-V0007
% #FIELD_VALUE: 备注
% #HANDWRITTEN: /
\fieldvalue{\handwritten{/}}
&
\\[18pt]
\hline
\end{tabularx}

\vspace{8pt}

\noindent
确认人/日期：
% #VALUE_ID: LEX-P0040-V0008
% #FIELD_VALUE: 确认人
% #HANDWRITTEN: 齐丽娟
\fieldvalue{\handwritten{齐丽娟}}
\quad
% #VALUE_ID: LEX-P0040-V0009
% #FIELD_VALUE: 确认日期
% #HANDWRITTEN: 2025.06.09
\fieldvalue{\handwritten{2025.06.09}}
\hfill
复核人/日期：
% #VALUE_ID: LEX-P0040-V0010
% #FIELD_VALUE: 复核人
% #HANDWRITTEN: 蒋丽娟
\fieldvalue{\handwritten{蒋丽娟}}
\quad
% #VALUE_ID: LEX-P0040-V0011
% #FIELD_VALUE: 复核日期
% #HANDWRITTEN: 2025.06.09
\fieldvalue{\handwritten{2025.06.09}}

\vfill

\begin{center}
第 22 页 共 222 页
\end{center}
\end{center}
% LEXOID_PAGE_COMPLETED: 40/59

\newpage

\begin{center}
\textbf{艾米奇托赛注射液（Z2）批生产记录}
\end{center}

\begin{tabularx}{\textwidth}{|>{\centering\arraybackslash}X|>{\centering\arraybackslash}X|>{\centering\arraybackslash}X|>{\centering\arraybackslash}X|}
\hline
品名 & 规格 & 生产批号 & 批量 \\
\hline
艾米奇托赛注射液 & 12ml；6.0$\times$10$^{7}$个细胞 &
% #VALUE_ID: LEX-P0041-V0001
% #FIELD_VALUE: 生产批号
% #TODO #HANDWRITTEN: R230?070; 手写批号部分字符难以辨认
\fieldvalue{\handwritten{R230?070}} &
200--480 瓶 \\
\hline
\end{tabularx}

\medskip

\textbf{1.4 清场}

操作依据：

\begin{enumerate}
\item 《生产区清场管理规程》\hfill YZ-SMP-PD-03-001
\item 《生产、检验状态标识管理规程》\hfill YZ-SMP-QA-04-014
\end{enumerate}

生产工序：一级生物反应器细胞培养

清场开始时间：
% #VALUE_ID: LEX-P0041-V0002
% #FIELD_VALUE: 清场开始时间（时）
% #HANDWRITTEN: 9
\fieldvalue{\handwritten{9}} 时
% #VALUE_ID: LEX-P0041-V0003
% #FIELD_VALUE: 清场开始时间（分）
% #HANDWRITTEN: 1
\fieldvalue{\handwritten{1}} 分

清场房间编码：
% #VALUE_ID: LEX-P0041-V0004
% #FIELD_VALUE: 清场房间编码
% #TODO #HANDWRITTEN: S406; 手写房间编码中间字符不清晰
\fieldvalue{\handwritten{S406}}
\hfill
75\%酒精批号：
% #VALUE_ID: LEX-P0041-V0005
% #FIELD_VALUE: 75\%酒精批号
% #TODO #HANDWRITTEN: 023?0?01; 手写批号字符难以辨认
\fieldvalue{\handwritten{023?0?01}}

\medskip

\begin{tabularx}{\textwidth}{|X|>{\centering\arraybackslash}p{0.24\textwidth}|}
\hline
\multicolumn{1}{|c|}{清场项目及标准} &
\multicolumn{1}{c|}{执行情况} \\
\hline
将本班次使用的各类工具配件恢复原位， &
% #VALUE_ID: LEX-P0041-V0006
% #FIELD_VALUE: 清场项目一执行情况
% #HANDWRITTEN: 已执行（勾选）
\fieldvalue{\handwritten{☑}}已执行 \quad $\square$否 \\
\hline
将本班次生产有关的其它物品或物料清出生产现场； &
% #VALUE_ID: LEX-P0041-V0007
% #FIELD_VALUE: 清场项目二执行情况
% #HANDWRITTEN: 已执行（勾选）
\fieldvalue{\handwritten{☑}}已执行 \quad $\square$否 \\
\hline
使用75\%酒精清洁设备、操作台面、不锈钢篮子、转移车、传递窗、门把手等； &
% #VALUE_ID: LEX-P0041-V0008
% #FIELD_VALUE: 清场项目三执行情况
% #HANDWRITTEN: 已执行（勾选）
\fieldvalue{\handwritten{☑}}已执行 \quad $\square$否 \\
\hline
使用75\%酒精喷洒墙面，然后使用拖把擦拭地面； &
% #VALUE_ID: LEX-P0041-V0009
% #FIELD_VALUE: 清场项目四执行情况
% #HANDWRITTEN: 已执行（勾选）
\fieldvalue{\handwritten{☑}}已执行 \quad $\square$否 \\
\hline
按照《生产、检验状态标识管理规程》更换房间、设备状态标识； &
% #VALUE_ID: LEX-P0041-V0010
% #FIELD_VALUE: 清场项目五执行情况
% #HANDWRITTEN: 已执行（勾选）
\fieldvalue{\handwritten{☑}}已执行 \quad $\square$否 \\
\hline
将本班次操作产生的废弃物移出洁净生产区； &
% #TODO: 表格下方内容在本页可见，按原页面顺序保留。
\end{tabularx}

\medskip

清场结束时间：
% #VALUE_ID: LEX-P0041-V0011
% #FIELD_VALUE: 清场结束时间（时）
% #HANDWRITTEN: 09
\fieldvalue{\handwritten{09}} 时
% #VALUE_ID: LEX-P0041-V0012
% #FIELD_VALUE: 清场结束时间（分）
% #HANDWRITTEN: 41
\fieldvalue{\handwritten{41}} 分

清场人员/日期：
% #VALUE_ID: LEX-P0041-V0013
% #FIELD_VALUE: 清场人员
% #TODO #HANDWRITTEN: 袁超; 姓名手写字迹不完全清晰
\fieldvalue{\handwritten{袁超}}
\quad
% #VALUE_ID: LEX-P0041-V0014
% #FIELD_VALUE: 清场日期
% #HANDWRITTEN: 2022.06.09
\fieldvalue{\handwritten{2022.06.09}}

\medskip

\begin{tabularx}{\textwidth}{|X|}
\hline
清场检查 \\
\hline
% #VALUE_ID: LEX-P0041-V0015
% #FIELD_VALUE: 清场检查结果
% #HANDWRITTEN: 合格（勾选）
\fieldvalue{\handwritten{☑}}合格 \\
\hline
$\square$不合格，具体内容为： \\
\hline
整改情况及整改人： \\
\hline
\end{tabularx}

确认人/日期：
% #VALUE_ID: LEX-P0041-V0016
% #FIELD_VALUE: 确认人
% #TODO #HANDWRITTEN: 方莉; 手写姓名部分字符难以辨认
\fieldvalue{\handwritten{方莉}}
\quad
% #VALUE_ID: LEX-P0041-V0017
% #FIELD_VALUE: 确认日期
% #HANDWRITTEN: 2022.06.09
\fieldvalue{\handwritten{2022.06.09}}

\vfill

\begin{center}
第23页 共222页
\end{center}
% LEXOID_PAGE_COMPLETED: 41/59

\newpage

\section*{艾米迈托赛注射液（Z2）批生产记录}

记录编号：YZ-BPR-MF-PD-01-020-a

\begin{tabularx}{\textwidth}{|c|X|c|X|c|X|}
\hline
品名 & 艾米迈托赛注射液 & 规格 & 12ml：$6.0\times10^{7}$个细胞 & 生产批号 &
% #VALUE_ID: LEX-P0042-V0001
% #FIELD_VALUE: 生产批号
% #TODO #HANDWRITTEN: A3...; 手写批号辨识不清
\fieldvalue{\handwritten{A3...}} \\
\hline
\multicolumn{4}{|c|}{} & 批量 & 200--480袋 \\
\hline
\end{tabularx}

\subsection*{2.\ 一级生物反应器组培培养第一天}

\subsection*{2.1\ 生产前确认}

\begin{tabularx}{\textwidth}{|p{0.54\textwidth}|X|}
\hline
\centering 项目 & \centering 操作记录 \tabularnewline
\hline
操作人员按《人员进出洁净生产区管理规程》要求进入洁净生产区作业。 &
操作间名称及编号：
% #VALUE_ID: LEX-P0042-V0002
% #FIELD_VALUE: 操作间名称及编号
% TODO #HANDWRITTEN: Z... SOBC; 手写内容辨识不清
\fieldvalue{\handwritten{Z... SOBC}} \\
\hline
确认生产区操作间在清场有效期内（C级日清洁效期3天）； &
是否在清洁有效期内：\\
% #VALUE_ID: LEX-P0042-V0003
% #FIELD_VALUE: 是否在清洁有效期内
\fieldvalue{☑ 是}\quad $\square$ 否\\
清洁效期：
% #VALUE_ID: LEX-P0042-V0004
% #FIELD_VALUE: 清洁效期
% #HANDWRITTEN: 2025.06.12
\fieldvalue{\handwritten{2025.06.12}} \\
\hline
确认操作间温度18--26℃，湿度30\%--75\%，相对压差$>0$Pa（操作间对C级走廊）； &
温度：
% #VALUE_ID: LEX-P0042-V0005
% #FIELD_VALUE: 温度
% #HANDWRITTEN: 20.9
\fieldvalue{\handwritten{20.9}}℃\quad
湿度：
% #VALUE_ID: LEX-P0042-V0006
% #FIELD_VALUE: 湿度
% #HANDWRITTEN: 49\%
\fieldvalue{\handwritten{49\%}}\\
压差：
% #VALUE_ID: LEX-P0042-V0007
% #FIELD_VALUE: 压差
% #HANDWRITTEN: 16
\fieldvalue{\handwritten{16}}Pa\\
是否符合要求：
% #VALUE_ID: LEX-P0042-V0008
% #FIELD_VALUE: 是否符合要求
\fieldvalue{☑ 是}\quad $\square$ 否 \\
\hline
确认操作间内无上批产品遗留物、文件及与本产品生产无关物料； &
% #VALUE_ID: LEX-P0042-V0009
% #FIELD_VALUE: 操作间内无遗留物料
\fieldvalue{☑ 是}\quad $\square$ 否 \\
\hline
检查发放的批记录及其它相关配套文件均与生产指令相对应，且均受控； &
% #VALUE_ID: LEX-P0042-V0010
% #FIELD_VALUE: 批记录及配套文件受控
\fieldvalue{☑ 是}\quad $\square$ 否 \\
\hline
确认设备完好，且在验证有效期内，计量器具有“合格证”且在有效期内，详见《主要设备仪/表确认表》； &
% #VALUE_ID: LEX-P0042-V0011
% #FIELD_VALUE: 设备及计量器具状态
\fieldvalue{☑ 是}\quad $\square$ 否 \\
\hline
反应器气源压力控制在如下范围：压缩空气进气压力0.05--0.15MPa、二氧化碳进气压力0.05--0.15MPa、氧气进气压力0.05--0.15MPa； &
压缩空气：
% #VALUE_ID: LEX-P0042-V0012
% #FIELD_VALUE: 压缩空气压力
% #HANDWRITTEN: 0.2
\fieldvalue{\handwritten{0.2}}MPa\newline
二氧化碳：
% #VALUE_ID: LEX-P0042-V0013
% #FIELD_VALUE: 二氧化碳压力
% #HANDWRITTEN: 0.2
\fieldvalue{\handwritten{0.2}}MPa\quad
氧气：
% #VALUE_ID: LEX-P0042-V0014
% #FIELD_VALUE: 氧气压力
% #HANDWRITTEN: 0.10
\fieldvalue{\handwritten{0.10}}MPa \\
\hline
将操作间状态标识更换为“生产加工状态卡”，写明品名、规格、批号等，并确认与本批生产指令一致； &
% #VALUE_ID: LEX-P0042-V0015
% #FIELD_VALUE: 生产加工状态卡
\fieldvalue{☑ 是}\quad $\square$ 否 \\
\hline
核对本次生产用物料耗材的品名、规格、批号、数量，检查物料外包装的完整性以及是否存在有效期内，并记录相关信息，详见《物料确认表》； &
% #VALUE_ID: LEX-P0042-V0016
% #FIELD_VALUE: 物料耗材核对
\fieldvalue{☑ 是}\quad $\square$ 否 \\
\hline
备注： &
% #VALUE_ID: LEX-P0042-V0017
% #HANDWRITTEN: /
\handwritten{/} \\
\hline
\end{tabularx}

\vspace{0.4em}

确认人/日期：
% #VALUE_ID: LEX-P0042-V0018
% #FIELD_VALUE: 确认人
% #HANDWRITTEN: 江峰
\fieldvalue{\handwritten{江峰}}
\quad
% #VALUE_ID: LEX-P0042-V0019
% #FIELD_VALUE: 确认日期
% #HANDWRITTEN: 2025.06.10
\fieldvalue{\handwritten{2025.06.10}}
\qquad
复核人/日期：
% #VALUE_ID: LEX-P0042-V0020
% #FIELD_VALUE: 复核人
% #HANDWRITTEN: 孙群
\fieldvalue{\handwritten{孙群}}
\quad
% #VALUE_ID: LEX-P0042-V0021
% #FIELD_VALUE: 复核日期
% #HANDWRITTEN: 2025.06.10
\fieldvalue{\handwritten{2025.06.10}}

\begin{center}
第24页 共222页
\end{center}
% LEXOID_PAGE_COMPLETED: 42/59

\newpage

\begin{center}
\textbf{艾米洛托莫注射液（Z2）批生产记录}
\end{center}

\noindent
\begin{tabularx}{\textwidth}{|l|X|l|X|}
\hline
原件 & 记录编号 &
% #VALUE_ID: LEX-P0043-V0001
% #FIELD_VALUE: 文件状态
\fieldvalue{原件} &
% #VALUE_ID: LEX-P0043-V0002
% #FIELD_VALUE: 记录编号
\fieldvalue{YZ-BPR-MF-PD-01-020-a} \\
\hline
\end{tabularx}

\noindent
\begin{tabularx}{\textwidth}{|l|X|l|X|l|X|}
\hline
品名 &
% #VALUE_ID: LEX-P0043-V0003
% #FIELD_VALUE: 品名
\fieldvalue{艾米洛托莫注射液} &
规格 &
% #VALUE_ID: LEX-P0043-V0004
% #FIELD_VALUE: 规格
\fieldvalue{12ml；$6.0\times10^{7}$个细胞} &
生产批号 &
% #VALUE_ID: LEX-P0043-V0005
% #FIELD_VALUE: 生产批号
% #TODO #HANDWRITTEN: A?23?29?6?41; handwriting partially illegible
\fieldvalue{\handwritten{A?23?29?6?41}} \\
\hline
\multicolumn{5}{|l|}{装量：%
% #VALUE_ID: LEX-P0043-V0006
% #FIELD_VALUE: 装量
\fieldvalue{200--480袋}} \\
\hline
\end{tabularx}

\subsection*{2.1.1 主要设备/仪表确认表}

\small
\begin{tabularx}{\textwidth}{|c|X|X|X|X|X|}
\hline
序号 & 设备名称 & 品牌/型号 & 设备编号 & 检查项 & 确认结果 \\
\hline
1 &
% #VALUE_ID: LEX-P0043-V0007
% #FIELD_VALUE: 设备名称
\fieldvalue{生物反应器控制台（主机）} &
% #VALUE_ID: LEX-P0043-V0008
% #FIELD_VALUE: 品牌/型号
\fieldvalue{华鑫生物} &
% #VALUE_ID: LEX-P0043-V0009
% #FIELD_VALUE: 设备编号
% #HANDWRITTEN: 1123610
\fieldvalue{\handwritten{1123610}} &
\multirow{5}{*}{1.设备应有“完好”标识；\newline
2.设备在验证有效期内；\newline
3.计量器具应有合格证，并且在有效期内。} &
% #VALUE_ID: LEX-P0043-V0010
% #FIELD_VALUE: 确认结果（是）
\fieldvalue{\ensuremath{\boxtimes}\,是}\quad$\square$ 否 \\
\hline
2 &
% #VALUE_ID: LEX-P0043-V0011
% #FIELD_VALUE: 设备名称
\fieldvalue{生物反应器罐体} &
% #VALUE_ID: LEX-P0043-V0012
% #FIELD_VALUE: 品牌/型号
\fieldvalue{华鑫生物} &
% #VALUE_ID: LEX-P0043-V0013
% #FIELD_VALUE: 设备编号
% #TODO #HANDWRITTEN: V?510?-202?; handwriting partially illegible
\fieldvalue{\handwritten{V?510?-202?}} &
&
% #VALUE_ID: LEX-P0043-V0014
% #FIELD_VALUE: 确认结果（是）
\fieldvalue{\ensuremath{\boxtimes}\,是}\quad$\square$ 否 \\
\hline
3 &
% #VALUE_ID: LEX-P0043-V0015
% #FIELD_VALUE: 设备名称
\fieldvalue{接管机} &
% #VALUE_ID: LEX-P0043-V0016
% #FIELD_VALUE: 品牌/型号
\fieldvalue{上海英启} &
% #VALUE_ID: LEX-P0043-V0017
% #FIELD_VALUE: 设备编号
% #HANDWRITTEN: 113610
\fieldvalue{\handwritten{113610}} &
&
% #VALUE_ID: LEX-P0043-V0018
% #FIELD_VALUE: 确认结果（是）
\fieldvalue{\ensuremath{\boxtimes}\,是}\quad$\square$ 否 \\
\hline
4 &
% #VALUE_ID: LEX-P0043-V0019
% #FIELD_VALUE: 设备名称
\fieldvalue{封管机} &
% #VALUE_ID: LEX-P0043-V0020
% #FIELD_VALUE: 品牌/型号
\fieldvalue{上海英启} &
% #VALUE_ID: LEX-P0043-V0021
% #FIELD_VALUE: 设备编号
% #HANDWRITTEN: 113261
\fieldvalue{\handwritten{113261}} &
&
% #VALUE_ID: LEX-P0043-V0022
% #FIELD_VALUE: 确认结果（是）
\fieldvalue{\ensuremath{\boxtimes}\,是}\quad$\square$ 否 \\
\hline
5 & & &
% #VALUE_ID: LEX-P0043-V0023
% #FIELD_VALUE: 设备编号
% #TODO #HANDWRITTEN: 乙级 2025-06-10; handwriting location and first characters uncertain
\fieldvalue{\handwritten{乙级 2025-06-10}} &
& $\square$ 是\quad$\square$ 否 \\
\hline
\multicolumn{6}{|l|}{备注：%
% #VALUE_ID: LEX-P0043-V0024
% #HANDWRITTEN: ／
% #TODO #HANDWRITTEN: ／; brief handwritten mark
\handwritten{／}} \\
\hline
\end{tabularx}
\normalsize

\noindent
确认人/日期：
% #VALUE_ID: LEX-P0043-V0025
% #FIELD_VALUE: 确认人
% #TODO #HANDWRITTEN: 汪?峰; name partially illegible
\fieldvalue{\handwritten{汪?峰}}
\quad
% #VALUE_ID: LEX-P0043-V0026
% #FIELD_VALUE: 确认日期
% #HANDWRITTEN: 2025.06.10
\fieldvalue{\handwritten{2025.06.10}}
\qquad
复核人/日期：
% #VALUE_ID: LEX-P0043-V0027
% #FIELD_VALUE: 复核人
% #TODO #HANDWRITTEN: 李?军; name partially illegible
\fieldvalue{\handwritten{李?军}}
\quad
% #VALUE_ID: LEX-P0043-V0028
% #FIELD_VALUE: 复核日期
% #HANDWRITTEN: 2025.06.10
\fieldvalue{\handwritten{2025.06.10}}

\subsection*{2.1.2 物料确认表}

\small
\begin{tabularx}{\textwidth}{|c|X|X|X|X|c|X|}
\hline
序号 & 名称 & 品牌 & 规格 & 批号 & 数量 & 是否符合要求 \\
\hline
1 &
% #VALUE_ID: LEX-P0043-V0029
% #FIELD_VALUE: 名称
% #HANDWRITTEN: 50ml-一次性储液袋
\fieldvalue{\handwritten{50ml-一次性储液袋}} &
% #VALUE_ID: LEX-P0043-V0030
% #FIELD_VALUE: 品牌
% #HANDWRITTEN: 金坛瑞
\fieldvalue{\handwritten{金坛瑞}} &
% #VALUE_ID: LEX-P0043-V0031
% #FIELD_VALUE: 规格
% #TODO #HANDWRITTEN: 1/1?; handwriting partially illegible
\fieldvalue{\handwritten{1/1?}} &
% #VALUE_ID: LEX-P0043-V0032
% #FIELD_VALUE: 批号
% #HANDWRITTEN: R230202023022
\fieldvalue{\handwritten{R230202023022}} &
% #VALUE_ID: LEX-P0043-V0033
% #FIELD_VALUE: 数量
% #HANDWRITTEN: 3
\fieldvalue{\handwritten{3}} &
% #VALUE_ID: LEX-P0043-V0034
% #FIELD_VALUE: 是否符合要求（是）
\fieldvalue{\ensuremath{\boxtimes}\,是}\quad$\square$ 否 \\
\hline
2 & & & &
% #VALUE_ID: LEX-P0043-V0035
% #FIELD_VALUE: 批号
% #HANDWRITTEN: 2025.06.10
\fieldvalue{\handwritten{2025.06.10}} &
& $\square$ 是\quad$\square$ 否 \\
\hline
\multicolumn{7}{|l|}{备注：} \\
\hline
\end{tabularx}
\normalsize

\noindent
确认人/日期：
% #VALUE_ID: LEX-P0043-V0036
% #FIELD_VALUE: 确认人
% #TODO #HANDWRITTEN: 汪?峰; name partially illegible
\fieldvalue{\handwritten{汪?峰}}
\quad
% #VALUE_ID: LEX-P0043-V0037
% #FIELD_VALUE: 确认日期
% #HANDWRITTEN: 2025.06.10
\fieldvalue{\handwritten{2025.06.10}}
\qquad
复核人/日期：
% #VALUE_ID: LEX-P0043-V0038
% #FIELD_VALUE: 复核人
% #TODO #HANDWRITTEN: 李?军; name partially illegible
\fieldvalue{\handwritten{李?军}}
\quad
% #VALUE_ID: LEX-P0043-V0039
% #FIELD_VALUE: 复核日期
% #HANDWRITTEN: 2025.06.10
\fieldvalue{\handwritten{2025.06.10}}

\subsection*{2.2 操作内容}

\begin{tabularx}{\textwidth}{|X|X|}
\hline
项目 & 操作记录 \\
\hline
\end{tabularx}

% TODO: The operation-content table continues beyond the visible portion of this page.
% LEXOID_PAGE_COMPLETED: 43/59

\newpage

\begin{center}
{\Large \textbf{艾米迈托赛注射液（Z2）批生产记录}}
\end{center}

\noindent
\begin{tabularx}{\textwidth}{|X|X|X|X|}
\hline
品名 & 规格 & 生产批号 & 批量 \\
\hline
艾米迈托赛注射液 & 12ml：$6.0\times10^{7}$个细胞 &
% #VALUE_ID: LEX-P0044-V0001
% #FIELD_VALUE: 生产批号
% #TODO #HANDWRITTEN: A?3?2?0?; handwriting is unclear
\fieldvalue{\handwritten{A?3?2?0?}} &
200--400袋 \\
\hline
\end{tabularx}

\vspace{2mm}

\begin{tabularx}{\textwidth}{|p{0.48\textwidth}|X|}
\hline
\textbf{1）参数检查：} &
检查时间：
% #VALUE_ID: LEX-P0044-V0002
% #FIELD_VALUE: 检查时间
% #HANDWRITTEN: 09：14
\fieldvalue{\handwritten{09：14}}\\
\parbox[t]{0.46\textwidth}{
检查培养参数控制值是否在如下范围。表面通气20
$\pm$5ml/min，底部通气 $\leq$10$\pm$3$ml/min，转速
$35$--$50$rpm$，温度$37.0\pm2.0^\circ$C，溶氧
$>30\%$，pH$6.70$--$7.60$；
} &
搅拌转速：
% #VALUE_ID: LEX-P0044-V0003
% #FIELD_VALUE: 搅拌转速
% #HANDWRITTEN: 25
\fieldvalue{\handwritten{25}} rpm\\
&
温度：
% #VALUE_ID: LEX-P0044-V0004
% #FIELD_VALUE: 温度
% #HANDWRITTEN: 37.0
\fieldvalue{\handwritten{37.0}}$^\circ$C\\
&
DO：
% #VALUE_ID: LEX-P0044-V0005
% #FIELD_VALUE: DO
% #HANDWRITTEN: 80.0
\fieldvalue{\handwritten{80.0}}\%\\
&
pH：
% #VALUE_ID: LEX-P0044-V0006
% #FIELD_VALUE: pH
% #HANDWRITTEN: 7.40
\fieldvalue{\handwritten{7.40}}\\
&
表面流量：
% #VALUE_ID: LEX-P0044-V0007
% #FIELD_VALUE: 表面流量
% #HANDWRITTEN: 20
\fieldvalue{\handwritten{20}}ml/min\\
&
底部流量：
% #VALUE_ID: LEX-P0044-V0008
% #FIELD_VALUE: 底部流量
% #HANDWRITTEN: 20
\fieldvalue{\handwritten{20}}ml/min\\
&
培养体积：
% #VALUE_ID: LEX-P0044-V0009
% #FIELD_VALUE: 培养体积
% #HANDWRITTEN: 2.0
\fieldvalue{\handwritten{2.0}}L\\
&
是否符合要求：□是\quad □否；
% #VALUE_ID: LEX-P0044-V0010
% #FIELD_VALUE: 是否符合要求
\fieldvalue{☑是}\quad □否\\
\hline
\textbf{2）观察缸体内微载体分布情况，微载体上下分布不均或一段时间出现较多结团时，提高转速至40--50rpm，每次2--5rpm；} &
是否调整转速：□是\quad □否；
% #VALUE_ID: LEX-P0044-V0011
% #FIELD_VALUE: 是否调整转速
\fieldvalue{☑是}\quad □否\\
&
调整时间：\underline{\hspace{25mm}}\\
&
调整后转速：\underline{\hspace{20mm}} rpm\\
\hline
\textbf{3）取样检测：} &
取样时间：
% #VALUE_ID: LEX-P0044-V0012
% #FIELD_VALUE: 取样时间
% #HANDWRITTEN: 09：27
\fieldvalue{\handwritten{09：27}}\\
\parbox[t]{0.46\textwidth}{
取样两份，每份体积$\geq20$ml，分别检测细胞密度、
活率、pH、葡萄糖浓度和细胞生长状态（镜检）。

注：pH检测与细胞密度/活率检测分装检测，且pH为
取样开口后首次检测项目；\par
当离线pH与在线pH差值$>0.2$时，应对在线电极做单点校
准。
} &
培养时长：
% #VALUE_ID: LEX-P0044-V0013
% #FIELD_VALUE: 培养时长
% #HANDWRITTEN: 24.5
\fieldvalue{\handwritten{24.5}}h\\
&
细胞密度：
% #VALUE_ID: LEX-P0044-V0014
% #FIELD_VALUE: 细胞密度
% #TODO #HANDWRITTEN: 1.83?; handwritten digits are unclear
\fieldvalue{\handwritten{1.83?}}\ $\times10^{?}$个细胞/ml\\
&
细胞活率：
% #VALUE_ID: LEX-P0044-V0015
% #FIELD_VALUE: 细胞活率
% #HANDWRITTEN: 86
\fieldvalue{\handwritten{86}}\%\\
&
葡萄糖浓度：
% #VALUE_ID: LEX-P0044-V0016
% #FIELD_VALUE: 葡萄糖浓度
% #HANDWRITTEN: 1.6
\fieldvalue{\handwritten{1.6}} mmol/L\\
&
离线pH：
% #VALUE_ID: LEX-P0044-V0017
% #FIELD_VALUE: 离线pH
% #HANDWRITTEN: 7.30
\fieldvalue{\handwritten{7.30}}\\
&
电极是否校准：\\
&
□是\quad
% #VALUE_ID: LEX-P0044-V0018
% #FIELD_VALUE: 电极是否校准
\fieldvalue{☑否}\\
&
镜检正常：\\
&
% #VALUE_ID: LEX-P0044-V0019
% #FIELD_VALUE: 镜检正常
\fieldvalue{☑是}\quad □否\\
\hline
\end{tabularx}

\vspace{2mm}

备注：
% #VALUE_ID: LEX-P0044-V0020
% #FIELD_VALUE: 备注
% #HANDWRITTEN: /
\fieldvalue{\handwritten{/}}

\vspace{3mm}

操作人/日期：
% #VALUE_ID: LEX-P0044-V0021
% #FIELD_VALUE: 操作人
% #TODO #HANDWRITTEN: 刘?鹏; name is unclear
\fieldvalue{\handwritten{刘?鹏}}
\quad
% #VALUE_ID: LEX-P0044-V0022
% #FIELD_VALUE: 操作日期
% #HANDWRITTEN: 2025-06-10
\fieldvalue{\handwritten{2025-06-10}}
\hfill
复核人/日期：
% #VALUE_ID: LEX-P0044-V0023
% #FIELD_VALUE: 复核人
% #TODO #HANDWRITTEN: 李?婷; name is unclear
\fieldvalue{\handwritten{李?婷}}
\quad
% #VALUE_ID: LEX-P0044-V0024
% #FIELD_VALUE: 复核日期
% #HANDWRITTEN: 2025-06-10
\fieldvalue{\handwritten{2025-06-10}}

\vfill

\begin{center}
第 26 页 共 222 页
\end{center}
% LEXOID_PAGE_COMPLETED: 44/59

\newpage

\begin{center}
\textbf{艾米达托美注射液（Z2）批生产记录}
\end{center}

记录编号：YZ-BPR-MF-PD-01-020-a

\begin{tabularx}{\textwidth}{|c|X|c|X|c|X|}
\hline
品名 & 艾米达托美注射液 & 规格 & 12ml：$6.0\times10^7$个细胞 & 生产批号 &
% #VALUE_ID: LEX-P0045-V0001
% #FIELD_VALUE: 生产批号
% #TODO #HANDWRITTEN: A3?3?2?5?6?; 字迹不清
\fieldvalue{\handwritten{A3?3?2?5?6?}} \\
\hline
\multicolumn{4}{|r|}{批量} & \multicolumn{2}{c|}{200--480袋} \\
\hline
\end{tabularx}

\subsection*{2.3 清场}

操作依据：

\begin{enumerate}
\item 《生产区内清场管理规程》\hfill YZ-SMP-PD-03-001
\item 《生产、检验状态标识管理规程》\hfill YZ-SMP-QA-04-014
\end{enumerate}

生产工序：一级生物反应器细胞培养

清场开始时间：
% #VALUE_ID: LEX-P0045-V0002
% #FIELD_VALUE: 清场开始时间（时）
% #HANDWRITTEN: 09
\fieldvalue{\handwritten{09}}时
% #VALUE_ID: LEX-P0045-V0003
% #FIELD_VALUE: 清场开始时间（分）
% #HANDWRITTEN: 45
\fieldvalue{\handwritten{45}}分

清场房间编码：
% #VALUE_ID: LEX-P0045-V0004
% #FIELD_VALUE: 清场房间编码
% #TODO #HANDWRITTEN: S?boc; 字迹不清
\fieldvalue{\handwritten{S?boc}}

\hfill 75\%酒精料号：
% #VALUE_ID: LEX-P0045-V0005
% #FIELD_VALUE: 75\%酒精料号
% #HANDWRITTEN: 20250610
\fieldvalue{\handwritten{20250610}}

\begin{tabularx}{\textwidth}{|X|c|}
\hline
\multicolumn{1}{|c|}{清场项目及标准} & 执行情况 \\
\hline
将本班次使用的各类工具配件恢复原位放置；
&
% #VALUE_ID: LEX-P0045-V0006
% #FIELD_VALUE: 清场项目第1项执行情况
% #HANDWRITTEN: 已执行
\fieldvalue{\handwritten{已执行}}\quad $\square$否 \\
\hline
将本班次生产有关的其它物品或物料清出生产现场；
&
% #VALUE_ID: LEX-P0045-V0007
% #FIELD_VALUE: 清场项目第2项执行情况
% #HANDWRITTEN: 已执行
\fieldvalue{\handwritten{已执行}}\quad $\square$否 \\
\hline
使用75\%酒精清洁设备、操作台面、不锈钢架子、转移车、传递窗、门把手等；
&
% #VALUE_ID: LEX-P0045-V0008
% #FIELD_VALUE: 清场项目第3项执行情况
% #HANDWRITTEN: 已执行
\fieldvalue{\handwritten{已执行}}\quad $\square$否 \\
\hline
使用75\%酒精清洁墙面，然后使用拖把擦拭地面；
&
% #VALUE_ID: LEX-P0045-V0009
% #FIELD_VALUE: 清场项目第4项执行情况
% #HANDWRITTEN: 已执行
\fieldvalue{\handwritten{已执行}}\quad $\square$否 \\
\hline
按照《生产、检验状态标识管理规程》更换房间、设备状态标识；
&
% #VALUE_ID: LEX-P0045-V0010
% #FIELD_VALUE: 清场项目第5项执行情况
% #HANDWRITTEN: 已执行
\fieldvalue{\handwritten{已执行}}\quad $\square$否 \\
\hline
将本班次操作产生的废弃物移出洁净生产区；
& \\
\hline
\end{tabularx}

清场结束时间：
% #VALUE_ID: LEX-P0045-V0011
% #FIELD_VALUE: 清场结束时间（时）
% #HANDWRITTEN: 01
\fieldvalue{\handwritten{01}}时
% #VALUE_ID: LEX-P0045-V0012
% #FIELD_VALUE: 清场结束时间（分）
% #HANDWRITTEN: 50
\fieldvalue{\handwritten{50}}分

清场人/日期：
% #VALUE_ID: LEX-P0045-V0013
% #FIELD_VALUE: 清场人
% #TODO #HANDWRITTEN: 刘?晓; 字迹不清
\fieldvalue{\handwritten{刘?晓}}
\quad
% #VALUE_ID: LEX-P0045-V0014
% #FIELD_VALUE: 清场日期
% #HANDWRITTEN: 2025.06.10
\fieldvalue{\handwritten{2025.06.10}}

\begin{tabularx}{\textwidth}{|X|}
\hline
清场检查 \\
\hline
% #VALUE_ID: LEX-P0045-V0015
% #FIELD_VALUE: 清场检查结果
% #HANDWRITTEN: 合格
\fieldvalue{\handwritten{合格}} \\
\hline
$\square$不合格，具体内容为： \\
\hline
整改情况及整改人： \\
\hline
\end{tabularx}

确认人/日期：
% #VALUE_ID: LEX-P0045-V0016
% #FIELD_VALUE: 确认人
% #TODO #HANDWRITTEN: 不详; 签名难以辨认
\fieldvalue{\handwritten{不详}}
\quad
% #VALUE_ID: LEX-P0045-V0017
% #FIELD_VALUE: 确认日期
% #HANDWRITTEN: 2025.06.10
\fieldvalue{\handwritten{2025.06.10}}

\begin{center}
第 27 页 共 222 页
\end{center}
% LEXOID_PAGE_COMPLETED: 45/59

\newpage

\section*{艾米迈托赛注射液（Z2）批生产记录}

记录编号：YZ-BPR-MF-PD-01-020-a\hfill
版本号：
% #VALUE_ID: LEX-P0046-V0001
% #FIELD_VALUE: 版本号
\fieldvalue{1.3}

\begin{tabularx}{\textwidth}{|X|X|X|X|}
\hline
产品名 & 规格 & 生产批号 & 批量 \\
\hline
% #VALUE_ID: LEX-P0046-V0002
% #FIELD_VALUE: 产品名
\fieldvalue{艾米迈托赛注射液}
&
% #VALUE_ID: LEX-P0046-V0003
% #FIELD_VALUE: 规格
\fieldvalue{12ml：$6.0\times10^{7}$个细胞}
&
% #VALUE_ID: LEX-P0046-V0004
% #FIELD_VALUE: 生产批号
% #TODO #HANDWRITTEN: A333...; 手写批号部分字迹不清
\fieldvalue{\handwritten{A333...}}
&
% #VALUE_ID: LEX-P0046-V0005
% #FIELD_VALUE: 批量
\fieldvalue{200--480袋}
\\
\hline
\end{tabularx}

\subsection*{3.\ 一级生物反应器细胞培养第二天}

\subsubsection*{3.1 生产前确认}

\begin{tabularx}{\textwidth}{|p{0.56\textwidth}|X|}
\hline
\centering 项目 & \centering 操作记录 \arraybackslash \\
\hline
操作人员按《人员进出洁净生产区管理规程》要求进入洁净区生产车间：
&
操作员名称及编号：
% #VALUE_ID: LEX-P0046-V0006
% #FIELD_VALUE: 操作员名称及编号
% TODO #HANDWRITTEN: 手写姓名及编号难以辨认
\fieldvalue{\handwritten{手写内容}}
\\
\hline
确认生产区操作间在清场有效期内（C级日清洁有效期3天）；
&
是否在清洁有效期内：\\
% #VALUE_ID: LEX-P0046-V0007
% #FIELD_VALUE: 清洁有效期内
% #HANDWRITTEN: ✓
\fieldvalue{\handwritten{✓}}\ 是\quad $\square$ 否\\
清洁日期：
% #VALUE_ID: LEX-P0046-V0008
% #FIELD_VALUE: 清洁日期
% #HANDWRITTEN: 2025.06.13
\fieldvalue{\handwritten{2025.06.13}}
\\
\hline
确认操作间温度18--26℃、湿度30\%--75\%、相对压差$>0$Pa（操作间对C级走廊向）；
&
温度：
% #VALUE_ID: LEX-P0046-V0009
% #FIELD_VALUE: 温度
% #HANDWRITTEN: 19.9
\fieldvalue{\handwritten{19.9}}℃\quad
湿度：
% #VALUE_ID: LEX-P0046-V0010
% #FIELD_VALUE: 湿度
% #HANDWRITTEN: 51.9
\fieldvalue{\handwritten{51.9}}\%\\
&
压差：
% #VALUE_ID: LEX-P0046-V0011
% #FIELD_VALUE: 压差
% #TODO #HANDWRITTEN: 1b; 单位前数字字迹不清
\fieldvalue{\handwritten{1b}}Pa\\
\hline
确认操作间内无上批产品遗留物、文件及与本产品生产无关物料；
&
% #VALUE_ID: LEX-P0046-V0012
% #FIELD_VALUE: 操作间无遗留物料
% #HANDWRITTEN: ✓
\fieldvalue{\handwritten{✓}}\ 是\quad $\square$ 否
\\
\hline
检查发放的批记录及其它相关配套文件均与生产指令相对应，且均受控；
&
% #VALUE_ID: LEX-P0046-V0013
% #FIELD_VALUE: 批记录及配套文件受控
% #HANDWRITTEN: ✓
\fieldvalue{\handwritten{✓}}\ 是\quad $\square$ 否
\\
\hline
确认设备完好，且在验证有效期内，计量器具有“合格证”且在有效期内，详见《主要设备仪器清单》；
&
% #VALUE_ID: LEX-P0046-V0014
% #FIELD_VALUE: 设备及计量器具状态
% #HANDWRITTEN: ✓
\fieldvalue{\handwritten{✓}}\ 是\quad $\square$ 否
\\
\hline
反应器气源压力应控制在如下范围：压缩空气进气压力
0.05--0.15MPa、二氧化碳进气压力0.05--0.15MPa、氧气进气压力0.05--0.15MPa；
&
压缩空气：
% #VALUE_ID: LEX-P0046-V0015
% #FIELD_VALUE: 压缩空气压力
% #HANDWRITTEN: 0.2
\fieldvalue{\handwritten{0.2}}MPa\newline
二氧化碳：
% #VALUE_ID: LEX-P0046-V0016
% #FIELD_VALUE: 二氧化碳压力
% #HANDWRITTEN: 0.08
\fieldvalue{\handwritten{0.08}}MPa\quad
氧气：
% #VALUE_ID: LEX-P0046-V0017
% #FIELD_VALUE: 氧气压力
% #HANDWRITTEN: 0.0
\fieldvalue{\handwritten{0.0}}MPa
\\
\hline
将操作间状态标识更换为“生产加工状态卡”，写明品名、规格、批号等，并确认与本批生产指令一致；
&
% #VALUE_ID: LEX-P0046-V0018
% #FIELD_VALUE: 生产加工状态卡
% #HANDWRITTEN: ✓
\fieldvalue{\handwritten{✓}}\ 是\quad $\square$ 否
\\
\hline
核对本次生产物料耗材的品名、规格、批号、数量，检查物料外包装的完整性以及是否在有效期及领牌内，并记录相关信息，详见《物料确认表》；
&
% #VALUE_ID: LEX-P0046-V0019
% #FIELD_VALUE: 物料及耗材确认
% #HANDWRITTEN: ✓
\fieldvalue{\handwritten{✓}}\ 是\quad $\square$ 否
\\
\hline
备注：
&
% #VALUE_ID: LEX-P0046-V0020
% #FIELD_VALUE: 备注
% #HANDWRITTEN: /
\fieldvalue{\handwritten{/}}
\\
\hline
\end{tabularx}

\vspace{0.5em}

确认人/日期：
% #VALUE_ID: LEX-P0046-V0021
% #FIELD_VALUE: 确认人
% #TODO #HANDWRITTEN: 手写签名难以辨认
\fieldvalue{\handwritten{手写签名}}
\quad
% #VALUE_ID: LEX-P0046-V0022
% #FIELD_VALUE: 确认日期
% #HANDWRITTEN: 2025.06.11
\fieldvalue{\handwritten{2025.06.11}}

\hfill
复核人/日期：
% #VALUE_ID: LEX-P0046-V0023
% #FIELD_VALUE: 复核人
% #TODO #HANDWRITTEN: 手写签名有涂改且难以辨认
\fieldvalue{\handwritten{手写签名}}
\quad
% #VALUE_ID: LEX-P0046-V0024
% #FIELD_VALUE: 复核日期
% #HANDWRITTEN: 2025.06.11
\fieldvalue{\handwritten{2025.06.11}}

\vfill

\begin{center}
第28页 共222页
\end{center}
% LEXOID_PAGE_COMPLETED: 46/59

\newpage

\section*{艾米迈托赛注射液（Z2）批生产记录}

\begin{tabularx}{\textwidth}{|l|X|l|X|}
\hline
品名 & 艾米迈托赛注射液 & 规格 & $12\mathrm{mL}: 6.0\times10^{7}$个细胞 \\
\hline
生产批号 & \multicolumn{3}{l|}{%
% #VALUE_ID: LEX-P0047-V0001
% #FIELD_VALUE: 生产批号
% #TODO #HANDWRITTEN: A?3?D?0?; 字迹不清
\fieldvalue{\handwritten{A?3?D?0?}}} \\
\hline
\end{tabularx}

\subsection*{3.1.1 主要设备/仪表确认表}

\small
\begin{tabularx}{\textwidth}{|c|X|X|X|X|c|}
\hline
序号 & 设备名称 & 品牌/型号 & 设备编号 & 检查项 & 确认结果 \\
\hline
1 & 生物反应器控制台（主机） & 华鑫生物 &
% #VALUE_ID: LEX-P0047-V0002
% #FIELD_VALUE: 设备编号
% #TODO #HANDWRITTEN: 113?610; 字迹不清
\fieldvalue{\handwritten{113?610}} &
\begin{tabular}[t]{@{}l@{}}
1.设备是否有“完好”标识；\\
2.设备在验证有效期内；\\
3.计量器具有“合格”证且在有效期内。
\end{tabular} &
% #VALUE_ID: LEX-P0047-V0003
% #FIELD_VALUE: 确认结果
\fieldvalue{是（已选）} \\
\hline
2 & 生物反应器罐体 & 华鑫生物 &
% #VALUE_ID: LEX-P0047-V0004
% #FIELD_VALUE: 设备编号
% #TODO #HANDWRITTEN: VG?T?D?--?; 字迹不清
\fieldvalue{\handwritten{VG?T?D?--?}} &
 & 
% #VALUE_ID: LEX-P0047-V0005
% #FIELD_VALUE: 确认结果
\fieldvalue{是（已选）} \\
\hline
3 & 接管机 & 上海英启 &
% #VALUE_ID: LEX-P0047-V0006
% #FIELD_VALUE: 设备编号
% #TODO #HANDWRITTEN: 112?610; 字迹不清
\fieldvalue{\handwritten{112?610}} &
 &
% #VALUE_ID: LEX-P0047-V0007
% #FIELD_VALUE: 确认结果
\fieldvalue{是（已选）} \\
\hline
4 & 封管机 & 上海英启 &
% #VALUE_ID: LEX-P0047-V0008
% #FIELD_VALUE: 设备编号
% #TODO #HANDWRITTEN: 112?61; 字迹不清
\fieldvalue{\handwritten{112?61}} &
 &
% #VALUE_ID: LEX-P0047-V0009
% #FIELD_VALUE: 确认结果
\fieldvalue{是（已选）} \\
\hline
5 & & &
% #VALUE_ID: LEX-P0047-V0010
% #TODO #HANDWRITTEN: ?? 2025-06-11; 表格空白行中的手写内容难以辨认
\handwritten{?? 2025-06-11} &
 & \\
\hline
\multicolumn{6}{|l|}{备注：%
% #VALUE_ID: LEX-P0047-V0011
% #HANDWRITTEN: /
\handwritten{/}} \\
\hline
\end{tabularx}
\normalsize

确认人/日期：
% #VALUE_ID: LEX-P0047-V0012
% #HANDWRITTEN: 江??（签名）
\handwritten{江??}
\quad
% #VALUE_ID: LEX-P0047-V0013
% #HANDWRITTEN: 2025.06.11
\handwritten{2025.06.11}
\hfill
复核人/日期：
% #VALUE_ID: LEX-P0047-V0014
% #HANDWRITTEN: 石??（签名）
\handwritten{石??}
\quad
% #VALUE_ID: LEX-P0047-V0015
% #HANDWRITTEN: 2025.06.11
\handwritten{2025.06.11}

\subsection*{3.1.2 物料确认表}

\small
\begin{tabularx}{\textwidth}{|c|X|X|X|X|c|c|}
\hline
序号 & 名称 & 品牌 & 规格 & 批号 & 数量 & 是否符合要求 \\
\hline
1 &
% #VALUE_ID: LEX-P0047-V0016
% #FIELD_VALUE: 名称
% #TODO #HANDWRITTEN: 50mL-一次性储液袋; 字迹不清
\fieldvalue{\handwritten{50mL-一次性储液袋}} &
% #VALUE_ID: LEX-P0047-V0017
% #FIELD_VALUE: 品牌
% #TODO #HANDWRITTEN: 金瑞达; 字迹不清
\fieldvalue{\handwritten{金瑞达}} &
% #VALUE_ID: LEX-P0047-V0018
% #FIELD_VALUE: 规格
% #TODO #HANDWRITTEN: 1/??; 字迹不清
\fieldvalue{\handwritten{1/??}} &
% #VALUE_ID: LEX-P0047-V0019
% #FIELD_VALUE: 批号
% #TODO #HANDWRITTEN: R130?202503030?; 字迹不清
\fieldvalue{\handwritten{R130?202503030?}} &
% #VALUE_ID: LEX-P0047-V0020
% #FIELD_VALUE: 数量
\fieldvalue{\handwritten{3}} &
% #VALUE_ID: LEX-P0047-V0021
% #FIELD_VALUE: 是否符合要求
\fieldvalue{是（已选）} \\
\hline
2 & &
% #VALUE_ID: LEX-P0047-V0022
% #TODO #HANDWRITTEN: 李春??; 字迹不清
\handwritten{李春??} &
& 
% #VALUE_ID: LEX-P0047-V0023
% #HANDWRITTEN: 2025.06.11
\handwritten{2025.06.11} &
& \\
\hline
\multicolumn{7}{|l|}{备注：%
% #VALUE_ID: LEX-P0047-V0024
% #HANDWRITTEN: /
\handwritten{/}} \\
\hline
\end{tabularx}
\normalsize

确认人/日期：
% #VALUE_ID: LEX-P0047-V0025
% #HANDWRITTEN: 江??（签名）
\handwritten{江??}
\quad
% #VALUE_ID: LEX-P0047-V0026
% #HANDWRITTEN: 2025.06.11
\handwritten{2025.06.11}
\hfill
复核人/日期：
% #VALUE_ID: LEX-P0047-V0027
% #HANDWRITTEN: 李春??（签名）
\handwritten{李春??}
\quad
% #VALUE_ID: LEX-P0047-V0028
% #HANDWRITTEN: 2025.06.11
\handwritten{2025.06.11}

\subsection*{3.2 操作内容}

\begin{tabularx}{\textwidth}{|X|X|}
\hline
项目 & 操作记录 \\
\hline
\end{tabularx}
% LEXOID_PAGE_COMPLETED: 47/59

\newpage

\begin{center}
\textbf{艾米迈托美注射液（Z2）批生产记录}
\end{center}

\noindent
记录编号：YZ-BPR-MF-PD-01-020-a

\begin{tabularx}{\textwidth}{|c|X|c|c|}
\hline
品名 & 规格 & 生产批号 & 批量 \\
\hline
艾米迈托美注射液 & 12ml；6.0$\times 10^{7}$个细胞 & 
% #VALUE_ID: LEX-P0048-V0001
% #FIELD_VALUE: 生产批号
% #TODO #HANDWRITTEN: A333333?; handwriting unclear
\fieldvalue{\handwritten{A333333?}} &
200--480 袋 \\
\hline
\end{tabularx}

\vspace{2pt}

\begin{tabularx}{\textwidth}{|p{0.48\textwidth}|X|}
\hline
\begin{minipage}[t]{\linewidth}
1）参数检查：

检查培养参数控制值是否在如下范围：表面通气20$\pm$5 ml/min，底部通气10$\pm$3 ml/min，转速35--50 rpm，温度37.0$\pm$2.0℃，溶氧$\geq$30\%，pH 6.70--7.60。
\end{minipage}
&
\begin{minipage}[t]{\linewidth}
检查时间：
% #VALUE_ID: LEX-P0048-V0002
% #FIELD_VALUE: 检查时间
% #HANDWRITTEN: 09：18
\fieldvalue{\handwritten{09：18}}

\vspace{2pt}
搅拌转速：
% #VALUE_ID: LEX-P0048-V0003
% #FIELD_VALUE: 搅拌转速
% #HANDWRITTEN: 30 rpm
\fieldvalue{\handwritten{30 rpm}}

\vspace{2pt}
温度：
% #VALUE_ID: LEX-P0048-V0004
% #FIELD_VALUE: 温度
% #HANDWRITTEN: 37.0℃
\fieldvalue{\handwritten{37.0℃}}

\vspace{2pt}
DO：
% #VALUE_ID: LEX-P0048-V0005
% #FIELD_VALUE: DO
% #HANDWRITTEN: 37.8\%
\fieldvalue{\handwritten{37.8\%}}

\vspace{2pt}
pH：
% #VALUE_ID: LEX-P0048-V0006
% #FIELD_VALUE: pH
% #HANDWRITTEN: 7.40
\fieldvalue{\handwritten{7.40}}

\vspace{2pt}
表面流量：
% #VALUE_ID: LEX-P0048-V0007
% #FIELD_VALUE: 表面流量
% #HANDWRITTEN: 20 ml/min
\fieldvalue{\handwritten{20}} ml/min

\vspace{2pt}
底部流量：
% #VALUE_ID: LEX-P0048-V0008
% #FIELD_VALUE: 底部流量
% #HANDWRITTEN: 10 ml/min
\fieldvalue{\handwritten{10}} ml/min

\vspace{2pt}
培养体积：
% #VALUE_ID: LEX-P0048-V0009
% #FIELD_VALUE: 培养体积
% #HANDWRITTEN: 30 L
\fieldvalue{\handwritten{30}} L

\vspace{2pt}
是否符合要求： 
% #VALUE_ID: LEX-P0048-V0010
% #FIELD_VALUE: 是否符合要求
% #HANDWRITTEN: 是（勾选）
\fieldvalue{\handwritten{是（勾选）}} \quad $\square$ 否
\end{minipage}
\\
\hline
\begin{minipage}[t]{\linewidth}
2）观察罐体内微载体分布情况，微载体上下分布不均--或开始出现较多绕团时，提高转速至40--50 rpm，每次调整至2--5rpm；
\end{minipage}
&
\begin{minipage}[t]{\linewidth}
是否调整转速：是 \quad
% #VALUE_ID: LEX-P0048-V0011
% #FIELD_VALUE: 是否调整转速
% #HANDWRITTEN: 否（勾选）
\fieldvalue{\handwritten{否（勾选）}}

\vspace{2pt}
调整时间：\underline{\hspace{3em}}/\underline{\hspace{3em}}

\vspace{2pt}
调整后转速：\underline{\hspace{3em}} rpm
\end{minipage}
\\
\hline
\begin{minipage}[t]{\linewidth}
3）取样检测：

取样前装，每袋体积$\geq$20ml，分别检测细胞密度、活率、pH、葡萄糖浓度和细胞生长状态（镜检）。

注：pH 检查需与细胞密度/活率检测分检检测，且 pH 为取样开口后首要检测项目；

当高碳 pH 与在线 pH 差值$>$0.2时，应对在线电极单点校准。
\end{minipage}
&
\begin{minipage}[t]{\linewidth}
取样时间：
% #VALUE_ID: LEX-P0048-V0012
% #FIELD_VALUE: 取样时间
% #HANDWRITTEN: 09：27
\fieldvalue{\handwritten{09：27}}

\vspace{2pt}
培养时长：
% #VALUE_ID: LEX-P0048-V0013
% #FIELD_VALUE: 培养时长
% #HANDWRITTEN: 48.3 h
\fieldvalue{\handwritten{48.3}} h

\vspace{2pt}
细胞密度：
% #VALUE_ID: LEX-P0048-V0014
% #FIELD_VALUE: 细胞密度
% #TODO #HANDWRITTEN: 6.6$\times10^{?}$; exponent unclear
\fieldvalue{\handwritten{6.6$\times10^{?}$}} 个细胞/ml

\vspace{2pt}
细胞活率：
% #VALUE_ID: LEX-P0048-V0015
% #FIELD_VALUE: 细胞活率
% #HANDWRITTEN: 86\%
\fieldvalue{\handwritten{86\%}}

\vspace{2pt}
葡萄糖浓度：
% #VALUE_ID: LEX-P0048-V0016
% #FIELD_VALUE: 葡萄糖浓度
% #HANDWRITTEN: 18.9 mmol/L
\fieldvalue{\handwritten{18.9}} mmol/L

\vspace{2pt}
葡萄糖 pH：
% #VALUE_ID: LEX-P0048-V0017
% #FIELD_VALUE: 葡萄糖 pH
% #TODO #HANDWRITTEN: 7.71; final digit uncertain
\fieldvalue{\handwritten{7.71}}

\vspace{2pt}
电极是否校准：
% #VALUE_ID: LEX-P0048-V0018
% #FIELD_VALUE: 电极是否校准
% #HANDWRITTEN: 是（勾选）
\fieldvalue{\handwritten{是（勾选）}} \quad 否

\vspace{2pt}
电极正常：
% #VALUE_ID: LEX-P0048-V0019
% #FIELD_VALUE: 电极正常
% #HANDWRITTEN: 是（勾选）
\fieldvalue{\handwritten{是（勾选）}} \quad $\square$ 否
\end{minipage}
\\
\hline
\multicolumn{2}{|l|}{
备注：
% #VALUE_ID: LEX-P0048-V0020
% #HANDWRITTEN: /
\fieldvalue{\handwritten{/}}
} \\
\hline
\end{tabularx}

\vspace{4pt}

操作人/日期：
% #VALUE_ID: LEX-P0048-V0021
% #FIELD_VALUE: 操作人
% #TODO #HANDWRITTEN: 乔某某; signature unclear
\fieldvalue{\handwritten{乔某某}}
\quad
% #VALUE_ID: LEX-P0048-V0022
% #FIELD_VALUE: 操作日期
% #HANDWRITTEN: 2025.06.1
\fieldvalue{\handwritten{2025.06.1}}

\hfill
复核人/日期：
% #VALUE_ID: LEX-P0048-V0023
% #FIELD_VALUE: 复核人
% #TODO #HANDWRITTEN: 刘岑岑; signature unclear
\fieldvalue{\handwritten{刘岑岑}}
\quad
% #VALUE_ID: LEX-P0048-V0024
% #FIELD_VALUE: 复核日期
% #HANDWRITTEN: 2025.06.1
\fieldvalue{\handwritten{2025.06.1}}
% LEXOID_PAGE_COMPLETED: 48/59

\newpage

\section*{米尔迈托赛注射液（Z2）批生产记录}

\noindent
\textbf{记录编号：} YZ-BPR-MF-PD-01-020-a \qquad
\textbf{版本号：} 1.3

\begin{tabularx}{\textwidth}{|X|X|X|X|}
\hline
品名 & 规格 & 生产批号 & 批量 \\ \hline
米尔迈托赛注射液 & 12ml；$6.0\times10^7$个细胞 &
% #VALUE_ID: LEX-P0049-V0001
% #FIELD_VALUE: 生产批号
% #TODO #HANDWRITTEN: A2323……；手写批号部分字迹不清
\fieldvalue{\handwritten{A2323……}} &
200--400袋 \\ \hline
\end{tabularx}

\subsection*{3.3 清场}

\noindent
\textbf{操作依据：}

\begin{enumerate}
\item 《生产区清场管理规程》 \hfill YZ-SMP-PD-03-001
\item 《生产、检验状态标识管理规程》 \hfill YZ-SMP-QA-04-014
\end{enumerate}

\noindent
生产工序：一级生物反应器细胞培养 \hfill 清场开始时间：

% #VALUE_ID: LEX-P0049-V0002
% #FIELD_VALUE: 清场开始时间
% #TODO #HANDWRITTEN: 07时？分；分钟数手写不清
\fieldvalue{\handwritten{07时？分}}

\noindent
清场房间编码：

% #VALUE_ID: LEX-P0049-V0003
% #FIELD_VALUE: 75\%酒精批号
% #TODO #HANDWRITTEN: D？050？01；手写批号部分字迹不清
\fieldvalue{\handwritten{D？050？01}}
\hfill 75\%酒精批号：

\vspace{0.5em}

\begin{tabularx}{\textwidth}{|X|c|}
\hline
\multicolumn{1}{|c|}{清场项目及标准} & 执行情况 \\ \hline
将本班次使用的各类工具配件恢复原位放置；
&
% #VALUE_ID: LEX-P0049-V0004
% #FIELD_VALUE: 将本班次使用的各类工具配件恢复原位放置——已执行
% #HANDWRITTEN: ✓
\fieldvalue{\handwritten{\checkmark}}已执行\quad □否
\\ \hline
将本班次生产有关的其它物品或物料清出生产现场；
&
% #VALUE_ID: LEX-P0049-V0005
% #FIELD_VALUE: 将本班次生产有关的其它物品或物料清出生产现场——已执行
% #HANDWRITTEN: ✓
\fieldvalue{\handwritten{\checkmark}}已执行\quad □否
\\ \hline
使用75\%酒精清洁设备、操作台面、不锈钢罐子、转移车、传递窗、门把手等；
&
% #VALUE_ID: LEX-P0049-V0006
% #FIELD_VALUE: 使用75\%酒精清洁设备、操作台面、不锈钢罐子、转移车、传递窗、门把手等——已执行
% #HANDWRITTEN: ✓
\fieldvalue{\handwritten{\checkmark}}已执行\quad □否
\\ \hline
使用75\%酒精喷洒地面，然后使用拖把擦拭地面；
&
% #VALUE_ID: LEX-P0049-V0007
% #FIELD_VALUE: 使用75\%酒精喷洒地面，然后使用拖把擦拭地面——已执行
% #HANDWRITTEN: ✓
\fieldvalue{\handwritten{\checkmark}}已执行\quad □否
\\ \hline
按照《生产、检验状态标识管理规程》更换房间、设备状态标识；
&
% #VALUE_ID: LEX-P0049-V0008
% #FIELD_VALUE: 按照《生产、检验状态标识管理规程》更换房间、设备状态标识——已执行
% #HANDWRITTEN: ✓
\fieldvalue{\handwritten{\checkmark}}已执行\quad □否
\\ \hline
将本班次操作产生的废弃物移出洁净生产区；
&
% #VALUE_ID: LEX-P0049-V0009
% #FIELD_VALUE: 将本班次操作产生的废弃物移出洁净生产区——已执行
% #HANDWRITTEN: ✓
\fieldvalue{\handwritten{\checkmark}}已执行\quad □否
\\ \hline
\end{tabularx}

\noindent
清场结束时间：

% #VALUE_ID: LEX-P0049-V0010
% #FIELD_VALUE: 清场结束时间
% #TODO #HANDWRITTEN: 09时？分；分钟数手写不清
\fieldvalue{\handwritten{09时？分}}

\noindent
清场人员/日期：

% #VALUE_ID: LEX-P0049-V0011
% #FIELD_VALUE: 清场人员
% #TODO #HANDWRITTEN: 姓名不清；手写姓名字迹无法辨认
\fieldvalue{\handwritten{（姓名不清）}}
\quad

% #VALUE_ID: LEX-P0049-V0012
% #FIELD_VALUE: 清场日期
% #HANDWRITTEN: 2025.06.01
\fieldvalue{\handwritten{2025.06.01}}

\vspace{0.5em}

\begin{tabularx}{\textwidth}{|X|}
\hline
清场检查 \\ \hline
% #VALUE_ID: LEX-P0049-V0013
% #FIELD_VALUE: 清场检查——合格
% #HANDWRITTEN: ✓
\fieldvalue{\handwritten{\checkmark}}合格 \\ \hline
□不合格，具体内容为： \\[1.8em] \hline
整改情况及整改人： \\[1.8em] \hline
\end{tabularx}

\vspace{0.5em}

\noindent
确认人/日期：

% #VALUE_ID: LEX-P0049-V0014
% #FIELD_VALUE: 确认人
% #TODO #HANDWRITTEN: 姓名不清；手写姓名字迹无法辨认
\fieldvalue{\handwritten{（姓名不清）}}
\quad

% #VALUE_ID: LEX-P0049-V0015
% #FIELD_VALUE: 确认日期
% #HANDWRITTEN: 2025.06.01
\fieldvalue{\handwritten{2025.06.01}}

\vfill

\begin{center}
第31页 共222页
\end{center}
% LEXOID_PAGE_COMPLETED: 49/59

\newpage

\section*{技术近托囊注射液（Z2）批生产记录}

\noindent
记录编号：YZ-BPR-MF-PD-01-020-a

\begin{tabularx}{\textwidth}{|>{\centering\arraybackslash}X|>{\centering\arraybackslash}X|>{\centering\arraybackslash}X|>{\centering\arraybackslash}X|}
\hline
品名 & 规格 & 生产批号 & 批量 \\
\hline
注射用盐酸金霉素 & 12ml；6.0$\times 10^7$个细胞 &
% #VALUE_ID: LEX-P0050-V0001
% #FIELD_VALUE: 生产批号
% #TODO #HANDWRITTEN: 4322?2W?b?01; 字迹不清
\fieldvalue{\handwritten{4322?2W?b?01}} &
200--480袋 \\
\hline
\end{tabularx}

\subsection*{4.\ 一级生物反应器细胞培养第三天}

\subsubsection*{4.1 生产前确认：}

\begin{tabularx}{\textwidth}{|>{\raggedright\arraybackslash}p{0.55\textwidth}|>{\raggedright\arraybackslash}X|}
\hline
\multicolumn{1}{|c|}{项目} & \multicolumn{1}{c|}{操作记录} \\
\hline
操作人员按《人员进出洁净生产区管理规程》要求进入洁净生产区； &
操作间名称及编码：
% #VALUE_ID: LEX-P0050-V0002
% #FIELD_VALUE: 操作间名称及编码
% TODO #HANDWRITTEN: 字迹不清
\fieldvalue{\handwritten{字迹不清}} \\
\hline
确认生产区操作间在清场有效期内（C级日清洁期3天）： &
是否在清洁有效期内：\quad
% #VALUE_ID: LEX-P0050-V0003
% #FIELD_VALUE: 是否在清洁有效期内——是
% #HANDWRITTEN: 已勾选
是\ \fieldvalue{\handwritten{\checkmark}}\quad 否

\vspace{2pt}
清洁日期：
% #VALUE_ID: LEX-P0050-V0004
% #FIELD_VALUE: 清洁日期
% #HANDWRITTEN: 2025.06.14
\fieldvalue{\handwritten{2025.06.14}} \\
\hline
确认操作间温度18--26$^\circ$C，湿度30\%--75\%，相对压差$>0$Pa（操作间对C级走廊）： &
温度：
% #VALUE_ID: LEX-P0050-V0005
% #FIELD_VALUE: 温度
% #HANDWRITTEN: 21.1
\fieldvalue{\handwritten{21.1}} $^\circ$C

湿度：
% #VALUE_ID: LEX-P0050-V0006
% #FIELD_VALUE: 湿度
% #HANDWRITTEN: 55.5
\fieldvalue{\handwritten{55.5}}\%

\vspace{2pt}
压差：
% #VALUE_ID: LEX-P0050-V0007
% #FIELD_VALUE: 压差
% #HANDWRITTEN: 16
\fieldvalue{\handwritten{16}}\ Pa

\vspace{2pt}
是否符合要求：
% #VALUE_ID: LEX-P0050-V0008
% #FIELD_VALUE: 是否符合要求——是
% #HANDWRITTEN: 已勾选
是\ \fieldvalue{\handwritten{\checkmark}}\quad 否 \\
\hline
确认操作间内无上批产品遗留物、文件及与本产品生产无关物料； &
% #VALUE_ID: LEX-P0050-V0009
% #FIELD_VALUE: 上批产品及无关物料检查
% #HANDWRITTEN: 已勾选
是\ \fieldvalue{\handwritten{\checkmark}}\quad 否 \\
\hline
检查废弃的批记录及其它相关配套文件均与生产指令相对应，且均受控； &
% #VALUE_ID: LEX-P0050-V0010
% #FIELD_VALUE: 批记录及配套文件检查
% #HANDWRITTEN: 已勾选
是\ \fieldvalue{\handwritten{\checkmark}}\quad 否 \\
\hline
确认仪器、且在验证有效期内，计量器具有“合格证”且在有效期内，详见《主要设备/仪表确认表》； &
% #VALUE_ID: LEX-P0050-V0011
% #FIELD_VALUE: 仪器及计量器具确认
% #HANDWRITTEN: 已勾选
是\ \fieldvalue{\handwritten{\checkmark}}\quad 否 \\
\hline
反应器气源压力应在如下范围：压缩空气进气压力0.05--0.15MPa、二氧化碳进气压力0.05--0.15MPa、氧气进气压力0.05--0.15MPa； &
压缩空气：
% #VALUE_ID: LEX-P0050-V0012
% #FIELD_VALUE: 压缩空气进气压力
% #HANDWRITTEN: 0.12
\fieldvalue{\handwritten{0.12}}\ MPa

二氧化碳：
% #VALUE_ID: LEX-P0050-V0013
% #FIELD_VALUE: 二氧化碳进气压力
% #HANDWRITTEN: 0.12
\fieldvalue{\handwritten{0.12}}\ MPa

氧气：
% #VALUE_ID: LEX-P0050-V0014
% #FIELD_VALUE: 氧气进气压力
% #HANDWRITTEN: 0.10
\fieldvalue{\handwritten{0.10}}\ MPa \\
\hline
将操作间状态标识更换为“生产加工状态卡”，写明品名、规格、批号等，并确认与本批生产指令一致； &
% #VALUE_ID: LEX-P0050-V0015
% #FIELD_VALUE: 生产加工状态卡确认
% #HANDWRITTEN: 已勾选
是\ \fieldvalue{\handwritten{\checkmark}}\quad 否 \\
\hline
核对本次生产用物料耗材的品名、规格、批号、数量，检查物料外包装的完整性以及是否有有效期/复验期内，并记录相关信息，详见《物料确认表》； &
% #VALUE_ID: LEX-P0050-V0016
% #FIELD_VALUE: 物料耗材确认
% #HANDWRITTEN: 已勾选
是\ \fieldvalue{\handwritten{\checkmark}}\quad 否 \\
\hline
备注： &
% #VALUE_ID: LEX-P0050-V0017
% #HANDWRITTEN: \checkmark
\handwritten{\checkmark} \\
\hline
\end{tabularx}

\vspace{0.5em}

确认人/日期：
% #VALUE_ID: LEX-P0050-V0018
% #FIELD_VALUE: 确认人
% #HANDWRITTEN: 李楠
\fieldvalue{\handwritten{李楠}}
\quad
% #VALUE_ID: LEX-P0050-V0019
% #FIELD_VALUE: 确认日期
% #HANDWRITTEN: 2025.06.12
\fieldvalue{\handwritten{2025.06.12}}
\hfill
复核人/日期：
% #VALUE_ID: LEX-P0050-V0020
% #FIELD_VALUE: 复核人
% #HANDWRITTEN: 李军
\fieldvalue{\handwritten{李军}}
\quad
% #VALUE_ID: LEX-P0050-V0021
% #FIELD_VALUE: 复核日期
% #HANDWRITTEN: 2025.6.12
\fieldvalue{\handwritten{2025.6.12}}

\vfill

\begin{center}
第32页 共222页
\end{center}
% LEXOID_PAGE_COMPLETED: 50/59

\newpage

\section*{艾米达托菲注射液（Z2）批生产记录}

\noindent
记录编号：YZ-BPR-MF-PD-01-020-a\hfill 版本号：1.3

\begin{tabularx}{\textwidth}{|X|X|X|X|}
\hline
品名 & 规格 & 生产批号 & 批量\\
\hline
% #VALUE_ID: LEX-P0051-V0001
% #FIELD_VALUE: 品名
\fieldvalue{艾米达托菲注射液} &
% #VALUE_ID: LEX-P0051-V0002
% #FIELD_VALUE: 规格
\fieldvalue{12ml；$6.0\times10^{7}$个细胞} &
% #VALUE_ID: LEX-P0051-V0003
% #FIELD_VALUE: 生产批号
% #TODO #HANDWRITTEN: 4332102056...; 手写批号部分字符难以辨认
\fieldvalue{\handwritten{4332102056...}} &
% #VALUE_ID: LEX-P0051-V0004
% #FIELD_VALUE: 批量
\fieldvalue{200--480袋}\\
\hline
\end{tabularx}

\subsection*{4.1.1 主要设备/仪表确认表}

\begin{tabularx}{\textwidth}{|c|X|X|X|X|c|}
\hline
序号 & 设备名称 & 品牌/型号 & 设备编号 & 检查项 & 确认结果\\
\hline
1 &
% #VALUE_ID: LEX-P0051-V0005
% #FIELD_VALUE: 设备名称（第1项）
\fieldvalue{生物反应器控制台（主机）} &
% #VALUE_ID: LEX-P0051-V0006
% #FIELD_VALUE: 品牌/型号（第1项）
\fieldvalue{华鑫生物} &
% #VALUE_ID: LEX-P0051-V0007
% #FIELD_VALUE: 设备编号（第1项）
% #HANDWRITTEN: 1123610
\fieldvalue{\handwritten{1123610}} &
\multirow{4}{*}{%
\begin{minipage}{\linewidth}
1. 设备器具有“良好”标识；\\
2. 设备铭牌在有效期内；\\
3. 计量器具有“合格”且在有效期内。
\end{minipage}} &
% #VALUE_ID: LEX-P0051-V0008
% #FIELD_VALUE: 确认结果（第1项）
\fieldvalue{是}\quad 否\\
\hline
2 &
% #VALUE_ID: LEX-P0051-V0009
% #FIELD_VALUE: 设备名称（第2项）
\fieldvalue{生物反应器罐体} &
% #VALUE_ID: LEX-P0051-V0010
% #FIELD_VALUE: 品牌/型号（第2项）
\fieldvalue{华鑫生物} &
% #VALUE_ID: LEX-P0051-V0011
% #FIELD_VALUE: 设备编号（第2项）
% #TODO #HANDWRITTEN: V4560502-22038; 首字符及中间字符较模糊
\fieldvalue{\handwritten{V4560502-22038}} &
&
% #VALUE_ID: LEX-P0051-V0012
% #FIELD_VALUE: 确认结果（第2项）
\fieldvalue{是}\quad 否\\
\hline
3 &
% #VALUE_ID: LEX-P0051-V0013
% #FIELD_VALUE: 设备名称（第3项）
\fieldvalue{接管机} &
% #VALUE_ID: LEX-P0051-V0014
% #FIELD_VALUE: 品牌/型号（第3项）
\fieldvalue{上海赛芙} &
% #VALUE_ID: LEX-P0051-V0015
% #FIELD_VALUE: 设备编号（第3项）
% #HANDWRITTEN: 1128556
\fieldvalue{\handwritten{1128556}} &
&
% #VALUE_ID: LEX-P0051-V0016
% #FIELD_VALUE: 确认结果（第3项）
\fieldvalue{是}\quad 否\\
\hline
4 &
% #VALUE_ID: LEX-P0051-V0017
% #FIELD_VALUE: 设备名称（第4项）
\fieldvalue{封管机} &
% #VALUE_ID: LEX-P0051-V0018
% #FIELD_VALUE: 品牌/型号（第4项）
\fieldvalue{上海赛芙} &
% #VALUE_ID: LEX-P0051-V0019
% #FIELD_VALUE: 设备编号（第4项）
% #HANDWRITTEN: 1123612
\fieldvalue{\handwritten{1123612}} &
&
% #VALUE_ID: LEX-P0051-V0020
% #FIELD_VALUE: 确认结果（第4项）
\fieldvalue{是}\quad 否\\
\hline
5 & & &
% #VALUE_ID: LEX-P0051-V0021
% #FIELD_VALUE: 设备编号（第5项）
% #TODO #HANDWRITTEN: 王? 03-06-12; 字迹及字段含义不清
\fieldvalue{\handwritten{王? 03-06-12}} &
& \\
\hline
\multicolumn{6}{|l|}{备注：%
% #VALUE_ID: LEX-P0051-V0022
% #HANDWRITTEN: /
\handwritten{/}}\\
\hline
\end{tabularx}

\noindent
确认人/日期：%
% #VALUE_ID: LEX-P0051-V0023
% #FIELD_VALUE: 确认人
% #TODO #HANDWRITTEN: 李?娟; 姓名部分字迹模糊
\fieldvalue{\handwritten{李?娟}}\quad
% #VALUE_ID: LEX-P0051-V0024
% #FIELD_VALUE: 确认日期
% #HANDWRITTEN: 2025.06.12
\fieldvalue{\handwritten{2025.06.12}}
\hfill
复核人/日期：%
% #VALUE_ID: LEX-P0051-V0025
% #FIELD_VALUE: 复核人
% #TODO #HANDWRITTEN: 李?宇; 姓名部分字迹模糊
\fieldvalue{\handwritten{李?宇}}\quad
% #VALUE_ID: LEX-P0051-V0026
% #FIELD_VALUE: 复核日期
% #HANDWRITTEN: 2025.06.12
\fieldvalue{\handwritten{2025.06.12}}

\subsection*{4.1.2 物料确认表}

\begin{tabularx}{\textwidth}{|c|X|X|X|X|c|X|}
\hline
序号 & 名称 & 品牌 & 规格 & 批号 & 数量 & 是否符合要求\\
\hline
1 &
% #VALUE_ID: LEX-P0051-V0027
% #FIELD_VALUE: 名称（第1项）
% #TODO #HANDWRITTEN: 50ml/袋注射用...; 手写内容部分难以辨认
\fieldvalue{\handwritten{50ml/袋注射用...}} &
% #VALUE_ID: LEX-P0051-V0028
% #FIELD_VALUE: 品牌（第1项）
% #TODO #HANDWRITTEN: 全成生物; 字迹略模糊
\fieldvalue{\handwritten{全成生物}} &
% #VALUE_ID: LEX-P0051-V0029
% #FIELD_VALUE: 规格（第1项）
% #HANDWRITTEN: 50ml/袋
\fieldvalue{\handwritten{50ml/袋}} &
% #VALUE_ID: LEX-P0051-V0030
% #FIELD_VALUE: 批号（第1项）
% #TODO #HANDWRITTEN: M120?050?02-02; 批号部分字符难以辨认
\fieldvalue{\handwritten{M120?050?02-02}} &
% #VALUE_ID: LEX-P0051-V0031
% #FIELD_VALUE: 数量（第1项）
% #HANDWRITTEN: 3
\fieldvalue{\handwritten{3}} &
% #VALUE_ID: LEX-P0051-V0032
% #FIELD_VALUE: 是否符合要求（第1项）
\fieldvalue{是}\quad 否\\
\hline
2 & & &
&
% #VALUE_ID: LEX-P0051-V0033
% #FIELD_VALUE: 第2项表内手写记录
% #TODO #HANDWRITTEN: 李?宇 2025.06.12; 位于第2行批号区域，字段含义不明确
\fieldvalue{\handwritten{李?宇 2025.06.12}} &
& \\
\hline
\multicolumn{7}{|l|}{备注：}\\
\hline
\end{tabularx}

\noindent
确认人/日期：%
% #VALUE_ID: LEX-P0051-V0034
% #FIELD_VALUE: 确认人
% #TODO #HANDWRITTEN: 李?娟; 姓名部分字迹模糊
\fieldvalue{\handwritten{李?娟}}\quad
% #VALUE_ID: LEX-P0051-V0035
% #FIELD_VALUE: 确认日期
% #HANDWRITTEN: 2025.06.12
\fieldvalue{\handwritten{2025.06.12}}
\hfill
复核人/日期：%
% #VALUE_ID: LEX-P0051-V0036
% #FIELD_VALUE: 复核人
% #TODO #HANDWRITTEN: 李?宇; 姓名部分字迹模糊
\fieldvalue{\handwritten{李?宇}}\quad
% #VALUE_ID: LEX-P0051-V0037
% #FIELD_VALUE: 复核日期
% #HANDWRITTEN: 2025.06.12
\fieldvalue{\handwritten{2025.06.12}}

\subsection*{4.2 操作内容}

\begin{tabularx}{\textwidth}{|X|X|}
\hline
项目 & 操作记录\\
\hline
\end{tabularx}

\begin{center}
第 33 页 共 222 页
\end{center}
% LEXOID_PAGE_COMPLETED: 51/59

\newpage

\begin{center}
\textbf{广州赛莱拉生物科技有限公司}\\
\textbf{\large 艾米迈托昔注射液（Z2）批生产记录}
\end{center}

\noindent
\begin{tabularx}{\textwidth}{@{}lXlX@{}}
记录编号：&
% #VALUE_ID: LEX-P0052-V0001
% #FIELD_VALUE: 记录编号
\fieldvalue{YZ-BPR-MF-PD-01-020-a}
&
版本号：&
% #VALUE_ID: LEX-P0052-V0002
% #FIELD_VALUE: 版本号
\fieldvalue{1.3}
\end{tabularx}

\vspace{2mm}

\renewcommand{\arraystretch}{1.35}
\begin{tabularx}{\textwidth}{|>{\centering\arraybackslash}p{0.22\textwidth}|>{\centering\arraybackslash}p{0.29\textwidth}|>{\centering\arraybackslash}p{0.22\textwidth}|>{\centering\arraybackslash}X|}
\hline
\textbf{品名} & \textbf{规格} & \textbf{生产批号} & \textbf{批量} \\
\hline
% #VALUE_ID: LEX-P0052-V0003
% #FIELD_VALUE: 品名
\fieldvalue{艾米迈托昔注射液}
&
% #VALUE_ID: LEX-P0052-V0004
% #FIELD_VALUE: 规格
\fieldvalue{12ml：$6.0\times10^{7}$个细胞}
&
% #VALUE_ID: LEX-P0052-V0005
% #FIELD_VALUE: 生产批号
% #TODO #HANDWRITTEN: 052?22?50?; 手写批号部分字符难以辨认
\fieldvalue{\handwritten{052?22?50?}}
&
% #VALUE_ID: LEX-P0052-V0006
% #FIELD_VALUE: 批量
\fieldvalue{200--480袋}
\\
\hline
\end{tabularx}

\vspace{1mm}

\begin{tabularx}{\textwidth}{|p{0.48\textwidth}|X|}
\hline
\textbf{1）参数检查：}

检查培养参数控制值是否在如下范围：表面通气20--25ml/min，底部通气10$\pm$3ml/min，转速35--50rpm，温度37.0$\pm$2.0℃，溶氧$>$30\%，pH6.70--7.60。
&
\begin{tabular}{@{}ll@{}}
检查时间：&
% #VALUE_ID: LEX-P0052-V0007
% #FIELD_VALUE: 检查时间
% #HANDWRITTEN: 06：50
\fieldvalue{\handwritten{06：50}}\\
搅拌转速：&
% #VALUE_ID: LEX-P0052-V0008
% #FIELD_VALUE: 搅拌转速
% #HANDWRITTEN: 25
\fieldvalue{\handwritten{25}}\ rpm\\
温度：&
% #VALUE_ID: LEX-P0052-V0009
% #FIELD_VALUE: 温度
% #HANDWRITTEN: 37.0
\fieldvalue{\handwritten{37.0}}℃\\
DO：&
% #VALUE_ID: LEX-P0052-V0010
% #FIELD_VALUE: DO
% #HANDWRITTEN: 66.6
\fieldvalue{\handwritten{66.6}}\%\\
pH：&
% #VALUE_ID: LEX-P0052-V0011
% #FIELD_VALUE: pH
% #HANDWRITTEN: 7.41
\fieldvalue{\handwritten{7.41}}\\
表面流量：&
% #VALUE_ID: LEX-P0052-V0012
% #FIELD_VALUE: 表面流量
% #HANDWRITTEN: 20
\fieldvalue{\handwritten{20}}\ ml/min\\
底部流量：&
% #VALUE_ID: LEX-P0052-V0013
% #FIELD_VALUE: 底部流量
% #HANDWRITTEN: 10
\fieldvalue{\handwritten{10}}\ ml/min\\
培养体积：&
% #VALUE_ID: LEX-P0052-V0014
% #FIELD_VALUE: 培养体积
% #HANDWRITTEN: 1.9
\fieldvalue{\handwritten{1.9}}\ L\\
是否符合要求：&
% #VALUE_ID: LEX-P0052-V0015
% #FIELD_VALUE: 是否符合要求
% #HANDWRITTEN: 是（勾选）
\fieldvalue{\handwritten{是（勾选）}}
\end{tabular}
\\
\hline
\textbf{2）观察罐体内微镜体分布情况，微载体上下分布不均一或开始出现较多黏团时，提高转速至40--50rpm，调校2--5rpm；}
&
\begin{tabular}{@{}ll@{}}
是否调整转速：&
% #VALUE_ID: LEX-P0052-V0016
% #FIELD_VALUE: 是否调整转速
% #HANDWRITTEN: 否（勾选）
\fieldvalue{\handwritten{否（勾选）}}\\
调整时间：&\underline{\hspace{25mm}}：\underline{\hspace{15mm}}\\
调整后转速：&\underline{\hspace{35mm}}\ rpm
\end{tabular}
\\
\hline
\textbf{3）取样检测：}

取样两袋，包袋体积$>$20ml，分别检测细胞密度、活率、pH、葡萄糖浓度和细胞生长状态（镜检）。

注：pH检测需与细胞密度/活率检测分袋检测，且pH为取样经开口后直接检测项目！

当离线pH与在线pH差值$\geq0.2$时，应对在线电极做单点校准。
&
\begin{tabular}{@{}ll@{}}
取样时间：&
% #VALUE_ID: LEX-P0052-V0017
% #FIELD_VALUE: 取样时间
% #HANDWRITTEN: 06：51
\fieldvalue{\handwritten{06：51}}\\
培养时长：&
% #VALUE_ID: LEX-P0052-V0018
% #FIELD_VALUE: 培养时长
% #HANDWRITTEN: 70
\fieldvalue{\handwritten{70}}\ h\\
细胞密度：&
% #VALUE_ID: LEX-P0052-V0019
% #FIELD_VALUE: 细胞密度
% #HANDWRITTEN: 1.46$\times10^{5}$
\fieldvalue{\handwritten{$1.46\times10^{5}$}}\ 个细胞/ml\\
细胞活率：&
% #VALUE_ID: LEX-P0052-V0020
% #FIELD_VALUE: 细胞活率
% #HANDWRITTEN: 90
\fieldvalue{\handwritten{90}}\%\\
葡萄糖浓度：&
% #VALUE_ID: LEX-P0052-V0021
% #FIELD_VALUE: 葡萄糖浓度
% #HANDWRITTEN: 17.0
\fieldvalue{\handwritten{17.0}}\ mmol/L\\
离线pH：&
% #VALUE_ID: LEX-P0052-V0022
% #FIELD_VALUE: 离线pH
% #HANDWRITTEN: 7.24
\fieldvalue{\handwritten{7.24}}\\
电镜是否检查：&
% #VALUE_ID: LEX-P0052-V0023
% #FIELD_VALUE: 电镜是否检查
% #HANDWRITTEN: 否（勾选）
\fieldvalue{\handwritten{否（勾选）}}\\
镜检正常：&
% #VALUE_ID: LEX-P0052-V0024
% #FIELD_VALUE: 镜检正常
% #HANDWRITTEN: 是（勾选）
\fieldvalue{\handwritten{是（勾选）}}
\end{tabular}
\\
\hline
备注：&
% #VALUE_ID: LEX-P0052-V0025
% #HANDWRITTEN: /
\handwritten{/}
\\
\hline
\end{tabularx}

\vspace{2mm}

\noindent
操作人/日期：
% #VALUE_ID: LEX-P0052-V0026
% #FIELD_VALUE: 操作人
% #TODO #HANDWRITTEN: 徐梅; 姓名部分笔画不完全清晰
\fieldvalue{\handwritten{徐梅}}
\quad
% #VALUE_ID: LEX-P0052-V0027
% #FIELD_VALUE: 操作日期
% #HANDWRITTEN: 2025.06.12
\fieldvalue{\handwritten{2025.06.12}}
\qquad
复核人/日期：
% #VALUE_ID: LEX-P0052-V0028
% #FIELD_VALUE: 复核人
% #TODO #HANDWRITTEN: 赵宇; 姓名部分笔画不完全清晰
\fieldvalue{\handwritten{赵宇}}
\quad
% #VALUE_ID: LEX-P0052-V0029
% #FIELD_VALUE: 复核日期
% #HANDWRITTEN: 2025.06.12
\fieldvalue{\handwritten{2025.06.12}}

\vfill

\begin{center}
第34页 共222页
\end{center}
% LEXOID_PAGE_COMPLETED: 52/59

\newpage

\begin{center}
\textbf{艾来托莫注射液（Z2）批生产记录}
\end{center}

\noindent
记录编号：YZ-BPR-MF-PD-01-020-4

\begin{center}
\begin{tabular}{|c|c|c|c|}
\hline
品名 & 规格 & 生产批号 & 批量 \\
\hline
艾来托莫注射液 & 12ml：$6.0\times10^{7}$个细胞 &
% #VALUE_ID: LEX-P0053-V0002
% #FIELD_VALUE: 生产批号
% #TODO #HANDWRITTEN: 453?22?06?001; 蓝色手写批号部分字迹不清
\fieldvalue{\handwritten{453?22?06?001}} &
200--480袋 \\
\hline
\end{tabular}
\end{center}

\noindent
部门：
% #VALUE_ID: LEX-P0053-V0001
% #FIELD_VALUE: 部门
\fieldvalue{原液}
\hfill
版本号：1/3

\section*{4.3 清场}

\noindent
操作依据：

\begin{enumerate}
\item 《生产区清场管理规程》\hfill YZ-SMP-PD-03-001
\item 《生产、检验状态标识管理规程》\hfill YZ-SMP-QA-04-014
\item 生产工序：一级生物反应器细胞培养
\end{enumerate}

\noindent
清场开始时间：
% #VALUE_ID: LEX-P0053-V0003
% #FIELD_VALUE: 清场开始时间--时
% #HANDWRITTEN: 07
\fieldvalue{\handwritten{07}}时
% #VALUE_ID: LEX-P0053-V0004
% #FIELD_VALUE: 清场开始时间--分
% #HANDWRITTEN: 14
\fieldvalue{\handwritten{14}}分

\noindent
清场房间编码：
% #VALUE_ID: LEX-P0053-V0005
% #FIELD_VALUE: 清场房间编码
% #HANDWRITTEN: 5260
\fieldvalue{\handwritten{5260}}
\hfill
75\%酒精批号：
% #VALUE_ID: LEX-P0053-V0006
% #FIELD_VALUE: 75\%酒精批号
% #HANDWRITTEN: 2025060701
\fieldvalue{\handwritten{2025060701}}

\begin{center}
\begin{tabular}{|p{0.70\textwidth}|p{0.22\textwidth}|}
\hline
\centering 清场项目及标准 & \centering 执行情况 \arraybackslash \\
\hline
将本班次使用的各类工具配件恢复原位放置；
&
% #VALUE_ID: LEX-P0053-V0007
% #FIELD_VALUE: 清场项目1执行情况
% #HANDWRITTEN: 勾选已执行
\fieldvalue{\handwritten{\checkmark}}已执行\quad $\square$否
\\
\hline
将本班次生产有关的其它物品或物料清出生产现场；
&
% #VALUE_ID: LEX-P0053-V0008
% #FIELD_VALUE: 清场项目2执行情况
% #HANDWRITTEN: 勾选已执行
\fieldvalue{\handwritten{\checkmark}}已执行\quad $\square$否
\\
\hline
使用75\%酒精清洁设备、操作台面、不锈钢器皿、转移车、传递窗、门把手等；
&
% #VALUE_ID: LEX-P0053-V0009
% #FIELD_VALUE: 清场项目3执行情况
% #HANDWRITTEN: 勾选已执行
\fieldvalue{\handwritten{\checkmark}}已执行\quad $\square$否
\\
\hline
使用75\%酒精喷洒地面，然后使用拖把擦拭地面；
&
% #VALUE_ID: LEX-P0053-V0010
% #FIELD_VALUE: 清场项目4执行情况
% #HANDWRITTEN: 勾选已执行
\fieldvalue{\handwritten{\checkmark}}已执行\quad $\square$否
\\
\hline
按照《生产、检验状态标识管理规程》更换房间、设备状态标识；
&
% #VALUE_ID: LEX-P0053-V0011
% #FIELD_VALUE: 清场项目5执行情况
% #HANDWRITTEN: 勾选已执行
\fieldvalue{\handwritten{\checkmark}}已执行\quad $\square$否
\\
\hline
将本班次操作产生的废弃物移出洁净生产区；
&
% #VALUE_ID: LEX-P0053-V0012
% #FIELD_VALUE: 清场项目6执行情况
% #HANDWRITTEN: 勾选已执行
\fieldvalue{\handwritten{\checkmark}}已执行\quad $\square$否
\\
\hline
\end{tabular}
\end{center}

\noindent
清场结束时间：
% #VALUE_ID: LEX-P0053-V0013
% #FIELD_VALUE: 清场结束时间--时
% #HANDWRITTEN: 07
\fieldvalue{\handwritten{07}}时
% #VALUE_ID: LEX-P0053-V0014
% #FIELD_VALUE: 清场结束时间--分
% #HANDWRITTEN: 21
\fieldvalue{\handwritten{21}}分

\noindent
清场人/日期：
% #VALUE_ID: LEX-P0053-V0015
% #FIELD_VALUE: 清场人
% #TODO #HANDWRITTEN: 梅梅; 姓名手写字迹不完全清晰
\fieldvalue{\handwritten{梅梅}}
\quad
% #VALUE_ID: LEX-P0053-V0016
% #FIELD_VALUE: 清场日期
% #HANDWRITTEN: 2025.06.12
\fieldvalue{\handwritten{2025.06.12}}

\begin{center}
\begin{tabular}{|p{0.92\textwidth}|}
\hline
清场检查 \\
\hline
% #VALUE_ID: LEX-P0053-V0017
% #FIELD_VALUE: 清场检查结果
% #HANDWRITTEN: 勾选合格
\fieldvalue{\handwritten{\checkmark}}合格 \\
\hline
$\square$不合格，具体内容为： \\
\hline
整改情况及整改人： \\
\hline
\end{tabular}
\end{center}

\noindent
确认人/日期：
% #VALUE_ID: LEX-P0053-V0018
% #FIELD_VALUE: 确认人
% #TODO #HANDWRITTEN: 李平; 姓名手写字迹不清
\fieldvalue{\handwritten{李平}}
\quad
% #VALUE_ID: LEX-P0053-V0019
% #FIELD_VALUE: 确认日期
% #HANDWRITTEN: 2025.06.12
\fieldvalue{\handwritten{2025.06.12}}

\begin{center}
第35页 共222页
\end{center}
% LEXOID_PAGE_COMPLETED: 53/59

\newpage

\begin{center}
\textbf{艾米迈托赛注射液（Z2）批生产记录}

\vspace{2mm}
\begin{tabularx}{\textwidth}{|>{\centering\arraybackslash}X|>{\centering\arraybackslash}X|>{\centering\arraybackslash}X|>{\centering\arraybackslash}X|}
\hline
品名 & 规格 & 生产批号 & 批量 \\
\hline
% #VALUE_ID: LEX-P0054-V0001
% #FIELD_VALUE: 品名
\fieldvalue{艾米迈托赛注射液} &
% #VALUE_ID: LEX-P0054-V0002
% #FIELD_VALUE: 规格
\fieldvalue{12ml；$6.0\times10^{7}$个细胞} &
% #VALUE_ID: LEX-P0054-V0003
% #FIELD_VALUE: 生产批号
% #HANDWRITTEN: 253220250601
\fieldvalue{\handwritten{253220250601}} &
% #VALUE_ID: LEX-P0054-V0004
% #FIELD_VALUE: 批量
\fieldvalue{200--480袋} \\
\hline
\end{tabularx}
\end{center}

\noindent\textbf{5.\ 一级生物反应器细胞培养第四天}

\medskip
\noindent\textbf{5.1 生产前确认：}

\begin{tabularx}{\textwidth}{|p{0.56\textwidth}|X|}
\hline
\centering 项目 & \centering 操作记录 \tabularnewline
\hline
操作人员按《人员进出洁净生产区管理规程》要求进入洁净生产间； &
操作间名称及编码：\rule{0.55\linewidth}{0.4pt}
\\
\hline
确认生产区操作间在清洁有效期内（C级日清洁期为3天）； &
是否在清洁有效期内：\quad $\square$是\quad 清洁有效期：\rule{0.25\linewidth}{0.4pt}\quad $\square$否
\\
\hline
确认操作间温度$18$--$26^\circ$C，湿度$30\%$--$75\%$，相对压差$>$0Pa（操作间对C级走廊）； &
温度：\rule{0.12\linewidth}{0.4pt}$^\circ$C\quad 湿度：\rule{0.12\linewidth}{0.4pt}\%\\
& 压差：\rule{0.12\linewidth}{0.4pt}Pa\\
& 是否符合要求：$\square$是\quad $\square$否
\\
\hline
确认操作间内无尘土、日落物、文件及与本产品生产无关的物料； &
% #VALUE_ID: LEX-P0054-V0005
% #FIELD_VALUE: 确认操作间内无尘土、日落物、文件及与本产品生产无关的物料
% #HANDWRITTEN: √
\fieldvalue{\handwritten{√}}\quad $\square$是\quad $\square$否\\[-1mm]
% #VALUE_ID: LEX-P0054-V0006
% #FIELD_VALUE: 确认日期
% #HANDWRITTEN: 25.06.12
\fieldvalue{\handwritten{25.06.12}}
\\
\hline
检查发放的批记录及其它相关配套文件均与生产指令相对应，且均受控； &
$\square$是\quad $\square$否
\\
\hline
确认设备完好，且在验证有效期内，计量器具有“合格证”且在有效期内，详见《主要设备仪/仪表确认表》； &
$\square$是\quad $\square$否
\\
\hline
反应器气源压力控制在如下范围：压缩空气进气压力
$0.05$--$0.15$MPa、二氧化碳进气压力$0.05$--$0.15$MPa、氧气进气压力
$0.05$--$0.15$MPa； &
压缩空气：\rule{0.12\linewidth}{0.4pt}MPa\\
& 二氧化碳：\rule{0.12\linewidth}{0.4pt}MPa\quad 氧气：\rule{0.12\linewidth}{0.4pt}MPa
\\
\hline
将操作间状态标识更换为“生产加工状态卡”，写明品名、规格、批号等，并确认与本批生产指令一致； &
$\square$是\quad $\square$否
\\
\hline
核对本次生产用物料耗材的品名、规格、批号、数量，检查物料外包装的完好性以及是否在有效期/复验期内，并记录相关信息，详见《物料确认表》； &
$\square$是\quad $\square$否
\\
\hline
备注： &
\\[8mm]
\hline
\end{tabularx}

\vspace{4mm}

\noindent
确认人/日期：
% #VALUE_ID: LEX-P0054-V0007
% #FIELD_VALUE: 确认人/日期
% #HANDWRITTEN: /（字迹不清）
\fieldvalue{\handwritten{/}}%
\hfill
复核人/日期：
% #VALUE_ID: LEX-P0054-V0008
% #FIELD_VALUE: 复核人/日期
% #HANDWRITTEN: /（字迹不清）
\fieldvalue{\handwritten{/}}

\vfill
\begin{center}
第 36 页 共 222 页
\end{center}

% TODO: A blue diagonal pen stroke visibly crosses the form; it is not a discernible logical field value and is not transcribed.
% LEXOID_PAGE_COMPLETED: 54/59

\newpage

\begin{center}
\textbf{华生生物制药科技（北京）有限公司}\\
\textit{艾米迭托赛注射液（Z2）批生产记录}
\end{center}

记录编号：YZ-BPR-MF-PD-01-020-a\hfill 原件

\renewcommand{\arraystretch}{1.35}
\small
\begin{tabularx}{\textwidth}{|>{\centering\arraybackslash}X|>{\centering\arraybackslash}X|>{\centering\arraybackslash}X|>{\centering\arraybackslash}X|}
\hline
品名 & 规格 & 生产批号 & 批量\\
\hline
% #VALUE_ID: LEX-P0055-V0001
% #FIELD_VALUE: 品名
\fieldvalue{艾米迭托赛注射液} &
% #VALUE_ID: LEX-P0055-V0002
% #FIELD_VALUE: 规格
\fieldvalue{12ml；6.0$\times 10^{7}$个细胞} &
% #VALUE_ID: LEX-P0055-V0003
% #FIELD_VALUE: 生产批号
% #HANDWRITTEN: B?31?2?6?01
\fieldvalue{\handwritten{B?31?2?6?01}} &
% #VALUE_ID: LEX-P0055-V0004
% #FIELD_VALUE: 批量
\fieldvalue{200--480\,瓶}\\
\hline
\end{tabularx}

\normalsize
\medskip
\textbf{5.1.1 主要设备/仪器确认表}

\small
\begin{tabularx}{\textwidth}{|c|X|X|X|X|c|}
\hline
序号 & 设备名称 & 品牌/型号 & 设备编号 & 检查项 & 确认结果\\
\hline
1 &
% #VALUE_ID: LEX-P0055-V0005
% #FIELD_VALUE: 设备名称
\fieldvalue{生物反应器控制台（主机）} &
% #VALUE_ID: LEX-P0055-V0006
% #FIELD_VALUE: 品牌/型号
\fieldvalue{华鑫生物} &
 & & $\square$是\quad $\square$否\\
\hline
2 &
% #VALUE_ID: LEX-P0055-V0007
% #FIELD_VALUE: 设备名称
\fieldvalue{生物反应器罐体} &
% #VALUE_ID: LEX-P0055-V0008
% #FIELD_VALUE: 品牌/型号
\fieldvalue{华鑫生物} &
 & & $\square$是\quad $\square$否\\
\hline
3 &
% #VALUE_ID: LEX-P0055-V0009
% #FIELD_VALUE: 设备名称
\fieldvalue{接管机} &
% #VALUE_ID: LEX-P0055-V0010
% #FIELD_VALUE: 品牌/型号
\fieldvalue{上海英莳} &
 &
% #VALUE_ID: LEX-P0055-V0013
% #FIELD_VALUE: 检查项中的手写记录
% #HANDWRITTEN: 无菌？2025.12
\fieldvalue{\handwritten{无菌？2025.12}} &
$\square$是\quad $\square$否\\
\hline
4 &
% #VALUE_ID: LEX-P0055-V0011
% #FIELD_VALUE: 设备名称
\fieldvalue{封管机} &
% #VALUE_ID: LEX-P0055-V0012
% #FIELD_VALUE: 品牌/型号
\fieldvalue{上海英莳} &
 & & $\square$是\quad $\square$否\\
\hline
5 & & & & & \\
\hline
\multicolumn{4}{|l|}{备注：} &
\multicolumn{2}{l|}{}\\
\hline
\end{tabularx}

\vspace{0.4em}

确认人/日期：
% #VALUE_ID: LEX-P0055-V0014
% #FIELD_VALUE: 确认人/日期
% #HANDWRITTEN: /
\fieldvalue{\handwritten{/}}
\hfill
复核人/日期：
% #VALUE_ID: LEX-P0055-V0015
% #FIELD_VALUE: 复核人/日期
% #HANDWRITTEN: /
\fieldvalue{\handwritten{/}}

\medskip
\textbf{5.1.2 物料确认表}

\begin{tabularx}{\textwidth}{|c|X|X|X|X|X|X|}
\hline
序号 & 名称 & 品牌 & 规格 & 批号 & 数量 & 是否符合要求\\
\hline
1 & & & & & & $\square$是\quad $\square$否\\
\hline
2 & & & & 
% #VALUE_ID: LEX-P0055-V0016
% #FIELD_VALUE: 批号
% #HANDWRITTEN: 2025.12
\fieldvalue{\handwritten{2025.12}} &
 & $\square$是\quad $\square$否\\
\hline
\multicolumn{6}{|l|}{备注：} &
\\
\hline
\end{tabularx}

\vspace{0.4em}

确认人/日期：
% #VALUE_ID: LEX-P0055-V0017
% #FIELD_VALUE: 确认人/日期
% #HANDWRITTEN: /
\fieldvalue{\handwritten{/}}
\hfill
复核人/日期：
% #VALUE_ID: LEX-P0055-V0018
% #FIELD_VALUE: 复核人/日期
% #HANDWRITTEN: /
\fieldvalue{\handwritten{/}}

\medskip
\textbf{5.2 操作内容}

\begin{tabularx}{\textwidth}{|X|X|}
\hline
项目 & 操作记录\\
\hline
 & \\
\hline
\end{tabularx}

% TODO：页面底部可见页码“第37页 共222页”。
% LEXOID_PAGE_COMPLETED: 55/59

\newpage

\begin{center}
\small
铂生物科技（北京）有限公司\\
{\tiny PLATINUM LIFE EXCELLENCE BIOTECH}\\[2pt]
\textbf{艾米近托赛注射液（Z2）批生产记录}
\end{center}

\noindent
\begin{tabularx}{\textwidth}{|>{\centering\arraybackslash}p{0.17\textwidth}|>{\centering\arraybackslash}p{0.25\textwidth}|>{\centering\arraybackslash}p{0.28\textwidth}|>{\centering\arraybackslash}X|}
\hline
品名 & 规格 & 生产批号 & 批量\\
\hline
艾米近托赛注射液 & 12ml：$6.0\times10^{7}$个细胞 &
% #VALUE_ID: LEX-P0056-V0002
% #FIELD_VALUE: 生产批号
% #TODO #HANDWRITTEN: 难辨；蓝色手写批号字符不清晰
\fieldvalue{\handwritten{难辨}} &
200--480袋\\
\hline
\end{tabularx}

\vspace{2pt}

\noindent
\begin{tabularx}{\textwidth}{|p{0.51\textwidth}|X|}
\hline
\begin{minipage}[t]{\linewidth}
\vspace{2pt}
1）参数检查：

检查培养参数控制值是否在如下范围：表面通气
20--50ml/min，底部通气 $10\pm3$ml/min，转速
35--50rpm，温度 $37.0\pm2.0^\circ\mathrm{C}$，溶氧
$\geq30\%$，pH 6.70--7.60。
\vspace{4pt}
\end{minipage}
&
\begin{minipage}[t]{\linewidth}
\vspace{2pt}
检查时间：\underline{\hspace{0.45cm}}年\underline{\hspace{0.45cm}}月\underline{\hspace{0.45cm}}日

搅拌转速：\underline{\hspace{0.7cm}}rpm

温度：\underline{\hspace{0.8cm}}$^\circ$C

DO：\underline{\hspace{0.8cm}}\%

pH：\underline{\hspace{1.2cm}}

表面流量：\underline{\hspace{0.8cm}}ml/min

底部流量：\underline{\hspace{0.8cm}}ml/min

培养体积：\underline{\hspace{0.8cm}}L

是否符合要求：
% #VALUE_ID: LEX-P0056-V0003
% #FIELD_VALUE: 是否符合要求
% #HANDWRITTEN: 是
\fieldvalue{\handwritten{是}} $\Box$ \quad 否 $\Box$
\vspace{2pt}
\end{minipage}
\\
\hline
\begin{minipage}[t]{\linewidth}
\vspace{2pt}
2）观察罐体内微载体分布情况，微载体上下分布不均
匀或一开始出现较多结团时，提高转速至40--50rpm，
每次2--5rpm。
\vspace{4pt}
\end{minipage}
&
\begin{minipage}[t]{\linewidth}
\vspace{2pt}
是否调整转速：是 $\Box$ \quad 否 $\Box$

调整时间：\underline{\hspace{0.45cm}}年\underline{\hspace{0.45cm}}月\underline{\hspace{0.45cm}}日

调整后转速：\underline{\hspace{0.9cm}}rpm
\vspace{2pt}
\end{minipage}
\\
\hline
\begin{minipage}[t]{\linewidth}
\vspace{2pt}
3）取样检测：

取样两袋，每袋体积$\geq20$ml，分别检测细胞密度、
活率、pH、葡萄糖浓度和细胞生长状态（镜检）。

注：pH检项与细胞密度/活率检验分袋检测，且pH为
取样袋开口后首次检测项目。

当离线pH与在线pH差值$>0.2$时，应对在线电极做
单点校准。
\vspace{4pt}
\end{minipage}
&
\begin{minipage}[t]{\linewidth}
\vspace{2pt}
取样时间：\underline{\hspace{0.45cm}}年\underline{\hspace{0.45cm}}月\underline{\hspace{0.45cm}}日

培养时长：\underline{\hspace{0.8cm}}h

细胞密度：\underline{\hspace{1.0cm}}个细胞/ml

细胞活率：\underline{\hspace{1.0cm}}\%

葡萄糖浓度：\underline{\hspace{1.0cm}}

离线pH：\underline{\hspace{1.0cm}}

电极是否校准：

是 $\Box$ \quad 否 $\Box$

镜检正常：

是 $\Box$ \quad 否 $\Box$
\vspace{2pt}
\end{minipage}
\\
\hline
\end{tabularx}

\vspace{3pt}

备注：\hrulefill

\vspace{8pt}

操作人/日期：
% #VALUE_ID: LEX-P0056-V0004
% #FIELD_VALUE: 操作人/日期
% #HANDWRITTEN: /
\fieldvalue{\handwritten{/}}
\hfill
复核人/日期：
% #VALUE_ID: LEX-P0056-V0005
% #FIELD_VALUE: 复核人/日期
% #HANDWRITTEN: /
\fieldvalue{\handwritten{/}}

\vfill

\begin{center}
{\tiny 第 38 页 共 222 页}
\end{center}
% LEXOID_PAGE_COMPLETED: 56/59

\newpage

\begin{center}
\textbf{\Large 艾米迈托美注射液（Z2）批生产记录}
\end{center}

\noindent
\begin{tabularx}{\textwidth}{|c|X|c|c|}
\hline
\multicolumn{2}{|c|}{原件} & 记录编号：YZ-BPR-MF-PD-01-020-a & 版本号：1.3 \\
\hline
品名 & 规格 & 生产批号 & 批量 \\
\hline
艾米迈托美注射液 & 12ml：6.0$\times 10^{7}$个/瓶 &
% #VALUE_ID: LEX-P0057-V0001
% #FIELD_VALUE: 生产批号
% #HANDWRITTEN: A352-?2025-06-?
\fieldvalue{\handwritten{A352-?2025-06-?}} &
200--480支 \\
\hline
\end{tabularx}

\noindent\textbf{5.3 清场}

\medskip
\noindent 操作依据：

\begin{enumerate}
  \item 《生产区清场管理规程》\hfill YZ-SMP-PD-03-001
  \item 《生产、检验状态标识管理规程》\hfill YZ-SMP-QA-04-014
\end{enumerate}

\noindent
生产工序：一级生物反应器细胞培养
\hfill
清场开始时间：
% #VALUE_ID: LEX-P0057-V0002
% #FIELD_VALUE: 清场开始时间
% #HANDWRITTEN: /
\fieldvalue{\handwritten{/}} 时 \rule{1.2em}{0.4pt} 分

\medskip
\noindent
清场房间编码：
% #VALUE_ID: LEX-P0057-V0003
% #FIELD_VALUE: 清场房间编码
% #HANDWRITTEN: /
\fieldvalue{\handwritten{/}}
\hfill
75\%酒精批号：
% #VALUE_ID: LEX-P0057-V0004
% #FIELD_VALUE: 75\%酒精批号
% #HANDWRITTEN: /
\fieldvalue{\handwritten{/}}

\medskip
\noindent
\begin{tabularx}{\textwidth}{|X|c|}
\hline
\multicolumn{1}{|c|}{清场项目及标准} & \multicolumn{1}{c|}{执行情况} \\
\hline
将本班次使用的各类工具配件恢复原位放置。 &
$\square$已执行\quad $\square$否 \\
\hline
将本班次生产有关的其它物品或物料清出生产现场。 &
$\square$已执行\quad $\square$否 \\
\hline
使用75\%酒精清洁设备、操作台面、不锈钢架子、转移车、传递窗等；
% #VALUE_ID: LEX-P0057-V0005
% #HANDWRITTEN: 签名及日期（字迹不清）
&
\multicolumn{1}{|l|}{\hspace{1em}\handwritten{签名及日期（字迹不清）}\hfill $\square$已执行\quad $\square$否} \\
\hline
使用75\%酒精喷洒地面，然后使用拖把擦拭地面； &
$\square$已执行\quad $\square$否 \\
\hline
按照《生产、检验状态标识管理规程》更换房间、设备状态标识； &
$\square$已执行\quad $\square$否 \\
\hline
将本班次操作产生的废弃物移出洁净生产区； &
$\square$已执行\quad $\square$否 \\
\hline
\end{tabularx}

\medskip
\noindent
清场结束时间：
% #VALUE_ID: LEX-P0057-V0006
% #FIELD_VALUE: 清场结束时间
% #HANDWRITTEN: /
\fieldvalue{\handwritten{/}} 时 \rule{1.2em}{0.4pt} 分

\medskip
\noindent
清场人员/日期：
% #VALUE_ID: LEX-P0057-V0007
% #FIELD_VALUE: 清场人员/日期
% #HANDWRITTEN: /
\fieldvalue{\handwritten{/}}

\medskip
\noindent
\begin{tabularx}{\textwidth}{|X|}
\hline
清场检查 \\
\hline
$\square$合格 \\
\hline
$\square$不合格，具体内容为： \\
\hline
整改情况及整改人： \\
\hline
\end{tabularx}

\medskip
\noindent
确认人/日期：
% #VALUE_ID: LEX-P0057-V0008
% #FIELD_VALUE: 确认人/日期
% #HANDWRITTEN: /
\fieldvalue{\handwritten{/}}

\begin{center}
第 39 页 共 222 页
\end{center}
% LEXOID_PAGE_COMPLETED: 57/59

\newpage

\begin{center}
\textbf{艾米诺托赛赛注射液（Z2）批生产记录}
\end{center}

记录编号：YZ-BPR-MF-PD-01-020a

\begin{tabularx}{\textwidth}{|X|X|X|}
\hline
品名 & 规格 & 生产批号 \\
\hline
艾米诺托赛注射液 & 12ml：$6.0\times 10^{7}$个细胞 &
% #VALUE_ID: LEX-P0058-V0001
% #FIELD_VALUE: 生产批号
% #TODO #HANDWRITTEN: 432210250?01; 字迹不清
\fieldvalue{\handwritten{432210250?01}} \\
\hline
\end{tabularx}

\section{一级生物反应器细胞收获}

\subsection{生产前确认：}

\begin{tabularx}{\textwidth}{|p{0.55\textwidth}|X|}
\hline
\centering 项目 & \centering 操作记录 \arraybackslash \\
\hline
操作人员按《人员进出洁净生产区管理规程》要求进入洁净区生产车间； &
操作间名称及编号：
% #VALUE_ID: LEX-P0058-V0002
% #FIELD_VALUE: 操作间名称及编号
% #TODO #HANDWRITTEN: 31#净间 2606; 字迹不清
\fieldvalue{\handwritten{31\#净间 2606}} \\
\hline
确认生产区操作间在清场有效期内（C级日清洁周期3天）； &
是否在清洁有效期内：

% #VALUE_ID: LEX-P0058-V0003
% #FIELD_VALUE: 是否在清洁有效期内：是
% #HANDWRITTEN: ✓
\fieldvalue{\handwritten{✓}} 是 \quad $\square$ 否

清洁有效期：
% #VALUE_ID: LEX-P0058-V0004
% #FIELD_VALUE: 清洁有效期
% #HANDWRITTEN: 2025.06.14
\fieldvalue{\handwritten{2025.06.14}} \\
\hline
确认操作间温度18--26℃、湿度30\%--75\%、相对压差$\geq 0$Pa（操作间对C级走廊）； &
温度：
% #VALUE_ID: LEX-P0058-V0005
% #FIELD_VALUE: 温度
% #HANDWRITTEN: 21.1
\fieldvalue{\handwritten{21.1}}\,℃ \quad
湿度：
% #VALUE_ID: LEX-P0058-V0006
% #FIELD_VALUE: 湿度
% #HANDWRITTEN: 55.6
\fieldvalue{\handwritten{55.6}}\,\%

压力：
% #VALUE_ID: LEX-P0058-V0007
% #FIELD_VALUE: 压力
% #HANDWRITTEN: 16
\fieldvalue{\handwritten{16}}\,Pa

是否符合要求：
% #VALUE_ID: LEX-P0058-V0008
% #FIELD_VALUE: 是否符合要求：是
% #HANDWRITTEN: ✓
\fieldvalue{\handwritten{✓}} 是 \quad $\square$ 否 \\
\hline
确认操作间无上批产品遗留物、文件及与本产品生产无关物料； &
% #VALUE_ID: LEX-P0058-V0009
% #FIELD_VALUE: 确认操作间无遗留物：是
% #HANDWRITTEN: ✓
\fieldvalue{\handwritten{✓}} 是 \quad $\square$ 否 \\
\hline
检查发放的批记录及其它相关配套文件均与生产指令相对应，且均接受； &
% #VALUE_ID: LEX-P0058-V0010
% #FIELD_VALUE: 检查批记录及配套文件：是
% #HANDWRITTEN: ✓
\fieldvalue{\handwritten{✓}} 是 \quad $\square$ 否 \\
\hline
确认设备完好，且在验证有效期内，计量器具有“合格证”且在有效期内，详见《主要设备/仪表确认表》； &
% #VALUE_ID: LEX-P0058-V0011
% #FIELD_VALUE: 设备及计量器具确认：是
% #HANDWRITTEN: ✓
\fieldvalue{\handwritten{✓}} 是 \quad $\square$ 否 \\
\hline
反应器氮源压力应控制在如下范围：压缩空气进气压力0.05--0.15MPa、二氧化碳进气压力0.05--0.15MPa、氧气进气压力0.05--0.15MPa； &
压缩空气：
% #VALUE_ID: LEX-P0058-V0012
% #FIELD_VALUE: 压缩空气压力
% #HANDWRITTEN: 0.02
\fieldvalue{\handwritten{0.02}}\,MPa

二氧化碳：
% #VALUE_ID: LEX-P0058-V0013
% #FIELD_VALUE: 二氧化碳压力
% #HANDWRITTEN: 0.12
\fieldvalue{\handwritten{0.12}}\,MPa

氧气：
% #VALUE_ID: LEX-P0058-V0014
% #FIELD_VALUE: 氧气压力
% #HANDWRITTEN: 0.1
\fieldvalue{\handwritten{0.1}}\,MPa

是否符合要求：
% #VALUE_ID: LEX-P0058-V0015
% #FIELD_VALUE: 是否符合要求：是
% #HANDWRITTEN: ✓
\fieldvalue{\handwritten{✓}} 是 \quad $\square$ 否 \\
\hline
将操作间状态标识更换为“生产加工状态卡”，写明品名、规格、批号等，并确认与本批生产指令一致； &
% #VALUE_ID: LEX-P0058-V0016
% #FIELD_VALUE: 操作间状态标识确认：是
% #HANDWRITTEN: ✓
\fieldvalue{\handwritten{✓}} 是 \quad $\square$ 否 \\
\hline
核对本次生产用物料耗材、溶液和试剂的品名、规格、批号、数量，检查物料外包装的完整性以及是否在有效期/复验期内，并记录相关信息，详见《物料确认表》和《溶液确认表》； &
% #VALUE_ID: LEX-P0058-V0017
% #FIELD_VALUE: 物料耗材、溶液和试剂确认：是
% #HANDWRITTEN: ✓
\fieldvalue{\handwritten{✓}} 是 \quad $\square$ 否 \\
\hline
备注： &
% #VALUE_ID: LEX-P0058-V0018
% #FIELD_VALUE: 备注
% #HANDWRITTEN: /
\fieldvalue{\handwritten{/}} \\
\hline
\end{tabularx}

\noindent
操作人/日期：
% #VALUE_ID: LEX-P0058-V0019
% #FIELD_VALUE: 操作人
% #HANDWRITTEN: 彭博
\fieldvalue{\handwritten{彭博}}
\quad
% #VALUE_ID: LEX-P0058-V0020
% #FIELD_VALUE: 操作日期
% #HANDWRITTEN: 2025.06.12
\fieldvalue{\handwritten{2025.06.12}}

\hfill
复核人/日期：
% #VALUE_ID: LEX-P0058-V0021
% #FIELD_VALUE: 复核人
% #TODO #HANDWRITTEN: 赵宇军; 字迹不清
\fieldvalue{\handwritten{赵宇军}}
\quad
% #VALUE_ID: LEX-P0058-V0022
% #FIELD_VALUE: 复核日期
% #HANDWRITTEN: 2025.06.12
\fieldvalue{\handwritten{2025.06.12}}

\subsubsection{6.1.1 主要设备/仪表确认表}

\begin{center}
第40页 共222页
\end{center}
% LEXOID_PAGE_COMPLETED: 58/59

\newpage

\begin{center}
\textbf{艾米迈托赛注射液（Z2）批生产记录}
\end{center}

记录编号：
% #VALUE_ID: LEX-P0059-V0001
% #FIELD_VALUE: 记录编号
\fieldvalue{YZ-BPR-MF-PD-01-020a-...}
\hfill
版本号：
% #VALUE_ID: LEX-P0059-V0002
% #FIELD_VALUE: 版本号
\fieldvalue{1.3}

\begin{tabularx}{\textwidth}{|c|X|X|X|}
\hline
品名 & 规格 & 生产批号 & 批量 \\
\hline
% #VALUE_ID: LEX-P0059-V0003
% #FIELD_VALUE: 品名
\fieldvalue{艾米迈托赛注射液} &
% #VALUE_ID: LEX-P0059-V0004
% #FIELD_VALUE: 规格
\fieldvalue{12ml：6.0$\times$10\textsuperscript{?}个细胞} &
% #VALUE_ID: LEX-P0059-V0005
% #FIELD_VALUE: 生产批号
% #TODO #HANDWRITTEN: 邯32220/060?0; 字迹不清
\fieldvalue{\handwritten{邯32220/060?0}} &
% #VALUE_ID: LEX-P0059-V0006
% #FIELD_VALUE: 批量
\fieldvalue{200--480袋} \\
\hline
\end{tabularx}

\begin{tabularx}{\textwidth}{|c|X|X|X|X|X|}
\hline
序号 & 设备名称 & 品牌/型号 & 设备编号 & 检查项 & 确认结果 \\
\hline
1 &
% #VALUE_ID: LEX-P0059-V0007
% #FIELD_VALUE: 设备名称
\fieldvalue{生物反应器控制台（主机）} &
% #VALUE_ID: LEX-P0059-V0008
% #FIELD_VALUE: 品牌/型号
\fieldvalue{华龛生物} &
% #VALUE_ID: LEX-P0059-V0009
% #FIELD_VALUE: 设备编号
% #HANDWRITTEN: 1123610
\fieldvalue{\handwritten{1123610}} &
设备应有“完好”标识；设备经验证有效期内；计量器具有“合格”证且在有效期内。 &
% #VALUE_ID: LEX-P0059-V0010
% #FIELD_VALUE: 确认结果
% #HANDWRITTEN: 是
\fieldvalue{\handwritten{是}} \\
\hline
2 &
% #VALUE_ID: LEX-P0059-V0011
% #FIELD_VALUE: 设备名称
\fieldvalue{生物反应器罐体} &
% #VALUE_ID: LEX-P0059-V0012
% #FIELD_VALUE: 品牌/型号
\fieldvalue{华龛生物} &
% #VALUE_ID: LEX-P0059-V0013
% #FIELD_VALUE: 设备编号
% #HANDWRITTEN: V5610502-220038
\fieldvalue{\handwritten{V5610502-220038}} &
\textit{同上} &
% #VALUE_ID: LEX-P0059-V0014
% #FIELD_VALUE: 确认结果
% #HANDWRITTEN: 是
\fieldvalue{\handwritten{是}} \\
\hline
3 &
% #VALUE_ID: LEX-P0059-V0015
% #FIELD_VALUE: 设备名称
\fieldvalue{接管机} &
% #VALUE_ID: LEX-P0059-V0016
% #FIELD_VALUE: 品牌/型号
\fieldvalue{上海英启} &
% #VALUE_ID: LEX-P0059-V0017
% #FIELD_VALUE: 设备编号
% #HANDWRITTEN: 1123856
\fieldvalue{\handwritten{1123856}} &
\textit{同上} &
% #VALUE_ID: LEX-P0059-V0018
% #FIELD_VALUE: 确认结果
% #HANDWRITTEN: 是
\fieldvalue{\handwritten{是}} \\
\hline
4 &
% #VALUE_ID: LEX-P0059-V0019
% #FIELD_VALUE: 设备名称
\fieldvalue{封管机} &
% #VALUE_ID: LEX-P0059-V0020
% #FIELD_VALUE: 品牌/型号
\fieldvalue{上海英启} &
% #VALUE_ID: LEX-P0059-V0021
% #FIELD_VALUE: 设备编号
% #HANDWRITTEN: 1123612
\fieldvalue{\handwritten{1123612}} &
\textit{同上} &
% #VALUE_ID: LEX-P0059-V0022
% #FIELD_VALUE: 确认结果
% #HANDWRITTEN: 是
\fieldvalue{\handwritten{是}} \\
\hline
5 & & &
% #VALUE_ID: LEX-P0059-V0023
% #FIELD_VALUE: 设备编号
% #TODO #HANDWRITTEN: 112?612; 字迹不清
\fieldvalue{\handwritten{112?612}} &
& \\
\hline
6 & & & & & \\
\hline
\end{tabularx}

备注：
% #VALUE_ID: LEX-P0059-V0024
% #HANDWRITTEN: /
\handwritten{/}

\noindent
确认人/日期：
% #VALUE_ID: LEX-P0059-V0025
% #FIELD_VALUE: 确认人
% #HANDWRITTEN: 陈梅
\fieldvalue{\handwritten{陈梅}}
\quad
% #VALUE_ID: LEX-P0059-V0026
% #FIELD_VALUE: 确认日期
% #HANDWRITTEN: 2025.06.12
\fieldvalue{\handwritten{2025.06.12}}
\hfill
复核人/日期：
% #VALUE_ID: LEX-P0059-V0027
% #FIELD_VALUE: 复核人
% #HANDWRITTEN: 李宇
\fieldvalue{\handwritten{李宇}}
\quad
% #VALUE_ID: LEX-P0059-V0028
% #FIELD_VALUE: 复核日期
% #HANDWRITTEN: 2025.06.12
\fieldvalue{\handwritten{2025.06.12}}

\subsection*{6.1.2 物料确认表}

\begin{tabularx}{\textwidth}{|c|X|X|X|X|c|X|}
\hline
序号 & 名称 & 品牌 & 规格 & 批号 & 数量 & 是否符合要求 \\
\hline
1 &
% #VALUE_ID: LEX-P0059-V0029
% #FIELD_VALUE: 名称
\fieldvalue{无菌葡萄糖（3L）} &
% #VALUE_ID: LEX-P0059-V0030
% #FIELD_VALUE: 品牌
% #HANDWRITTEN: 金诚世
\fieldvalue{\handwritten{金诚世}} &
% #VALUE_ID: LEX-P0059-V0031
% #FIELD_VALUE: 规格
\fieldvalue{1个/包} &
% #VALUE_ID: LEX-P0059-V0032
% #FIELD_VALUE: 批号
% #HANDWRITTEN: H138?0 20250601
\fieldvalue{\handwritten{H138?0 20250601}} &
% #VALUE_ID: LEX-P0059-V0033
% #FIELD_VALUE: 数量
% #HANDWRITTEN: 1
\fieldvalue{\handwritten{1}} &
% #VALUE_ID: LEX-P0059-V0034
% #FIELD_VALUE: 是否符合要求
% #HANDWRITTEN: 是
\fieldvalue{\handwritten{是}} \\
\hline
2 &
% #VALUE_ID: LEX-P0059-V0035
% #FIELD_VALUE: 名称
% #HANDWRITTEN: 250ml一次性储液袋
\fieldvalue{\handwritten{250ml一次性储液袋}} &
% #VALUE_ID: LEX-P0059-V0036
% #FIELD_VALUE: 品牌
% #HANDWRITTEN: 金诚世
\fieldvalue{\handwritten{金诚世}} &
% #VALUE_ID: LEX-P0059-V0037
% #FIELD_VALUE: 规格
% #HANDWRITTEN: 1/包
\fieldvalue{\handwritten{1/包}} &
% #VALUE_ID: LEX-P0059-V0038
% #FIELD_VALUE: 批号
% #HANDWRITTEN: R131?2025-3301
\fieldvalue{\handwritten{R131?2025-3301}} &
% #VALUE_ID: LEX-P0059-V0039
% #FIELD_VALUE: 数量
% #HANDWRITTEN: 1
\fieldvalue{\handwritten{1}} &
% #VALUE_ID: LEX-P0059-V0040
% #FIELD_VALUE: 是否符合要求
% #HANDWRITTEN: 是
\fieldvalue{\handwritten{是}} \\
\hline
3 & & & &
% #VALUE_ID: LEX-P0059-V0041
% #FIELD_VALUE: 批号
% #HANDWRITTEN: 李宇 2025.06.12
\fieldvalue{\handwritten{李宇 2025.06.12}} &
& \\
\hline
\end{tabularx}

备注：
% #VALUE_ID: LEX-P0059-V0042
% #HANDWRITTEN: /
\handwritten{/}

\noindent
确认人/日期：
% #VALUE_ID: LEX-P0059-V0043
% #FIELD_VALUE: 确认人
% #HANDWRITTEN: 陈梅
\fieldvalue{\handwritten{陈梅}}
\quad
% #VALUE_ID: LEX-P0059-V0044
% #FIELD_VALUE: 确认日期
% #HANDWRITTEN: 2025.06.12
\fieldvalue{\handwritten{2025.06.12}}
\hfill
复核人/日期：
% #VALUE_ID: LEX-P0059-V0045
% #FIELD_VALUE: 复核人
% #HANDWRITTEN: 李宇
\fieldvalue{\handwritten{李宇}}
\quad
% #VALUE_ID: LEX-P0059-V0046
% #FIELD_VALUE: 复核日期
% #HANDWRITTEN: 2025.06.12
\fieldvalue{\handwritten{2025.06.12}}

\subsection*{6.1.3 溶液确认表}

\begin{tabularx}{\textwidth}{|c|X|X|X|c|X|X|}
\hline
序号 & 名称 & 规格 & 批号 & 数量 & 有效期 & 是否符合要求 \\
\hline
1 &
% #VALUE_ID: LEX-P0059-V0047
% #FIELD_VALUE: 名称
\fieldvalue{裂解液溶液} &
% #VALUE_ID: LEX-P0059-V0048
% #FIELD_VALUE: 规格
\fieldvalue{150ml/袋} &
% #VALUE_ID: LEX-P0059-V0049
% #FIELD_VALUE: 批号
% #HANDWRITTEN: 2025061001
\fieldvalue{\handwritten{2025061001}} &
% #VALUE_ID: LEX-P0059-V0050
% #FIELD_VALUE: 数量
% #HANDWRITTEN: 1
\fieldvalue{\handwritten{1}} &
% #VALUE_ID: LEX-P0059-V0051
% #FIELD_VALUE: 有效期
% #HANDWRITTEN: 2025.07.07
\fieldvalue{\handwritten{2025.07.07}} &
% #VALUE_ID: LEX-P0059-V0052
% #FIELD_VALUE: 是否符合要求
% #HANDWRITTEN: 是
\fieldvalue{\handwritten{是}} \\
\hline
\end{tabularx}

\begin{center}
第 41 页 共 222 页
\end{center}

\end{document}
% LEXOID_PAGE_COMPLETED: 59/59
